% Options for packages loaded elsewhere
\PassOptionsToPackage{unicode}{hyperref}
\PassOptionsToPackage{hyphens}{url}
\documentclass[
  ukrainian,
]{article}
\usepackage{xcolor}
\usepackage[margin=2.5cm]{geometry}
\usepackage{amsmath,amssymb}
\setcounter{secnumdepth}{-\maxdimen} % remove section numbering
\usepackage{iftex}
\ifPDFTeX
  \usepackage[T1]{fontenc}
  \usepackage[utf8]{inputenc}
  \usepackage{textcomp} % provide euro and other symbols
\else % if luatex or xetex
  \usepackage{unicode-math} % this also loads fontspec
  \defaultfontfeatures{Scale=MatchLowercase}
  \defaultfontfeatures[\rmfamily]{Ligatures=TeX,Scale=1}
\fi
\usepackage{lmodern}
\ifPDFTeX\else
  % xetex/luatex font selection
\fi
% Use upquote if available, for straight quotes in verbatim environments
\IfFileExists{upquote.sty}{\usepackage{upquote}}{}
\IfFileExists{microtype.sty}{% use microtype if available
  \usepackage[]{microtype}
  \UseMicrotypeSet[protrusion]{basicmath} % disable protrusion for tt fonts
}{}
\makeatletter
\@ifundefined{KOMAClassName}{% if non-KOMA class
  \IfFileExists{parskip.sty}{%
    \usepackage{parskip}
  }{% else
    \setlength{\parindent}{0pt}
    \setlength{\parskip}{6pt plus 2pt minus 1pt}}
}{% if KOMA class
  \KOMAoptions{parskip=half}}
\makeatother
\usepackage{longtable,booktabs,array}
\newcounter{none} % for unnumbered tables
\usepackage{calc} % for calculating minipage widths
% Correct order of tables after \paragraph or \subparagraph
\usepackage{etoolbox}
\makeatletter
\patchcmd\longtable{\par}{\if@noskipsec\mbox{}\fi\par}{}{}
\makeatother
% Allow footnotes in longtable head/foot
\IfFileExists{footnotehyper.sty}{\usepackage{footnotehyper}}{\usepackage{footnote}}
\makesavenoteenv{longtable}
\usepackage{graphicx}
\graphicspath{{/Users/kalinicenkovitalijmikolajovic/.codex/worktrees/5c15/Script_for_translate_book/Output_Final/}}
\makeatletter
\newsavebox\pandoc@box
\newcommand*\pandocbounded[1]{% scales image to fit in text height/width
  \sbox\pandoc@box{#1}%
  \Gscale@div\@tempa{\textheight}{\dimexpr\ht\pandoc@box+\dp\pandoc@box\relax}%
  \Gscale@div\@tempb{\linewidth}{\wd\pandoc@box}%
  \ifdim\@tempb\p@<\@tempa\p@\let\@tempa\@tempb\fi% select the smaller of both
  \ifdim\@tempa\p@<\p@\scalebox{\@tempa}{\usebox\pandoc@box}%
  \else\usebox{\pandoc@box}%
  \fi%
}
% Set default figure placement to htbp
\def\fps@figure{htbp}
\makeatother
\ifLuaTeX
\usepackage[bidi=basic,shorthands=off,provide=*]{babel}
\else
\usepackage[bidi=default,shorthands=off,provide=*]{babel}
\fi
\ifLuaTeX
  \usepackage{selnolig} % disable illegal ligatures
\fi
\setlength{\emergencystretch}{3em} % prevent overfull lines
\providecommand{\tightlist}{%
  \setlength{\itemsep}{0pt}\setlength{\parskip}{0pt}}
\usepackage{bookmark}
\IfFileExists{xurl.sty}{\usepackage{xurl}}{} % add URL line breaks if available
\urlstyle{same}
\hypersetup{
  pdflang={uk},
  hidelinks,
  pdfcreator={LaTeX via pandoc}}

\author{}
\date{}

\begin{document}

ЧАСТКОВО ДИФЕРЕНЦІАЛ РІВНЯННЯ для Вчені та Інженери

Стенлі Дж. Фарлоу, професор математики, Університет Мену

\section{С. Фарлоу}\label{ux441.-ux444ux430ux440ux43bux43eux443}

\section{РІВНЯННЯ З ПОХІДНИМИ ПОХІДНИМИ ДЛЯ НАУКОВЦІВ І
ІНЖЕНЕРІВ}\label{ux440ux456ux432ux43dux44fux43dux43dux44f-ux437-ux43fux43eux445ux456ux434ux43dux438ux43cux438-ux43fux43eux445ux456ux434ux43dux438ux43cux438-ux434ux43bux44f-ux43dux430ux443ux43aux43eux432ux446ux456ux432-ux456-ux456ux43dux436ux435ux43dux435ux440ux456ux432}

Перекладено з англійської \textbf{A.I.} Плісом

Редактор

С. І. Похожаєва

5\textbar{} 66K 22.161.6 φ24 υ ΥΠΚ 517.2

\section{Фарлоу С.}\label{ux444ux430ux440ux43bux43eux443-ux441.}

F 24 Рівняння з частими похідними для науковців та інженерів. Москва,
видавництво Мир, 1985, 384 с.

King's American Mathematician, який є підручником з теорії
диференціальних рівнянь у похідних похідних. Він характеризується
компактністю, ясністю та ясністю подачі, а також неформальним підходом
до подачі матеріалу. Вона містить багато ілюстрацій, графіків і діаграм,
але замість суворих доказів часто даються роздуми, засновані на інтуїції
або аналогіях.

Для інженерів і нематематиків --- біологів, диміків, а також студентів
університету.
\texttt{\textbackslash{}Phi\ =\ \textbackslash{}frac\{1702050000-380\}\{041(01)-85\}\ 26-85,\ \textbackslash{}text\{\ q.\ 1\}}
BBK 22.161.6

Редакційна колегія журналу Literature on Mathematical Sciences

© переклад російською, Мир, 1985

© John Wiley \& Sons, Inc.~1982. Всі права захищені. Авторизований
переклад з англомовного видання, опублікований John Wiley \& Sons, Inc.

\section{З РЕДАКТОРА
ПЕРЕКЛАДУ}\label{ux437-ux440ux435ux434ux430ux43aux442ux43eux440ux430-ux43fux435ux440ux435ux43aux43bux430ux434ux443}

Книга, запропонована для читачів, слід розглядати як підручник для
курсів: «Диференціальні рівняння в похідних похідних» та «Рівняння
математичної фізики» для студентів технічних університетів. Вона також
буде корисною для інженерів різних спеціальностей. У книзі представлено
традиційні методи розв'язання диференціальних рівнянь у похідних
похідних, включаючи деякі чисельні методи. Водночас особлива увага
приділяється методу розділення змінних і методу інтегральних
перетворень, і цю частину книги можна використовувати на практичних
заняттях в університетах

Орієнтація книги на широку «нематематичну» аудиторію визначила стиль
викладання розглянутих питань: інтуїтивний підхід і велика увага до
фізичного значення не лише самих рівнянь, а й меж і початкових умов для
різних задач. Вважається, що читачі знайомі лише з основами
диференціального та інтегрального числення та елементами теорії
звичайних диференціальних рівнянь.

Фахівці можуть не задовольнитися інтуїтивним описом певних понять і
методів автором, але більш суворий виклад неминуче призведе до
ускладнення книги і звузить коло її читачів, що, очевидно, не було
наміром автора.

Більшим запереченням може бути відсутність загальних результатів щодо
задач для диференціальних рівнянь у похідних з змінними коефіцієнтами,
включно з рівняннями вищого порядку. Однак читач має чітко розуміти, що
ця книга є вступним курсом з теорії диференціальних рівнянь у часткових
похідних. Питання лінійної теорії диференціальних рівнянь, які автор не
включив, можна знайти в літературі, наведеній наприкінці книги.

Матеріал книги подається автором з великою методологічною майстерністю.
Однак не все було однаково успішним, і винний читач знайде це,
наприклад, у вивченні принципу Дюамеля, сферичних заповітів та деяких
інших уривків у книзі.

Без сумніву, запропонована книга С. Фарлоу знайде читача не лише серед
студентів і інженерів, а й серед викладачів, які можуть багато чого
навчитися з конкретного способу презентації та розробки спеціальних
математичних курсів.

\section{ПЕРЕДМОВА}\label{ux43fux435ux440ux435ux434ux43cux43eux432ux430}

За останні кілька років спостерігається значне зростання кількості
учнів, які починають вивчати диференціальні рівняння з похідними у своїх
одинадцятих класах. Крім того, багато з них спеціалізуються на сферах,
де інтуїція важливіша за математичну строгість. Тому, пишучи цю книгу, я
прагнув створити текст, який стимулює інтуїтивне мислення читачів, не
втрачаючи при цьому значної математичної строгості. Можливо, на одному
полюсі, подати предмет на високому рівні «епсилон-дельта», але тоді, як
це зазвичай буває, студенти не знатимуть, що робити з отриманою
інформацією. Інша крайність --- повне ігнорування математичних тонкощів,
але тоді ні учні, ні вчителі не знатимуть, як діяти. Я намагався знайти
розумний компроміс між цими крайнощами.

Ця книга виникла з курсу лекцій, які я читав протягом останніх п'яти
років. Розташування матеріалу не є загальноприйнятим --- вони майже
незалежні один від одного

Друг лекцій.

Основна увага приділяється двом найважливішим аналітичним методам:
розділенню змінних та інтегральним перетворенням. Також представлені
деякі нестандартні розділи, такі як метод Монте-Карло, варіаційне
числення, теорія керування, теорія потенціалів, інтегральні рівняння. Це
пов'язано з тим, що більшості студентів доведеться стикатися з цими
питаннями під час навчання, і якщо вони не ознайомляться з ними зараз,
навряд чи коли-небудь систематично їх вивчатимуть.

Цю книгу можна використовувати як навчальний посібник для одно- або
двосеместрового курсу як для початківців, так і для старших студентів.
Читачі повинні знати лише диференціальне та інтегральне числення та
звичайні диференціальні рівняння. Матеріал більшості лекцій подається в
одному-двох уроках. Типовий односеместровий курс може включати лекції
1-13, 15-17, 19-20, 22-23, 25-27, 30-32, 37-39. Усі 47 лекцій легко
читати протягом двох семестрів. Водночас часу буде достатньо для
вирішення проблем.

Автор вдячний Вайлі за пропозицію написати книгу, а також рецензентам
--- професорам Крісу Рорресу та М. Куршиду Алі, які зробили низку
корисних коментарів до тексту книги. Будь-які пропозиції щодо подальшого
покращення цього підручника будуть вдячні як від учнів, так і від
викладачів. Я також вдячний Дороті, Сюзан, Александер і Дезі Фарлоу.

Стенлі Дж. Фарлоу

\section{ВСТУП ДО ТЕОРІЇ ДИФЕРЕНЦІАЛЬНИХ РІВНЯНЬ У ПОХІДНИХ
ПОХІДНИХ}\label{ux432ux441ux442ux443ux43f-ux434ux43e-ux442ux435ux43eux440ux456ux457-ux434ux438ux444ux435ux440ux435ux43dux446ux456ux430ux43bux44cux43dux438ux445-ux440ux456ux432ux43dux44fux43dux44c-ux443-ux43fux43eux445ux456ux434ux43dux438ux445-ux43fux43eux445ux456ux434ux43dux438ux445}

МЕТА ЛЕКЦІЇ: Показати, що таке диференціальні рівняння в похідних
похідних, де і як вони виникають і як їх розв'язують. Коротко розгляньте
методи класифікації рівнянь і надайте список базових понять, які будуть
детально вивчені на наступних лекціях.

Більшість фізичних явищ у таких сферах, як гідродинаміка, електрика та
магнетизм, механіка, оптика та теплопередача, загалом можна описати за
допомогою рівнянь у похідних похідних (PNP). Більшість рівнянь у
математичній фізиці є диференціальними рівняннями з похідними похідними.
Правда, з деякими спрощеними припущеннями ці рівняння зводяться до
звичайних диференціальних рівнянь, але повний опис таких систем неминуче
призводить до використання диференціальних рівнянь у похідних похідних.

\section{Що таке часткові диференціальні
рівняння?}\label{ux449ux43e-ux442ux430ux43aux435-ux447ux430ux441ux442ux43aux43eux432ux456-ux434ux438ux444ux435ux440ux435ux43dux446ux456ux430ux43bux44cux43dux456-ux440ux456ux432ux43dux44fux43dux43dux44f}

Рівняння з частковими похідними --- це рівняння, яке містить
диференціали у похідних похідних. На відміну від звичайних
диференціальних рівнянь (TAC), де невідома функція залежить лише від
однієї змінної, у рівняннях часткових похідних вона залежить від кількох
змінних (наприклад, температура \texttt{u(x,\ t)} залежить від
координати x і часу t).

Тепер наведемо деякі з найважливіших диференціальних рівнянь у похідних
похідних. Для спрощення запису ми використаємо наступні нотації

\begin{verbatim}
u_t \equiv \frac{\partial u}{\partial t}
\end{verbatim}

\texttt{,\ \textquotesingle{}u\_x\ \textbackslash{}equiv\ \textbackslash{}frac\{\textbackslash{}partial\ u\}\{\textbackslash{}partial\ x\}\textquotesingle{}\ ,\ \textquotesingle{}u\_\{xx\}\ \textbackslash{}equiv\ \textbackslash{}frac\{\textbackslash{}partial\^{}2\ u\}\{\textbackslash{}partial\ x\^{}2\}\textquotesingle{}\ ,\ ...}

\section{Деякі диференціальні рівняння в похідних
похідних}\label{ux434ux435ux44fux43aux456-ux434ux438ux444ux435ux440ux435ux43dux446ux456ux430ux43bux44cux43dux456-ux440ux456ux432ux43dux44fux43dux43dux44f-ux432-ux43fux43eux445ux456ux434ux43dux438ux445-ux43fux43eux445ux456ux434ux43dux438ux445}

\begin{verbatim}
u_t = u_{xx}
\end{verbatim}

(одновимірне рівняння теплопровідності),
\texttt{u\_t\ =\ u\_\{xx\}\ +\ u\_\{yy\}} (двовимірне рівняння
теплопровідності),
\texttt{u\_\{rr\}\ +\ \textbackslash{}frac\{1\}\{r\}u\_r\ +\ \textbackslash{}frac\{1\}\{r\^{}2\}u\_\{00\}\ =\ 0}
(рівняння Лапласа за полярними координатами),
\texttt{u\_\{tt\}\ =\ u\_\{xx\}\ +\ u\_\{yy\}\ +\ u\_\{zz\}} (тривимірне
хвильове рівняння),
\texttt{u\_\{tt\}\ =\ u\_\{xx\}\ +\ zu\_t\ +\ \textbackslash{}beta\ u}
(телеграфне рівняння).

\section{Примітки щодо наведених
прикладів}\label{ux43fux440ux438ux43cux456ux442ux43aux438-ux449ux43eux434ux43e-ux43dux430ux432ux435ux434ux435ux43dux438ux445-ux43fux440ux438ux43aux43bux430ux434ux456ux432}

У всіх наведених вище прикладах функція невідома і залежить від більше
ніж однієї змінної. Змінна і (яку ми диференціюємо) називається залежною
змінною. Змінні, за якими відбувається диференціювання, називаються
незалежними змінними. Наприклад, у рівнянні.

\begin{verbatim}
u_t = u_{xx}
\end{verbatim}

Залежна змінна \texttt{u(x,\ t)} є функцією двох незалежних змінних x і
t; та в рівнянні

\begin{verbatim}
u_t = u_{rr} + \frac{1}{r} u_r + \frac{1}{r^2} u_{00}
\end{verbatim}

змінна \texttt{u(r,\ \textbackslash{}theta,\ t)} залежить від
\texttt{r,\ \textbackslash{}theta,\ t} .

\section{Чому необхідно вивчати диференціальні рівняння у похідних
похідних?}\label{ux447ux43eux43cux443-ux43dux435ux43eux431ux445ux456ux434ux43dux43e-ux432ux438ux432ux447ux430ux442ux438-ux434ux438ux444ux435ux440ux435ux43dux446ux456ux430ux43bux44cux43dux456-ux440ux456ux432ux43dux44fux43dux43dux44f-ux443-ux43fux43eux445ux456ux434ux43dux438ux445-ux43fux43eux445ux456ux434ux43dux438ux445}

Більшість фізичних законів природи можна сформулювати через рівняння з
похідними похідними. Прикладами є рівняння Максвелла, закон
теплопередачі Ньютона, рівняння Нав'є-Стокса, рівняння руху Ньютона та
рівняння Шредінгера в кватичній механіці. У всіх цих рівняннях фізичні
явища описуються мовою просторових і часових похідних. Похідні
з'являються в рівняннях, оскільки вони описують найважливіші фізичні
величини (такі як швидкість, прискорення, сила, тертя, потік, струм
тощо). Таким чином, виникають диференціальні рівняння в частинних
похідних, які містять невідому функцію, яку потрібно визначити.

Мета цієї книги --- показати читачеві:

\begin{enumerate}
\def\labelenumi{\arabic{enumi}.}
\item
  Як сформулювати фізичну задачу у формі диференціального рівняння у
  похідних похідних (побудова математичної моделі).
\item
  Як розв'язати диференціальне рівняння в похідних (з урахуванням
  початкових і граничних умов).
\end{enumerate}

Перш ніж перейти до побудови математичних моделей, давайте коротко
розглянемо методи розв'язання рівнянь у похідних похідних.

\section{Як розв'язати диференціальні рівняння в похідних
похідних?}\label{ux44fux43a-ux440ux43eux437ux432ux44fux437ux430ux442ux438-ux434ux438ux444ux435ux440ux435ux43dux446ux456ux430ux43bux44cux43dux456-ux440ux456ux432ux43dux44fux43dux43dux44f-ux432-ux43fux43eux445ux456ux434ux43dux438ux445-ux43fux43eux445ux456ux434ux43dux438ux445}

Гарне питання. Виявляється, існує цілий арсенал методів, придатних для
практичного використання. Найважливішими є ті, у яких диференціальні
рівняння з похідними похідними зводяться до звичайних диференціальних
рівнянь. Перелічимо десять методів розв'язання диференціальних рівнянь у
похідних похідних:

\begin{enumerate}
\def\labelenumi{\arabic{enumi}.}
\tightlist
\item
  Метод розділення змінних. Рівняння з факторами, отриманими з n
  незалежними змінними, зводиться до n звичайних-
\end{enumerate}

диференціальнірівняння 1).

\begin{enumerate}
\def\labelenumi{\arabic{enumi}.}
\setcounter{enumi}{1}
\item
  Метод інтегральних перетворень. Рівняння в похідних з n незалежними
  змінними зводиться до рівняння з похідними у (n-1) незалежних змінних;
  Відповідно, рівняння з частковими похідними з двома незалежними
  змінними можна звести до звичайного диференціального рівняння
  \texttt{\^{}2} ).
\item
  Метод перетворення координат. Початкове диференціальне рівняння в
  похідних похідних зводиться до звичайного рівняння з частковими
  похідними або до іншого, простішого рівняння у похідних похідних
  шляхом відповідного перетворення координат (наприклад, обертання осей
  координат тощо).
\item
  Перетворення залежної змінної. Початкове диференціальне рівняння в
  похідних похідних перетворюється у таке рівняння з похідними похідними
  для іншої невідомої функції, яка
\end{enumerate}

легше розв'язати, ніж початкові 3).

\begin{enumerate}
\def\labelenumi{\arabic{enumi}.}
\setcounter{enumi}{4}
\tightlist
\item
  Чисельні методи. Початкове диференціальне рівняння в похідних похідних
  зводиться до системи різницей, яку розв'язують методом ітерацій на
  комп'ютері. У багатьох випадках це єдиний спосіб розв'язати
  диференціальне рівняння з похідними похідними. Окрім різницьких
  методів для розв'язання рівнянь з частковими про-
\end{enumerate}

\begin{enumerate}
\def\labelenumi{\alph{enumi})}
\setcounter{enumi}{4}
\tightlist
\item
  Ефективніше використовувати загальні перетворення незалежних і
  залежних змінних, включаючи похідні. Ед.
\end{enumerate}

4) Це стосується повного розділення змінних. При частковому розділенні
змінних одне диференціальне рівняння в похідних похідних зводиться до
кількох рівнянь у часткових похідних з меншою кількістю незалежних
змінних. --- Ед. Ед.

2) Це відбувається у випадку одномірного інтегрального перетворення. У
випадку k-вимірного інтегрального перетворення LCP з n незалежними
змінними зводиться до рівняння часткових з n---k незалежними змінними.
--- Ед. Ед.

Існують також інші чисельні методи, зокрема ті, що базуються на
наближенні розв'язків за поліноміальними поверхнями (наближення за
сплайнами).

\begin{enumerate}
\def\labelenumi{\arabic{enumi}.}
\setcounter{enumi}{5}
\tightlist
\item
  Метод теорії збурень. Початкова нелінійна задача зводиться до
  послідовності лінійних задач, апрокенн-
\end{enumerate}

нелінійна задача 1).

\begin{itemize}
\tightlist
\item
  \begin{enumerate}
  \def\labelenumi{\arabic{enumi}.}
  \setcounter{enumi}{6}
  \tightlist
  \item
    Метод функцій Гріна. Початкові та граничні умови замінюються
    системою найпростіших джерел, і задача розв'язується для кожного
    найпростішого джерела. Повне розв'язання початкової задачі
    отримується шляхом підсумовування розв'язків елементарних джерел.
  \end{enumerate}
\item
  \begin{enumerate}
  \def\labelenumi{\arabic{enumi}.}
  \setcounter{enumi}{7}
  \tightlist
  \item
    Метод інтегральних рівнянь. Рівняння з частковими похідними
    зводяться до інтегрального рівняння (рівняння, в якому невідома
    функція знаходиться під знаком інтеграла). Існує багато різних
    методів розв'язання інтегральних рівнянь.
  \end{enumerate}
\item
  \begin{enumerate}
  \def\labelenumi{\arabic{enumi}.}
  \setcounter{enumi}{8}
  \tightlist
  \item
    Варіаційні методи. Замість рівняння з частковими похідними
    диференціальними розв'язують задачу мінімізації. Виявляється, що
    функція, яка дає мінімум певному виразу (наприклад, загальній
    енергії системи), одночасно є розв'язком початкового рівняння
    похідниху похідних 2.
  \end{enumerate}
\item
  \begin{enumerate}
  \def\labelenumi{\arabic{enumi}.}
  \setcounter{enumi}{9}
  \tightlist
  \item
    Метод розкладу за власними функціями. Розв'язок рівняння часткових
    похідних шукається як послідовність власних функцій. Ці власні
    функції знаходяться як розв'язки так званої задачі власних значень,
    що відповідає початковій задачі для рівняння часткових похідних.
  \end{enumerate}
\end{itemize}

\section{Типи рівнянь у похідних
похідних}\label{ux442ux438ux43fux438-ux440ux456ux432ux43dux44fux43dux44c-ux443-ux43fux43eux445ux456ux434ux43dux438ux445-ux43fux43eux445ux456ux434ux43dux438ux445}

Рівняння в похідних похідних можна класифікувати за багатьма
характеристиками. Класифікація рівнянь важлива, оскільки, як виявилося,
кожен клас має свою загальну теорію та методи розв'язання рівнянь.

Тут ми наведемо шість основних методів класифікації

рівняння.

\begin{enumerate}
\def\labelenumi{\arabic{enumi}.}
\tightlist
\item
  Порядок рівняння. Порядок рівняння --- це найвищий порядок часткових
  похідних, включених до рівняння. Наприклад,
\end{enumerate}

Примітка. Ед.

1) Точніше, це метод лінеаризації нелінійних задач. Загальна теорія
збурень містить лінійні та нелінійні теорії збурень, кожна з яких, у
свою чергу, є регулярною та особливою. --- Ед. Ед.

Ці методи застосовні до спеціальних задач, що зумовлено існуванням
відповідних виразів --- функціоналів для цих задач.

Застосовувався раніше

\begin{verbatim}
u_t = u_{xx}
\end{verbatim}

(рівняння другого порядку), \texttt{u\_t\ =\ u\_x} (рівняння першого
порядку), \texttt{u\_1\ =\ uu\_\{xxx\}\ +\ \textbackslash{}sin\ x}
(рівняння трьох порядків).

\begin{enumerate}
\def\labelenumi{\arabic{enumi}.}
\setcounter{enumi}{1}
\tightlist
\item
  Кількість змінних. Кількість змінних --- це кількість незалежних
  змінних. Наприклад,
\end{enumerate}

\begin{verbatim}
u_t = u_{xx}
\end{verbatim}

\texttt{(рівняння\ з\ двома\ змінними\ \textquotesingle{}x\textquotesingle{}\ і\ \textquotesingle{}t\textquotesingle{}\ ),\ \textquotesingle{}u\_t\ =\ u\_\{rr\}\ +\ \textbackslash{}frac\{1\}\{r\}\ u\_r\ +\ \textbackslash{}frac\{1\}\{r\^{}2\}\ u\_\{\textbackslash{}theta\textbackslash{}theta\}\textquotesingle{}\ (рівняння\ з\ трьома\ змінними\ \textquotesingle{}r\textquotesingle{}\ ,}

\begin{enumerate}
\def\labelenumi{\arabic{enumi}.}
\setcounter{enumi}{2}
\tightlist
\item
  Лінійність. Рівняння з частковими похідними можуть бути лінійними та
  нелінійними. У лінійних рівняннях залежна змінна та всі її часткові
  похідні включаються лінійно, зокрема, вони не множаться одна на одну,
  не множаться в квадраті тощо). Точніше, лінійне рівняння другого
  порядку з двома незалежними змінними є рівнянням вигляду
\end{enumerate}

\begin{verbatim}
(1.1) Au_{xx} + Bu_{xy} + Cu_{yy} + Du_x + Eu_y + Fu = G
\end{verbatim}

де A, B, C, D, E, F і G --- це константи або задані функції
пояснювальних змінних x і y.

\begin{verbatim}
u_{tt} = e^{-t}u_{xx} + \sin t
\end{verbatim}

(лінійне рівняння), \texttt{uu\_\{xx\}\ +\ u\_t\ =\ 0} (нелінійне
рівняння), \texttt{u\_\{xx\}\ +\ yu\_\{yy\}\ =\ 0} (лінійне рівняння),
\texttt{xu\_x\ +\ yu\_y\ +\ u\^{}2\ =\ 0} (нелінійне рівняння).

\begin{itemize}
\tightlist
\item
  \begin{enumerate}
  \def\labelenumi{\arabic{enumi}.}
  \setcounter{enumi}{3}
  \tightlist
  \item
    Однорідність. Рівняння (1.1) називають однорідним, якщо права
    частина \texttt{G(x,\ y)} ідентична нулю для всіх x і y. Якщо
    \texttt{G(x,\ y)} не ідентичне нулю, то рівняння називають
    неоднорідним.
  \end{enumerate}
\item
  \begin{enumerate}
  \def\labelenumi{\arabic{enumi}.}
  \setcounter{enumi}{4}
  \tightlist
  \item
    Типи шансів. Якщо коефіцієнти A, B, C, D, E і F рівняння (1.1) є
    сталими, то це рівняння називається рівнянням зі сталими
    коефіцієнтами (інакше рівнянням зі змінними коефіцієнтами).
  \end{enumerate}
\item
  \begin{enumerate}
  \def\labelenumi{\arabic{enumi}.}
  \setcounter{enumi}{5}
  \tightlist
  \item
    Три основні типи лінійних рівнянь. Усі лінійні рівняння з частковими
    похідними другого порядку виду (1.1) належать до одного з трьох
    типів: а) параболічних, б) гіперболічних, в) еліптичних.
  \end{enumerate}
\end{itemize}

Параболічний тип. Рівняння параболічного типу описують процеси
теплопровідності та дифузії і визначаються умовою
\texttt{B\^{}2\ -\ 4AC\ =\ 0} .

Гіперболічний тип. Рівняння гіперболічного типу описів

осцилюючі системи та хвильові рухи визначаються умовою
\texttt{B\^{}2-4AC\textgreater{}0} .

Еліпічний тип. Рівняння еліптичного типу описують усталені процеси і
визначаються умовним \texttt{B\^{}2-4AC\textless{}0} . Приклади

\begin{enumerate}
\def\labelenumi{\alph{enumi})}
\tightlist
\item
\end{enumerate}

\begin{verbatim}
u_i = u_{xx}
\end{verbatim}

, \texttt{B\^{}2\ -\ 4AC\ =\ 0} (параболічна), (гіперболічна),
(гіперболічна), (гіперболічна \texttt{y\ \textgreater{}\ 0} ),
(гіперболічна), (еліптична), (еліптична), параболічна на
\texttt{y\ =\ 0} , гіперболічна на \texttt{y\ \textless{}\ 0} .

(У випадку змінних коефіцієнтів тип рівняння може змінюватися від точки
до точки.)

\section{ЗАУВАЖЕННЯ}\label{ux437ux430ux443ux432ux430ux436ux435ux43dux43dux44f}

\begin{itemize}
\tightlist
\item
  \begin{enumerate}
  \def\labelenumi{\arabic{enumi}.}
  \tightlist
  \item
    Загалом, значення \texttt{B\^{}2-4AC} є функцією незалежних змінних.
    Відповідно, базовий тип рівняння може змінюватися в межах визначення
    рівняння, хоча це не є обов'язковим.
  \end{enumerate}
\item
  \begin{enumerate}
  \def\labelenumi{\arabic{enumi}.}
  \setcounter{enumi}{1}
  \tightlist
  \item
    У рівнянні (1.1) незалежними змінними є x і y. У багатьох задачах
    однією з двох незалежних змінних є час, і рівняння (1.1) можна
    записати через x і t.
  \end{enumerate}
\end{itemize}

\pandocbounded{\includegraphics[keepaspectratio]{images/_page_10_Figure_9.jpeg}}

PHC 1.1. Діаграма класифікації диференціальних рівнянь у похідних
похідних.

\begin{enumerate}
\def\labelenumi{\arabic{enumi}.}
\setcounter{enumi}{2}
\tightlist
\item
  Рис. 1.1. Представлено діаграму класифікації для диференціальних
  рівнянь у похідних похідних.
\end{enumerate}

\section{ЗАВДАННЯ}\label{ux437ux430ux432ux434ux430ux43dux43dux44f}

\begin{itemize}
\tightlist
\item
  \begin{enumerate}
  \def\labelenumi{\arabic{enumi}.}
  \tightlist
  \item
    Класифікуємо наступні рівняння відповідно до всіх характеристик,
    зазначених у днаграмі на рис. 1.1:
  \end{enumerate}
\item
  \begin{enumerate}
  \def\labelenumi{\alph{enumi})}
  \tightlist
  \item
    \texttt{u\_t\ =\ u\_\{xx\}\ +\ 2u\_x\ +\ u} ,
  \end{enumerate}
\end{itemize}

\texttt{6)\ \textquotesingle{}u\_t\ =\ u\_\{xx\}\ +\ e\^{}\{-t\}\textquotesingle{}\ ,\ B)\ \textquotesingle{}u\_\{xx\}\ +\ 3u\_\{xy\}\ +\ u\_\{yy\}\ =\ \textbackslash{}sin\ x\textquotesingle{}\ ,\ \textquotesingle{}u\_\{tt\}\ =\ uu\_\{xxx\}\ +\ e\^{}\{-t\}\textquotesingle{}\ .}

\texttt{-\ 2.\ Скільки\ розв\textquotesingle{}язків\ \textquotesingle{}u\_1\ =\ u\_\{xx\}\textquotesingle{}\ рівняння\ існує?\ Спробуйте\ знайти\ рішення,\ як\ \textquotesingle{}u(x,\ t)\ =\ e\^{}\{a\textbackslash{}lambda\ +\ bt\}\textquotesingle{}\ .}
- 3. Якщо функції \texttt{u\_1(x,\ t)} і \texttt{u\_2(x,\ t)}
задовольняють рівняння (1.1), чи задовольняє сума цих функцій його?
Доведи це. - 4. Ймовірно, найпростіше диференціальне рівняння в похідних
---

\begin{verbatim}
\frac{\partial u(x, y)}{\partial x} = 0.
\end{verbatim}

Чи можете ви розв'язати це рівняння? (Знайдіть усі функції, що
задовольняють це рівняння.)

\begin{enumerate}
\def\labelenumi{\arabic{enumi}.}
\setcounter{enumi}{4}
\tightlist
\item
  Що можна сказати про рівняння
\end{enumerate}

\begin{verbatim}
\frac{\partial^2 u(x, y)}{\partial x \partial y} = 0
\end{verbatim}

?

Чи можете ви знайти всі розв'язки цього рівняння? Скільки їх? Порівняйте
з кількістю розв'язків звичайного диференціального рівняння

\begin{verbatim}
\frac{d^2y}{dx^2} = 0.
\end{verbatim}

\section{ЗАДАЧІ ТИПУ ДИФУЗІЇ (ПАРАБОЛІЧНІ
РІВНЯННЯ)}\label{ux437ux430ux434ux430ux447ux456-ux442ux438ux43fux443-ux434ux438ux444ux443ux437ux456ux457-ux43fux430ux440ux430ux431ux43eux43bux456ux447ux43dux456-ux440ux456ux432ux43dux44fux43dux43dux44f}

МЕТА ЛЕКЦІЇ: Показати, як параболічні рівняння використовуються для
розв'язання задач теплопровідності та дифузії. Наведено кілька прикладів
параболічних рівнянь, а також пояснюється фізичне значення різних членів
\texttt{(u\_1,\ u\_2,\ u\_3,\ u\_4,\ u\_6)} , включених у ці рівняння.

На конкретному прикладі розглядається основна ідея зміщеної задачі для
параболічного рівняння. Одна з головних цілей цієї лекції --- надати
читачеві інтуїтивне уявлення

Розуміння параболічних проблем.

Почнемо з розгляньмо просту фізичну задачу і те, як її можна описати
математичною моделлю з рівнянням у похідних похідних.

Потім ми ускладнимо задачу і покажемо, як нові фізичні умови призводять
до появи нових рівнянь у похідних похідних. У цій лекції диференціальні
рівняння не виводяться і не розв'язуються. Про це ми поговоримо пізніше.

\section{Простий експеримент з
теплопровідністю}\label{ux43fux440ux43eux441ux442ux438ux439-ux435ux43aux441ux43fux435ux440ux438ux43cux435ux43dux442-ux437-ux442ux435ux43fux43bux43eux43fux440ux43eux432ux456ux434ux43dux456ux441ux442ux44e}

Припустимо, що ми провели простий експеримент, який розбили на такі
кроки:

КРОК 1. Візьміть достатньо довгий (скажімо,
\texttt{L=2\textbackslash{},\textbackslash{}mathrm\{m\}} ) стрижень
(наприклад, мідний) діаметром
\texttt{2\textbackslash{},\textbackslash{}mathrm\{cm\}} , бічна поверхня
(але не кінці) якої термічно ізольована. Замість стрижня можна
використовувати мідну трубку, покриту теплоізоляцією зовні та зсередини.

Іншими словами, тепло може проходити лише через кінці стрижня і не може
проходити через його бічну поверхню.

КРОК 2. Тепер помістимо цей стрижень у пристрій із фіксованою
температурою \texttt{T\_\{\textbackslash{}mathfrak\{o\}\}} (градуси
Цельсія) на достатньо довгий час, щоб температура всередині стрижня
стала

Так само, як і температура пристрою. Для простоти припустимо, що
температура пристрою \texttt{T\_0\ =\ 10\^{}\{\textbackslash{}circ\}} C.

КРОК 3. Давайте витягнемо стрижень з нашого пристрою у певний момент
часу, що зручно взяти як нуль t = 0, і приєднаємо до нього дві термопари
з кінців. Завдання цих елементів --- підтримувати фіксовані температури
\texttt{T\_1} і \texttt{T\_2} на кінцях (наприклад,
\texttt{T\_1=0\^{}\{\textbackslash{}circ\}} C і
\texttt{T\_2=50\^{}\{\textbackslash{}circ\}} C). Іншими словами, два
термостати постійно контролюють температуру на кінцях стрижня, і якщо
температура відхиляється від встановлених значень \texttt{T\_1} і
\texttt{T\_2} , активуються потужні нагрівальні або охолоджувальні
елементи (пристрої) для відповідного регулювання температури. Наш
експеримент ілюструє RPS, 2 1,

\pandocbounded{\includegraphics[keepaspectratio]{images/_page_13_Figure_4.jpeg}}

Рис. 2.1. Схема експерименту: a --- це термопара, яка підтримує
встановлену температуру \texttt{T\_1} на лівому кінці стрижня;
\texttt{\textbackslash{}delta} --- це термопара, яка підтримує
встановлену температуру \texttt{T\_2} на правому кінці стрижня.

КРОК 4. Тепер ми будемо слідкувати за температурним профілем у стрижні
за допомогою пристрою візуалізації. Чому ми хочемо провести такий
експеримент? Це інше питання, про яке ми поговоримо пізніше.

Це кінець опису нашого експерименту.

Головна мета цієї лекції --- показати, як цю фізичну проблему та її
модифікації можна змоделювати мовою параболічних рівнянь.

\section{Математична модель
теплопровідності}\label{ux43cux430ux442ux435ux43cux430ux442ux438ux447ux43dux430-ux43cux43eux434ux435ux43bux44c-ux442ux435ux43fux43bux43eux43fux440ux43eux432ux456ux434ux43dux43eux441ux442ux456}

Опис нашої фізичної проблеми вимагає трьох типів відносин.

\begin{itemize}
\tightlist
\item
  \begin{enumerate}
  \def\labelenumi{\arabic{enumi}.}
  \tightlist
  \item
    Рівняння з частковими похідними, які описують фізичне явище
    теплопровідності.
  \end{enumerate}
\item
  \begin{enumerate}
  \def\labelenumi{\arabic{enumi}.}
  \setcounter{enumi}{1}
  \tightlist
  \item
    Граничні умови, що описують процеси теплопередачі на межах.
  \end{enumerate}
\item
  \begin{enumerate}
  \def\labelenumi{\arabic{enumi}.}
  \setcounter{enumi}{2}
  \tightlist
  \item
    Початкові твердження, які описують стан системи на початку процесу.
  \end{enumerate}
\end{itemize}

\section{Рівняння
теплопровідності}\label{ux440ux456ux432ux43dux44fux43dux43dux44f-ux442ux435ux43fux43bux43eux43fux440ux43eux432ux456ux434ux43dux43eux441ux442ux456}

Базове одновимірне рівняння теплопровідності записується у вилці

(2.1)

\begin{verbatim}
u_t = \alpha^2 u_{xx}, \quad 0 < x < L, \quad 0 < t < \infty.
\end{verbatim}

Це рівняння пов'язує значення \texttt{u\_t} --- швидкість зміни
температури з часом (виміряна в градусах/с)

та величину \texttt{u\_\{xx\}} --- це увігнутий профілю температури
\texttt{u(x,\ t)} (який є мірою різниці між температурою в заданій точці
та температурою в сусідніх точках).

\pandocbounded{\includegraphics[keepaspectratio]{images/_page_14_Figure_14.jpeg}}

Рис. 2.2. Стрілки показують зміну температури відповідно до рівняння
\texttt{u\_t\ =\ \textbackslash{}alpha\^{}2\ u\_\{xx\}} .

\texttt{Це\ рівняння\ буде\ виведено\ з\ базового\ рівняння\ збереження\ тепла\ в\ наступних\ лекціях,\ але\ наразі\ ми\ розглянемо\ його\ заданим.\ Це\ рівняння\ каже,\ що\ температура\ \textquotesingle{}u\textbackslash{}left(x,\textbackslash{}\ t\textbackslash{}right)\textquotesingle{}\ (у\ певній\ точці\ t\ у\ певній\ точці\ x\ стрижня)\ зростає\ \textquotesingle{}(u\_t\textgreater{}0)\textquotesingle{}\ або\ зменшується\ \textquotesingle{}(u\_t\textless{}0)\textquotesingle{}\ залежно\ від\ того,\ чи\ є\ друга\ похідна\ \textquotesingle{}u\_\{xx\}\textquotesingle{}\ позитивною\ чи\ негативною.\ Фіг.\ Рисунок\ 2.2\ ілюструє\ зміну\ температури\ в\ різних\ точках\ на\ елементі.}

Давайте подивимось, як значення \texttt{u\_\{xx\}} можна інтерпретувати
мовою теплопровідності. Припустимо, ми наблизимо значення
\texttt{u\_\{xx\}} фінальних постів

\begin{verbatim}
u_{xx}(x, t) \cong \frac{1}{\Delta x^2} [u(x + \Delta x, t) - 2u(x, t) + u(x - \Delta x, t)].
\end{verbatim}

Це співвідношення можна переписати у вигляді

\begin{verbatim}
u_{\lambda x}(x, t) \simeq \frac{2}{\Delta x^2} \left[ \frac{u(x + \Delta x, t) + u(x - \Delta x, t)}{2} - u(x, t) \right].
\end{verbatim}

Тепер ми можемо дати наступну інтерпретацію значення \texttt{u\_\{xx\}}
.

\begin{enumerate}
\def\labelenumi{\arabic{enumi}.}
\item
  Якщо температура \texttt{u(x,\ t)} менша за середнє значення
  температури у двох сусідніх точках, то
  \texttt{u\_\{xx\}\ \textgreater{}\ 0} (тут загальний тепловий потік по
  осі x додатний).
\item
  Якщо температура \texttt{u(x,\ t)} дорівнює середньому значенню
  температур у двох сусідніх точках, то \texttt{u\_\{xx\}\ =\ 0} (тут
  загальний тепловий потік по осі x дорівнює нулю).
\end{enumerate}

\texttt{3.\ Якщо\ температура\ \textquotesingle{}u\textbackslash{}left(x,\textbackslash{}\ t\textbackslash{}right)\textquotesingle{}\ вища\ за\ середню\ температуру\ у\ двох\ сусідніх\ точках,\ то\ \textquotesingle{}u\_\{xx\}\textless{}0\textquotesingle{}\ (тут\ —\ повний\ потік}

тепла вздовж осі x є від'ємною).

\texttt{Усе\ це\ ілюструється\ на\ рис.\ 2.2.\ Іншими\ словами,\ якщо\ температура\ в\ точці\ x\ більша\ за\ середню\ температуру\ у\ двох\ сусідніх\ точках\ \textquotesingle{}x\ -\ \textbackslash{}Delta\ x\textquotesingle{}\ і\ \textquotesingle{}x\ +\ \textbackslash{}Delta\ x\textquotesingle{}\ ,\ то\ температура\ в\ точці\ x\ зменшиться.\ Відповідно,\ точна\ швидкість\ зниження\ температури,\ а\ саме\ величина\ \textquotesingle{}u\_t\textquotesingle{}\ ,\ пропорційна\ цій\ різниці.\ Коефіцієнт\ пропорційності\ \textquotesingle{}\textbackslash{}alpha\^{}2\textquotesingle{}\ визначається\ властивостями\ матеріалу.\ Ми\ розглянемо\ цю\ константу\ детальніше\ у\ майбутніх\ лекціях.}

\section{Граничні умови
(ГІ)}\label{ux433ux440ux430ux43dux438ux447ux43dux456-ux443ux43cux43eux432ux438-ux433ux456}

Усі фізичні проблеми характеризуються наявністю певних меж. Отже, щоб
рухатися вперед, ми повинні включити межі в математичну модель, щоб
адекватно описати фізичну проблему. У нашому експерименті граничні умови
(ГІ) дуже прості. Оскільки температура на кінцях x=0 і x=L завжди
залишається сталою і дорівнює \texttt{T\_1} та \texttt{T\_2} , можливо
записати

(2.2)

\begin{verbatim}
\left\{ \begin{array}{l} u\left(0,\ t\right) = T_{t}, \\ u\left(L,\ t\right) = T_{2}, \end{array} \right. \quad 0 < t < \infty,
\end{verbatim}

\section{Початкові умови
(NU)}\label{ux43fux43eux447ux430ux442ux43aux43eux432ux456-ux443ux43cux43eux432ux438-nu}

Усі фізичні процеси мають починатися у певний момент часу (зазвичай
приймається як нуль t=0), тому ми повинні встановити фізичні умови саме
в цю точку. Відтоді, як ми почали контролювати температуру з моменту,
коли паличка мала постійну температуру
\texttt{T\_\{\textbackslash{}mathrm\{\ \}o\}} , ми

(2.3)

\begin{verbatim}
u(x, 0) = T_0, \quad 0 \le x \le L
\end{verbatim}

Тепер ми створили математичну модель експерименту. Система співвідношень
(2.1), (2.2) та (2.3) називається задачею зміщення для рівняння
теплопровідності і зазвичай записується як

(2.4)

\begin{verbatim}
\begin{aligned} & (\text{YUII}) \quad u_t = \alpha^2 u_{xx}, & 0 < x < L, \ 0 < t < \infty, \\ & \left( \Gamma \text{Y} \right) \quad \begin{cases} u \ (0, \ t) = T_1, \\ u \ (L, \ t) = T_2, \end{cases} & 0 < t < \infty, \\ & (\text{HY}) \quad u \ (x, \ 0) = T_0, \quad 0 \leqslant x \leqslant L. \end{aligned}
\end{verbatim}

Цікаво, що це зовсім не очевидно: чи існує одна функція
\texttt{u(x,\ t)}, яка задовольняє умови (2.4), і чи описує ця функція
температуру в стрижні? Отже, наше безпосереднє завдання --- знайти цю
функцію \texttt{u(x,\ t)} з (2.4).

Перш ніж завершити цю лекцію, давайте розглянемо деякі варіації основної
проблеми. Почнемо з невеликих змін рівняння теплопровідності.

\section{Деякі рівняння дифузійного
типу}\label{ux434ux435ux44fux43aux456-ux440ux456ux432ux43dux44fux43dux43dux44f-ux434ux438ux444ux443ux437ux456ux439ux43dux43eux433ux43e-ux442ux438ux43fux443}

Теплопередача через бічну поверхню пропорційна різниці температур

Рівняння

\begin{verbatim}
u_t = \alpha^2 u_{xx} - \beta (u - u_0)
\end{verbatim}

Описує теплопровідність елемента, враховуючи не лише дифузію теплових
\texttt{\textbackslash{}alpha\^{}2\ u\_\{xx\}} вздовж елемента, а й
теплопередачу через бічну поверхню елемента. Відтік тепла (
\texttt{\textbackslash{}beta\ \textgreater{}\ 0} ) або його приплив (
\texttt{\textbackslash{}beta\ \textless{}\ 0} ) пропорційний різниці між
температурою u стрижня та температурою навколишнього середовища
\texttt{u\_0} ( \texttt{\textbackslash{}beta} --- це сталий коефіцієнт
пропорційності). Якщо коефіцієнт \texttt{\textbackslash{}beta} великий
порівняно з \texttt{\textbackslash{}alpha\^{}2} , то теплопотік уздовж
стрижня буде малим порівняно з

Протікання через бічну поверхню, і тому тепло буде проходити через неї
відповідно до закону
\texttt{u\_t\ =\ -\textbackslash{}beta\ (u\ -\ u\_a)} .

У хімії, де величина відіграє роль концентрації, рівняння

\begin{verbatim}
u_1 = \alpha^2 u_{xx} - \beta (u - u_0)
\end{verbatim}

означає, що швидкість зміни \texttt{u\_i} кількості речовини залежить як
від дифузії \texttt{\textbackslash{}alpha\^{}2\ u\_\{xx\}} (у напрямку
осі x), так і від появи ( \texttt{\textbackslash{}beta\ \textless{}\ 0}
) або розпаду ( \texttt{\textbackslash{}beta\ \textgreater{}\ 0} )
речовини в хімічній реакції, і пропорційна різниці між двома
концентраціями u і \texttt{u\_\{\textbackslash{}theta\}} .

\section{Внутрішнє джерело
тепла}\label{ux432ux43dux443ux442ux440ux456ux448ux43dux454-ux434ux436ux435ux440ux435ux43bux43e-ux442ux435ux43fux43bux430}

Неоднородне рівняння

\begin{verbatim}
u_t = \alpha^2 u_{xx} + f(x, t)
\end{verbatim}

відповідає випадку, коли всередині елемента є джерело тепла (розташоване
вздовж усього елемента і діє в будь-який момент t). Наприклад, може бути
так, що дріт з електричним струмом проходить через стрижень і опір є
постійним джерелом тепла \texttt{f(x,t)} = K.

\section{Рівняння конвективної
дифузії}\label{ux440ux456ux432ux43dux44fux43dux43dux44f-ux43aux43eux43dux432ux435ux43aux442ux438ux432ux43dux43eux457-ux434ux438ux444ux443ux437ux456ux457}

\texttt{Припустимо,\ що\ домішка\ поширюється\ вздовж\ потоку,\ що\ рухається\ зі\ швидкістю\ v.\ Очевидно,\ що\ концентрація\ \textquotesingle{}u\textbackslash{}left(x,\textbackslash{},t\textbackslash{}right)\textquotesingle{}\ домішок\ змінюється\ залежно\ від\ x\ (позитивний\ напрямок\ осі\ x\ обирається\ вздовж\ потоку)\ і\ часу\ t.\ Швидкість\ зміни\ концентрації\ \textquotesingle{}u\_t\textquotesingle{}\ визначається\ рівнянням\ конвективної\ дифузії}

\begin{verbatim}
u_t = \alpha^2 u_{xx} - v u_x.
\end{verbatim}

Термін \texttt{\textbackslash{}alpha\^{}2\ u\_\{xx\}} описує внесок
дифузії, а \texttt{vu\_x} --- конвективний компонент. Що важливіше,
дифузія або конвекція залежить від співвідношення
\texttt{\textbackslash{}alpha\^{}2} до v. Димові частинки піднімаються
конвективно через рух гарячого повітря і одночасно розсіюють через
вихорові рухи повітря. Окрім цих модифікацій рівняння теплопровідності,
можливо змінювати граничні умови на кінцях стрижня відповідно до
фізичної ситуації. Деякі з цих змін ми обговоримо у третій лекції.

\section{ЗАУВАЖЕННЯ}\label{ux437ux430ux443ux432ux430ux436ux435ux43dux43dux44f-1}

Рівняння теплопровідності
\texttt{u\_t\ =\ \textbackslash{}alpha\^{}2(x)\ u\_\{xx\}} з змінним
коефіцієнтом \texttt{\textbackslash{}alpha(x)} відповідає задачі, в якій
дифузія залежить від x (тобто матеріал є гетерогенним). Наприклад, мідь
і сталь---

\pandocbounded{\includegraphics[keepaspectratio]{images/_page_18_Picture_4.jpeg}}

PHC. 2.3.

Пластини контактують, як показано на рис. 2.3. Припустимо, що
температура \texttt{u(0,t)=0\^{}\{\textbackslash{}circ\}} C
підтримується на лівому боці міді, а температура
\texttt{u(L,t)=20\^{}\{\textbackslash{}circ\}} C --- з правого боку
сталі, тоді рівняння, що описує теплопровідність у цій системі, виглядає
так:

\begin{verbatim}
u_t = \alpha^2(x) \, u_{_{\! A \, X}}, \quad 0 < x < L,
\end{verbatim}

де

\begin{verbatim}
\alpha(x) = \begin{cases} \alpha_1 & \text{(коэффициент диффузии в меди)}, \quad 0 < x < L/2, \\ \alpha_2 & \text{(коэффициент диффузии в стали)}, \quad L/2 < x < L. \end{cases}
\end{verbatim}

\# ЗАВДАННЯ

\begin{enumerate}
\def\labelenumi{\arabic{enumi}.}
\tightlist
\item
  Якщо початкова температура в стрижні
\end{enumerate}

\begin{verbatim}
u(x, 0) = \sin \pi x, \quad 0 \leqslant x \leqslant 1,
\end{verbatim}

а граничні умови такі:

\begin{verbatim}
u(0, t) = 0,
\end{verbatim}

\texttt{u(1,\ t)\ =\ 0,} Тоді як поводитиметься температура в стрижні
\texttt{U(x,\ t)} при T \textgreater{} 0? 2. Припустимо, що в стрижні є
постійне внутрішнє джерело тепла, так що рівняння теплопровідності в
стрижні записується як

\begin{verbatim}
u_t = \alpha^2 u_{xx} + 1, \quad 0 < x < 1.
\end{verbatim}

Припустимо, що ми зафіксували температуру на межі з значеннями
\texttt{u(0,\ t)} = 0 і \texttt{u(1,\ t)} = 1. Якою буде стаціонарна
температура у стрижні? Іншими словами, який c-незалежний розподіл
\texttt{U(x,\ t)} перешкоджує \texttt{u(x,\ t)}?

\begin{enumerate}
\def\labelenumi{\arabic{enumi}.}
\setcounter{enumi}{2}
\tightlist
\item
  Припустимо, що металевий стрижень втрачає тепло через бічну поверхню
  згідно з рівнянням
\end{enumerate}

\begin{verbatim}
u_1 = \alpha^2 u_{xx} - \beta u
\end{verbatim}

, \texttt{0\ \textless{}\ x\ \textless{}\ 1} .

І припустимо, що ми підтримуємо кінці стрижня при температурах
\texttt{U(0,\ T)} = 1, \texttt{U(1,\ T)} = 1. Знайдіть стаціонарний
розподіл температури по стрижню. Побудуйте графік цього розподілу. У
яких моментах відбувається тепловий потік?

\texttt{4.\ Припустимо,\ що\ термоізольований\ металевий\ стрижень\ довжини\ L\ =\ 1\ має\ початкову\ температуру\ \textquotesingle{}\textbackslash{}sin(3\textbackslash{}pi\ x)\textquotesingle{}\ ,\ а\ фіксовану\ температуру\ —\ 0\ і\ \textquotesingle{}10\textbackslash{},\^{}\{\textbackslash{}circ\}\textquotesingle{}\ C\ на\ лівому\ та\ правому\ кінцях.}

\section{Лекція 3}\label{ux43bux435ux43aux446ux456ux44f-3}

\section{ГРАНИЧНІ УМОВИ В ЗАДАЧАХ ДИФУЗІЙНОГО
ТИПУ}\label{ux433ux440ux430ux43dux438ux447ux43dux456-ux443ux43cux43eux432ux438-ux432-ux437ux430ux434ux430ux447ux430ux445-ux434ux438ux444ux443ux437ux456ux439ux43dux43eux433ux43e-ux442ux438ux43fux443}

МЕТА ЛЕКЦІЇ: Показати, як виникають різні граничні умови у задачах
теплопровідності та дифузії, а також ввести важливе поняття теплового
потоку.

Розглядаються три найважливіші типи граничних умов:

\begin{itemize}
\tightlist
\item
  \begin{enumerate}
  \def\labelenumi{\arabic{enumi}.}
  \tightlist
  \item
    u = \texttt{g(t)} (температура задається на межі),
    \texttt{-\ \textquotesingle{}2.\ \textbackslash{}\ \textbackslash{}frac\{\textbackslash{}partial\ u\}\{\textbackslash{}partial\ n\}\ +\ \textbackslash{}lambda\ u\ =\ g\textbackslash{}left(t\textbackslash{}right)\textquotesingle{}\ (встановлення\ температури\ навколишнього\ середовища);}
  \end{enumerate}
\end{itemize}

n --- вектор зовнішньої нормалі до межі),
\texttt{3.\ \textbackslash{}frac\{\textbackslash{}partial\ u\}\{\textbackslash{}partial\ n\}\ =\ g(t)}
(визначений тепловий потік через межу).

У задачах теплопровідності зазвичай використовують три основні типи
граничних умов. Ця лекція наводить приклади того, як різні фізичні умови
експерименту призводять до різних типів граничних умов.

\section{Граничні умови першого виду (температура задається на
межі)}\label{ux433ux440ux430ux43dux438ux447ux43dux456-ux443ux43cux43eux432ux438-ux43fux435ux440ux448ux43eux433ux43e-ux432ux438ux434ux443-ux442ux435ux43cux43fux435ux440ux430ux442ux443ux440ux430-ux437ux430ux434ux430ux454ux442ux44cux441ux44f-ux43dux430-ux43cux435ux436ux456}

Розглянемо тепловий потік в одномірному стрижні, показаний на рис. 3 1,
і припустимо, що температура кінців змінюється відповідно до законів
\texttt{g\_1(t)} і \texttt{g\_2(t)} .

\pandocbounded{\includegraphics[keepaspectratio]{images/_page_20_Picture_4.jpeg}}

Рис. 31. Температура встановлюється на межі.

Як згадувалося в попередній лекції, кожен пристрій, що підтримує
встановлену температуру на кінцях, складається з термостата та
нагрівального елемента для підтримки відповідних теплових умов. Проблеми
з граничними умовами першого виду дуже поширені. У деяких випадках
складність полягає у пошуку контролерів граничних умов, тобто температур
\texttt{g\_1(t)} та \texttt{g\_2(t)} меж, які призводять до певної зміни
температури всередині стрижня. Наприклад, у металургії часто необхідно
обирати граничні умови контролю, щоб температура металу всередині печі
змінювалася з часом, а температурний градієнт був невеликим.

\pandocbounded{\includegraphics[keepaspectratio]{images/_page_20_Picture_7.jpeg}}

Рис. 3.2. Коливання температури на межі.

Граничні умови першого виду також виникають у багатовимірних задачах
теплопровідності. Як приклад можна навести цікаву задачу визначення
температурного поля всередині кругового диска з радіусом \emph{R} (див.
рис. 3.2), коли гранична температура задається у полярних координатах
формулою
\texttt{u(R,\ \textbackslash{}theta,\ t)\ =\ \textbackslash{}cos\ t\ \textbackslash{}sin\ \textbackslash{}theta}
.

Звісно, щоб розв'язати задачу, потрібно знати початкову температуру, але
в цьому випадку початкові умови перестають впливати на розв'язок через
дуже короткий час, і отримана температура всередині диска визначається
лише значеннями температури на межі.

\section{Граничні умови другого виду (задана температура навколишнього
середовища)}\label{ux433ux440ux430ux43dux438ux447ux43dux456-ux443ux43cux43eux432ux438-ux434ux440ux443ux433ux43eux433ux43e-ux432ux438ux434ux443-ux437ux430ux434ux430ux43dux430-ux442ux435ux43cux43fux435ux440ux430ux442ux443ux440ux430-ux43dux430ux432ux43aux43eux43bux438ux448ux43dux44cux43eux433ux43e-ux441ux435ux440ux435ux434ux43eux432ux438ux449ux430}

Припустимо, що ми знову розглядаємо термічно ізольований мідний
стержень, але тепер відмовляємося від заданого температурного режиму на
кінцях стрижня. Натомість давайте наблизимо кінці стрижня до контакту з
двома середовищами. Нехай температура одного з них \texttt{g\_1(t)} , а
іншого \texttt{g\_2(t)} . Іншими словами, припустимо, що лівий кінець
стрижня поміщений у контейнер із рідиною, температура якої змінюється
відповідно до закону \texttt{g\_1(t)} . А правий кінець поміщається в
іншу рідину з температурою \texttt{g\_2(t)} (рис. 3, 3)

\pandocbounded{\includegraphics[keepaspectratio]{images/_page_21_Picture_6.jpeg}}

Рис. 3.3. Конвективна теплопередача через межі, a --- це рідина при
температурі \texttt{g\_1(t)} , \texttt{\textbackslash{}delta} --- рідина
при температурі \texttt{g\_2(t)} .

Задаючи граничні умови такого типу, ми не можемо вважати, що граничні
температури стрижня збігаються з температурами рідини. Але ми знаємо
закон Ньютона: якщо температура одного з кінців стрижня менша за
температуру відповідної рідини, тепло потече в стрижень зі швидкістю,
пропорційною різницям температур. Іншими словами, для одномірного
стрижня з межами x=0 і x=1 закон теплопередачі Ньютона формулюється так:

(3.1)

\begin{verbatim}
\{
\end{verbatim}

\texttt{вихідний\ тепловий\ потік\ на\ \textquotesingle{}x=0\textquotesingle{}\ дорівнює\ \textquotesingle{}h\textbackslash{}left{[}u\textbackslash{}left(0,\textbackslash{},t\textbackslash{}right)-g\_\{1\}\textbackslash{}left(t\textbackslash{}right)\textbackslash{}right{]}\textquotesingle{}\ ,\ вихідний\ тепловий\ потік\ при\ \textquotesingle{}x=L\textquotesingle{}\ дорівнює\ \textquotesingle{}h\textbackslash{}left{[}u\textbackslash{}left(L,\textbackslash{},t\textbackslash{}right)-g\_\{2\}\textbackslash{}left(t\textbackslash{}right)\textbackslash{}right{]}\textquotesingle{}\ ,}

де h --- коефіцієнт теплопередачі, який показує, скільки калорій
проходить через межу за одну секунду при різниці температур у один
градус. Вихідний потік дорівнює кількості кало-

Пройшов крізь палочний спис за одну секунду. Зверніть увагу, що вихідний
тепловий потік буде позитивним на кінці стрижня, де температура стрижня
вища за температуру навколишнього середовища. Рівняння (3.1) разом із
відомим законом теплопровідності Фур'є тепер можна використовувати для
отримання Р и с

\pandocbounded{\includegraphics[keepaspectratio]{images/_page_22_Figure_3.jpeg}}

. 3.4. Ілюстрація закону Фур'є:
\texttt{\textbackslash{}frac\{\textbackslash{}partial\ u\}\{\textbackslash{}partial\ n\}\ \textgreater{}\ 0}
означає, що тепло надходить у цю область;
\texttt{\textbackslash{}frac\{\textbackslash{}partial\ u\}\{\textbackslash{}partial\ n\}\ \textless{}\ 0}
означає, що тепло виходить із цієї зони; n --- напрямок зовнішньої
нормалі;
\texttt{\textbackslash{}frac\{\textbackslash{}partial\ u\}\{\textbackslash{}partial\ n\}\ -}
Зміна температури в напрямку N;
\texttt{\textbackslash{}frac\{\textbackslash{}partial\ u\}\{\textbackslash{}partial\ n\}\ =\ -\textbackslash{}frac\{\textbackslash{}partial\ u\}\{\textbackslash{}partial\ (-n)\}}
є похідною у напрямку внутрішньої нормальної --- це похідна, взята
\texttt{\textbackslash{}mathbf\{e\}} протилежному знаку протилежного
знака пормалі.

граничні умови. Закон Фур'є дає друге представлення (перше міститься у
(3.1)) для вихідного теплового потоку. Прирівнюючи ці два вирази до
потоку, ми отримуємо граничні умови, які шукаємо. Сформулюємо закон
Фур'є (встановлений експериментально):

(3.2) Тепловий потік, що проходить через межу площі, пропорційний
нормальній похідній температури у напрямку внутрішньої норми.

Цей закон стверджує, що якщо температура швидко підвищується у напрямку
зовнішньої температури до межі D (рисунок 3.4), то потік буде текти з
навколишнього середовища до області D.

У нашій одновимірній задачі закон Фур'є має вигляд:

(3.3)

\begin{verbatim}
\begin{cases} \text{вытекающий поток тепла при } x=0 \text{ равен } k \frac{\partial u}{\partial x}, \\ \text{вытекающий поток тепла при } x=' \text{ равен } -k \frac{\partial u}{\partial x}, \end{cases}
\end{verbatim}

де k --- теплопровідність матеріалу, яка є мірою того, наскільки добре
матеріал проводить тепло. Погано провідні матеріали мають коефіцієнт
близький до нуля в одиницях GHS, тоді як мідь та алюміній мають значення
значно вище нуля.

\pandocbounded{\includegraphics[keepaspectratio]{images/_page_23_Figure_5.jpeg}}

Рис. 3.5. Ще одна ілюстрація закону Фур'є.

Закон Фур'є (3.3) фактично описує теплопровідність всюди всередині
стрижня, а не лише на його межі; Наприклад, (див. рис. 3.5).

(3.4) Потік, що проходить через точку

x

\texttt{зліва\ направо\ \textquotesingle{}=-k\textbackslash{}frac\{\textbackslash{}partial\ n\}\{\textbackslash{}partial\ x\}\textquotesingle{}\ Закон\ Фур\textquotesingle{}є\ (3.4)\ стверджує:\ якщо\ \textquotesingle{}u\_x\ \textless{}\ 0\textquotesingle{}\ ,\ то\ в\ точці\ \textquotesingle{}x\_0\textquotesingle{}\ тепло\ тече\ зліва\ направо,\ якщо\ \textquotesingle{}u\_x\ \textgreater{}\ 0\textquotesingle{}\ ,\ то\ в\ точці\ \textquotesingle{}x\_0\textquotesingle{}\ тепло\ тече\ справа\ наліво.\ (Тепло\ завжди\ переходить\ від\ вищих\ температур\ до\ нижчих.)}

Нарешті, якщо використовувати вирази (3.1) і (3.3) для теплового потоку,
то отримуємо бажані граничні умови для задачі, описаної на рис. 3.3, у
чисто математичній формі

\begin{verbatim}
(\Gamma Y) \quad \begin{cases} \frac{\partial u\left(0,\,t\right)}{\partial x} = \frac{h}{k}\left[u\left(0,\,t\right) - g_{1}\left(t\right)\right], \\ \frac{\partial u\left(L,\,t\right)}{\partial x} = -\frac{h}{k}\left[u\left(L,\,t\right) - g_{2}\left(t\right)\right]. \end{cases} \quad 0 < t < \infty,
\end{verbatim}

Дуже часто константа h/k позначається як \texttt{\textbackslash{}lambda}
, а граничні умови теплового потоку на межі записуються у вигляді

(3.5)

\begin{verbatim}
u_{x}(0, t) = \lambda [u(0, t) - g_{1}(t)], u_{x}(L, t) = -\lambda [u(L, t) - g_{2}(t)].
\end{verbatim}

Для випадку вищих розмірностей отримують подібні граничні умови.
Наприклад, якщо межа кругового диска осяяна рухомою рідиною з
температурою \texttt{g(\textbackslash{}theta,t)} , то граничні умови
записуються як

\begin{verbatim}
\frac{\partial u}{\partial r}(R, \theta, t) = -\frac{h}{k} [u(R, \theta, t) - g(\theta, t)].
\end{verbatim}

Тут
\texttt{\textbackslash{}frac\{\textbackslash{}partial\ u\}\{\textbackslash{}partial\ r\}(R,\ \textbackslash{}theta,\ t)}
позначає зовнішню нормальну похідну (у позитивному напрямку r-осі)
значення u, обчислену в точці \texttt{(R,\ \textbackslash{}theta)} межі.
Ми назвемо такі граничні умови лінійними (оскільки вони лінійні в U і
\texttt{u\_r} ), але неоднорідними, оскільки функція
\texttt{g(\textbackslash{}theta,\ t)} знаходиться праворуч.

\section{Граничні умови третього виду (враховуючи потік; зокрема,
враховується випадок термічно ізольованих
меж)}\label{ux433ux440ux430ux43dux438ux447ux43dux456-ux443ux43cux43eux432ux438-ux442ux440ux435ux442ux44cux43eux433ux43e-ux432ux438ux434ux443-ux432ux440ux430ux445ux43eux432ux443ux44eux447ux438-ux43fux43eux442ux456ux43a-ux437ux43eux43aux440ux435ux43cux430-ux432ux440ux430ux445ux43eux432ux443ux454ux442ux44cux441ux44f-ux432ux438ux43fux430ux434ux43eux43a-ux442ux435ux440ux43cux456ux447ux43dux43e-ux456ux437ux43eux43bux44cux43eux432ux430ux43dux438ux445-ux43cux435ux436}

Жоден потік не проходить через термоізольовані межі, тому нормальна
похідна (зовнішня чи внутрішня) має зрівнятися з нулем на межі, оскільки
потік пропорційний нормальній похідній. У випадку одномірного елемента з
ізольованими кінцями x=0 і x=L, граничні умови мають вигляд

\begin{verbatim}
u_{\lambda}(0, t) = 0,
\end{verbatim}

\texttt{\textquotesingle{}u\_\{\textbackslash{}lambda\}(L,\ t)\ =\ 0.\textquotesingle{}\ \textquotesingle{}0\ \textless{}\ t\ \textless{}\ \textbackslash{}infty,\textquotesingle{}\ У\ двовимірній\ області\ ізоляція\ межі\ означає,\ що\ нормальна\ похідна\ температури\ на\ межі\ перетворюється\ на\ нуль.\ Припустимо,\ наприклад,\ що\ круглий\ диск\ термічно\ ізольований\ уздовж\ межі.\ Тоді\ граничну\ умову\ слід\ записати\ так:\ \textquotesingle{}u\_r(R,\ \textbackslash{}theta,\ t)\ =\ 0\textquotesingle{}\ для\ всіх\ \textquotesingle{}0\ \textbackslash{}le\ \textbackslash{}theta\ \textless{}\ 2\textbackslash{}pi\textquotesingle{}\ і\ всіх\ \textquotesingle{}0\ \textless{}\ t\ \textless{}\ \textbackslash{}infty\textquotesingle{}\ .}

З іншого боку, якщо вказати кількість тепла, що проходить через межу
нашого диска, то граничні умови набувають вигляду

\begin{verbatim}
u_r(R, \theta, t) = f(\theta, t),
\end{verbatim}

де \texttt{f(0,\ t)} --- це кількість тепла, що переходить у круглий
диск із зовнішнього джерела.

Тепер проілюструємо різні типи граничних умов.

\section{Типові граничні умови для задачі одномірної
теплопровідності}\label{ux442ux438ux43fux43eux432ux456-ux433ux440ux430ux43dux438ux447ux43dux456-ux443ux43cux43eux432ux438-ux434ux43bux44f-ux437ux430ux434ux430ux447ux456-ux43eux434ux43dux43eux43cux456ux440ux43dux43eux457-ux442ux435ux43fux43bux43eux43fux440ux43eux432ux456ux434ux43dux43eux441ux442ux456}

Припустимо, що у нас є мідний стрижень довжиною 200 см, бікова поверхня
якого термічно ізольована, а початкова температура становить 0 °C.
Припустимо також, що верхній кінець стрижня x=0 термічно ізольований, а
нижній кінець x=200 змивається рухомою водою, яка має постійну
температуру \texttt{g\_2(t)=20} °C (див. рис. 3.6).

\pandocbounded{\includegraphics[keepaspectratio]{images/_page_25_Picture_6.jpeg}}

Рис. 3.6. Змішане завдання.

Математична модель цієї задачі описується наступними чотирма
відношеннями:

\begin{verbatim}
(3.6) \qquad \begin{array}{llll} (\mathrm{Y}\mathrm{H}\Pi) & u_t = \alpha^2 u_{xx}, & 0 < x < 200, & 0 < t < \infty, \\ u_x (0, \, t) = 0, & \\ u_x (200, \, t) = -\frac{h}{k} \left[ u \, (200, \, t) - 20 \right], & 0 < t < \infty, \\ (\mathrm{HY}) & u \, (x, \, 0) = 0 \, ^{\mathrm{o}}\mathrm{C}, & 0 \leqslant x \leqslant 200, & \end{array}
\end{verbatim}

де \texttt{\textbackslash{}alpha\^{}2\ =\ 1,16} см 2/с --- теплова
дифузія міді;

k = 0,93 кал/см·с·°C --- теплопровідність міді;

h --- коефіцієнт теплопередачі. Визначення цього співвідношення ---
велика проблема. Він вимірюється швидкістю виходу тепла від другого
кінця стрижня до води. Це залежить від швидкості потоку води, стану
поверхні, яку слід змивати тощо.

Читач може спробувати провести експеримент, щоб визначити

ділення значення цього коефіцієнта.

\section{ЗАУВАЖЕННЯ}\label{ux437ux430ux443ux432ux430ux436ux435ux43dux43dux44f-2}

\begin{enumerate}
\def\labelenumi{\arabic{enumi}.}
\tightlist
\item
  Типова задача внутрішньої теплопровідності для квадрата показана на
  рис. 3.7.
\end{enumerate}

\texttt{У\ цій\ задачі\ початкова\ температура\ \textquotesingle{}u\textbackslash{}left(x,\textbackslash{},y,\textbackslash{},t\textbackslash{}right)\textquotesingle{}\ при\ t\ =\ 0\ задається\ всередині\ квадрата,\ і\ рівняння\ з\ похідними\ з\ частковими\ має\ вигляд\ Р\ и\ с}

\pandocbounded{\includegraphics[keepaspectratio]{images/_page_26_Figure_9.jpeg}}

. 3.7. Типові граничні умови для задачі дифузії в квадраті: a ---
кількість тепла, \texttt{f\_1(t)} проходить через межу; b ---
термоізольована межа; a --- це навколишня температура Равії
\texttt{10\^{}\{\textbackslash{}circ\}} C; z --- температура
підтримується рівною \texttt{10\^{}\{\textbackslash{}circ\}} C.

а граничні умови будуть виконані на
\texttt{0\ \textless{}\ t\ \textless{}\ \textbackslash{}infty} і
визначатимуть температуру в будь-який момент. Якою б не була ця
температура, вона повинна відповідати граничним умовам, зазначеним на
рис. 3.7.

\begin{enumerate}
\def\labelenumi{\arabic{enumi}.}
\setcounter{enumi}{1}
\tightlist
\item
  Зверніть увагу, що гранична умова
\end{enumerate}

\begin{verbatim}
u_r(R, \theta, t) = -\frac{h}{k} [u(R, \theta, t) - g(\theta, t)]
\end{verbatim}

На колі не означає, що гранична температура дорівнює
\texttt{g(\textbackslash{}theta,t)} , але якщо коефіцієнт теплопередачі
h великий, то ці умови фактично еквівалентні вимогі
\texttt{g(\textbackslash{}theta,t)} температури на межі.

\section{ЗАВДАННЯ}\label{ux437ux430ux432ux434ux430ux43dux43dux44f-1}

\begin{enumerate}
\def\labelenumi{\arabic{enumi}.}
\item
  Намалюйте контури графіків для розв'язання змішаної задачі (3.6) для
  різних моментів часу. Чи відповідають ваші графи визначеним граничним
  умовам? Яка стаціонарна температура стрижня? Очевидно, ваша відповідь
  базується на інтуїції?
\item
  Як би ви інтерпретували наступне змішане завдання?
\end{enumerate}

\begin{verbatim}
\begin{aligned}
(\text{УЧП})\quad & u_t = \alpha^2 u_{xx}, && 0 < x < 1,\ 0 < t < \infty, \\
(\text{ГУ})\quad & u(0,t) = 0,\quad u_x(1,t) = 1, && 0 < t < \infty, \\
(\text{НУ})\quad & u(x,0) = \sin(\pi x), && 0 \leqslant x \leqslant 1.
\end{aligned}
\end{verbatim}

Чи можете ви графічно показати рішення цієї проблеми в різні моменти
часу? Чи буде це рішення, як правило, стаціонарним? Це очевидно?

\begin{enumerate}
\def\labelenumi{\arabic{enumi}.}
\setcounter{enumi}{2}
\tightlist
\item
  Яку фізичну інтерпретацію ви можете дати проблемі
\end{enumerate}

\begin{verbatim}
\begin{array}{llll} (\text{УЧП}) & u_t = \alpha^2 u_{xx}, & 0 < x < 1, & 0 < t < \infty, \\ (\text{ГУ}) & \begin{cases} u_x \, (0, \, t) = 0, & 0 < t < \infty, \\ u_x \, (1, \, t) = 0, & 0 < t < \infty, \end{cases} \\ (\text{НУ}) & u \, (x, \, 0) = \sin (\pi x), & 0 \leqslant x \leqslant 1? \\ \end{array}
\end{verbatim}

Чи можете ви графічно показати рішення цієї проблеми в різні моменти
часу?

Що можна сказати про стабільну температуру?

\begin{enumerate}
\def\labelenumi{\arabic{enumi}.}
\setcounter{enumi}{3}
\tightlist
\item
  Припустимо, що бічна поверхня металевого стрижня не ізольована і має
  початкову температуру 20°C, але температура одного кінця стає 50°C.
  Потім стрижень промивають рідиною з температурою 30°C. Як виглядає
  змішана задача, що відповідає цьому випадку?
\end{enumerate}

\section{ВИВЕДЕННЯ РІВНЯННЯ
ТЕПЛОПРОВІДНОСТІ}\label{ux432ux438ux432ux435ux434ux435ux43dux43dux44f-ux440ux456ux432ux43dux44fux43dux43dux44f-ux442ux435ux43fux43bux43eux43fux440ux43eux432ux456ux434ux43dux43eux441ux442ux456}

МЕТА ЛЕКЦІЇ: Показати, як можна отримати одномірне рівняння
теплопровідності

\begin{verbatim}
u_t = \alpha^2 u_{xx} + f(x, t),
\end{verbatim}

спираючись лише на фундаментальний закон збереження кількості тепла.
Розглядається залежність швидкості теплопередачі від основних
термофізичних параметрів: теплопровідності, теплоємності та густини
матеріалу. Розглядаються різні модифікації базового рівняння
теплопровідності.

У кожній галузі знань існує набір базових принципів, правильність яких
очевидна. Усі інші твердження мають бути виведені з цих базових
принципів. Звісно, те, що очевидно одній людині, може здатися сумнівним
для іншої. Історія будь-якої науки полягає у постійному пошуку більш
фундаментальних принципів, що лежать в основі цієї науки. Основні
принципи мають бути настільки універсальними, щоб ніхто не сумнівався.

? --- Припущення A --- Припущення B --- Припущення C --- Припущення D
--- ? Рис. 4.1. Аксіоматичний метод.

Наприклад, одна людина може вважати, що всі факти певної науки можна
вивести з одного базового припущення, яке ми позначимо як B. На основі
припущення B вона або вона може довести теорему C, що призведе до
теореми D, а яка, у свою чергу, призведе до доведення багатьох інших
теорем (див. рисунок 4.1). Ось як виглядає прогрес у науці --- дослідник
може отримувати все більше нових результатів. Найвидатніші фізики,
хіміки та біологи або саме на цьому шляху.

З іншого боку, замість доведення нових теорем можна спробувати, якщо
можливо, знайти новий принцип (назвемо його A), який є більш
фундаментальним за B. Припущення A природно вважається більш
фундаментальним, ніж B.

Якщо B можна вивести з A. Отже, якщо можна знайти новий принцип A, тоді
межа наших знань зміщується до сфери більш фундаментальних ідей про
певну науку. У теорії теплопровідності основним принципом є закон
збереження енергії (теплова енергія). Усі інші твердження випливають із
цього базового принципу (рис. 4.2).

Закон збереження енергії

\begin{verbatim}
u_t = \alpha^2 u_{xx} + f(x,t)
\end{verbatim}

Інші властивості теплопровідності Рис. 4.2. Енергозбереження: наріжний
камінь теплопровідності.

Звісно, ми могли б забути цю лекцію і використати рівняння
теплопровідності як фундаментальне поняття (дехто може так подумати,
якщо це очевидно), але тоді ми обдуримо більш серйозних студентів, адже
закон збереження вважається найфундаментальнішим принципом у науці.
Науковці, моделюючи свої задачі, часто записують відношення, що
виражають закони збереження, а потім переписують їх у вигляді рівнянь у
похідних похідних.

Повернемося до головної мети нашої лекції --- вивести рівняння
теплопровідності з рівняння збереження енергії.

\section{Виведення рівняння
теплопровідності}\label{ux432ux438ux432ux435ux434ux435ux43dux43dux44f-ux440ux456ux432ux43dux44fux43dux43dux44f-ux442ux435ux43fux43bux43eux43fux440ux43eux432ux456ux434ux43dux43eux441ux442ux456-1}

Розглянемо однорідний стрижень довжини L, щодо якого зробимо такі
припущення:

\begin{enumerate}
\def\labelenumi{\arabic{enumi}.}
\tightlist
\item
  Стрижень виготовлений з одного однорідного провідного матеріалу.
\end{enumerate}

\pandocbounded{\includegraphics[keepaspectratio]{images/_page_29_Figure_10.jpeg}}

Рис. 4.3. Тонкий теплопровідний стержень.

\begin{enumerate}
\def\labelenumi{\arabic{enumi}.}
\setcounter{enumi}{1}
\tightlist
\item
  Бікова поверхня стрижня термічно ізольована (тепло може поширюватися
  лише вздовж осі x).
\end{enumerate}

\texttt{3.\ Стрижень\ тонка,\ тобто\ температура\ всіх\ точок\ у\ кожному\ поперечному\ перерізі\ стрижня\ є\ сталою.\ Якщо\ розглядати\ частину\ вудки\ на\ сегменті\ \textquotesingle{}{[}x,\ x+\textbackslash{}Delta\ x{]}\textquotesingle{}\ і\ використовувати}

закон збереження кількості тепла, тоді можна написати:

\texttt{(4.1)\ Загальна\ зміна\ кількості\ тепла\ на\ сегменті\ \textquotesingle{}{[}x,\ x\ +\ \textbackslash{}Delta\ x{]}\ =\textquotesingle{}\ дорівнює\ загальній\ кількості\ тепла,\ що\ пройшло\ через\ межі\ +\ Нульова\ кількість\ тепла,\ утвореного\ всередині\ сегмента\ \textquotesingle{}{[}x,\ x\ +\ \textbackslash{}Delta\ x{]}\textquotesingle{}\ .}

\texttt{Оскільки\ загальна\ кількість\ тепла\ (у\ калоріях)\ всередині\ сегмента\ \textquotesingle{}\textbackslash{}\{x,\ x\ +\ \textbackslash{}Delta\ x\textbackslash{}\}\textquotesingle{}\ в\ будь-який\ момент\ часу\ t\ обчислюється\ формулою:}

\texttt{Загальна\ кількість\ тепла\ всередині\ сегмента\ \textquotesingle{}{[}x,\ x+\textbackslash{}Delta\ x{]}\ =\ \textbackslash{}int\textbackslash{}limits\_\{-\textbackslash{}infty\}\^{}\{x+\textbackslash{}Delta\ x\}\ coAu\textquotesingle{}\ (s,\ t)\ ds,}

де c --- питома теплоємність матеріалу (показує здатність матеріалу
накопичувати тепло),

Щільність матеріалу,

A --- це площа поперечного перерізу стрижня.

\textbf{va}con збереження енергії (4.1) можна надати у такій
математичній формі:

(4.2)

\begin{verbatim}
\frac{d}{dt} \int_{x}^{x+\Delta x} c\rho Au(s, t) ds = c\rho A \int_{x}^{x+\Delta x} u_{t}(s, t) ds = eA \left[ u_{x}(x+\Delta x, t) - u_{x}(x, t) \right] + A \int_{x+\Delta x}^{x+\Delta x} f(s, t) ds.
\end{verbatim}

\textbf{де} k --- теплопровідність матеріалу (вказує на здатність
матеріалу проводити тепло), \texttt{f(x,\ t)} --- об'ємна потужність
зовнішнього джерела тепла (калорії/см·с).

Задача полягає в тому, щоб записати рівняння (4.2) у формі, яка не
містить інтегралів. Щоб розв'язати цю задачу, нагадаємо читачеві теорему
про середнє значення з ходу інтегральної нечисельності.

\section{Теорема про
значення}\label{ux442ux435ux43eux440ux435ux43cux430-ux43fux440ux43e-ux437ux43dux430ux447ux435ux43dux43dux44f}

Якщо функція \texttt{f(x)} є неперервною на сегменті {[}a, b{]}, то
існує принаймні одна точка
\texttt{\textbackslash{}xi\ \textbackslash{}in\ {[}a,\ b{]}} така, що

\begin{verbatim}
\int_{a}^{b} f(x) dx = f(\xi) (b-a).
\end{verbatim}

Застосовуючи цей результат до рівняння (4.2), ми отримуємо наступний
результат

\begin{verbatim}
coAu_{t}(\xi, t) \Delta x = kA\{u_{x}(x + \Delta x, t) - u_{x}(x, t)\} + Af(\xi, t) \Delta x,
x < \xi < x + \Delta x,
\end{verbatim}

або

\begin{verbatim}
u_t(\xi, t) = \frac{k}{c\rho} \left\{ \frac{u_x(x + \Delta x, t) - u_x(x, t)}{\Delta x} \right\} + \frac{1}{c\rho} f(\xi, t).
\end{verbatim}

\begin{verbatim}
Спрямовуючи \Delta x до 0, отримуємо потрібне рівняння
\end{verbatim}

(4.3)

\begin{verbatim}
u_t(x, t) = \alpha^2 u_{xx}(x, t) + F(x, t),
\end{verbatim}

де

\begin{verbatim}
\alpha^{2}=\frac{k}{c\rho}
\end{verbatim}

\begin{verbatim}
— коефіцієнт теплової дифузії, F\left(x,\ t\right)=\frac{1}{c\rho}f\left(x,\ t\right) — густина джерел тепла.
\end{verbatim}

Ми могли б зупинитися на цьому. Однак ми також розглянемо випадок, коли
бічна поверхня стрижня не термічно ізольована. Припустимо, що кількість
теплового потоку через бічну поверхню стрижня в цьому випадку
пропорційна різниці між температурою стрижня \texttt{u(x,t)} і
температурою навколишнього середовища, яка залишається сталою і дорівнює
нулю. У цьому випадку закон збереження кількості тепла веде до рівняння

\begin{verbatim}
(4.4) u_t = \alpha^2 u_{xx} - \beta u + F(x, t),
\end{verbatim}

\begin{verbatim}
де \beta — коефіцієнт пропорційності потоку через бічну поверхню.
\end{verbatim}

\section{ЗАУВАЖЕННЯ}\label{ux437ux430ux443ux432ux430ux436ux435ux43dux43dux44f-3}

\begin{enumerate}
\def\labelenumi{\arabic{enumi}.}
\tightlist
\item
  Як зазначалося вище, константа k називається теплопровідністю речовини
  і чисельно дорівнює кількості тепла (у калоріях), що проходить за
  секунду через пластину товщиною 1 см і площею поперечного перерізу 1
  см³ при різниці температур на протилежних гранях пластини рівною 1°C.
  Числові значення теплопровідності для різних матеріалів можна знайти в
  довідниках з фізики або хімії. Типові значення теплопровідності
  варіюються від 1 (для міді) до значень, близьких до нуля для хороших
  теплоізоляторів. Якщо стрижень зроблений із однорідного матеріалу, то
  значення k не залежить від x. Для деяких матеріалів значення k
  залежить від температури, і тому рівняння теплопровідності в цьому
  випадку стає нелінійним:
\end{enumerate}

\begin{verbatim}
u_{t} = \frac{1}{c\rho} \frac{\partial}{\partial x} \{ k (u) | u_{x} \}.
\end{verbatim}

Однак найчастіше теплопровідність k дуже слабо впливає на температуру u,
і цю нелінійність можна ігнорувати.

\begin{enumerate}
\def\labelenumi{\arabic{enumi}.}
\setcounter{enumi}{1}
\item
  Значення c зазвичай називають питомою теплоємністю речовини. Він
  характеризує здатність речовини зберігати теплову енергію. Наприклад,
  запечена картопля має високу теплоємність, оскільки може зберігати
  велику кількість тепла на одиницю маси (саме тому вона зберігає високу
  температуру так довго). Чисельна теплоємність дорівнює кількості
  калорій тепла, які потрібно передати 1 граму речовини, щоб підвищити
  її температуру на 1°C. Для більшості розглянутих задач значення c
  можна вважати сталим, не записуючи ні x, ні i. Значення питомої
  теплоємності для конкретних матеріалів можна знайти у фізичних
  довідниках.
\item
  Одиниці вимірювання основних величин, включених у рівняння
\end{enumerate}

теплопровідність (у системі GHS) наступна:

i-температура (градуси Цельсія), \texttt{u\_t} --- швидкість зміни
температури (°C/s), \texttt{u\_x} --- нахил температурної кривої
(°C/см), \texttt{u\_\{xx\}} --- опуклість температурної кривої (°C/см²),

\begin{verbatim}
c = щільність тепла (калор/г \cdot °C)
\end{verbatim}

\begin{verbatim}
k — теплопровідність (кал/см \cdot при \cdot °C), \rho — щільність (г/<sup>см3</sup>), \alpha^2 — це теплова дифузія (см<sup>²</sup>/с).
\end{verbatim}

\begin{verbatim}
4. Слід зазначити, що теплова дифузія матеріалу \alpha^2 = \frac{k}{\epsilon \rho} прямо пропорційна теплопровідності матеріалу і обернено пропорційна густині \rho та питому теплоємності \epsilon ; Ми сподіваємося, що цей зв'язок відповідає інтуїтивним ідеям читача.
\end{verbatim}

\section{ЗАВДАННЯ}\label{ux437ux430ux432ux434ux430ux43dux43dux44f-2}

\begin{verbatim}
1. Підставити розмірність кожної з u,\ u_t,\ \dots величин у рівняння
\end{verbatim}

\begin{verbatim}
u_t = \alpha^2 u_{xx} - \beta u
\end{verbatim}

\begin{verbatim}
і переконайтеся, що всі терміни мають однаковий розмір {}^{\circ}\mathrm{C}/c .
\end{verbatim}

\begin{enumerate}
\def\labelenumi{\arabic{enumi}.}
\setcounter{enumi}{1}
\tightlist
\item
  Підставити розмірність кожної з величин у рівняння
\end{enumerate}

\begin{verbatim}
u_t = \alpha^2 u_{xx} - v u_x,
\end{verbatim}

\begin{verbatim}
де \upsilon вимірюється в одиницях швидкості, і переконайтеся, що всі терми мають однакову розмірність.
\end{verbatim}

3 Отримаємо рівняння теплопровідності

\begin{verbatim}
u_{t} = \frac{1}{c\rho} \frac{\partial}{\partial x} [h(x) u_{x}] + f(x, t)
\end{verbatim}

\begin{verbatim}
Для випадку, коли теплопровідність k\left( x\right) залежить від координати X.
\end{verbatim}

\begin{enumerate}
\def\labelenumi{\arabic{enumi}.}
\setcounter{enumi}{3}
\tightlist
\item
  Припустимо, що \texttt{u(x,\ t)} --- це значення концентрації речовини
  в потоку, що рухається зі швидкістю v. Припустимо, що концентрація
  змінюється внаслідок процесів дифузії та конвекції. Виведемо рівняння
\end{enumerate}

\begin{verbatim}
u_1 = \alpha^2 u_{xx} - v u_x,
\end{verbatim}

\begin{verbatim}
Виходячи з того, що в будь-який момент часу в певному [x, x + \Delta x] сегменті матерії не з'являється і не зникає.
\end{verbatim}

ПРИМІТКА. Закон збереження матерії виглядає так:

\begin{verbatim}
Зміни маси речовини на сегменті [x, x + \Delta x] = = Зміна внаслідок дифузії через межі + Зміна внаслідок переносу через межі.
\end{verbatim}

Лекція 5

\section{ЗМІННЕ
РОЗДІЛЕННЯ}\label{ux437ux43cux456ux43dux43dux435-ux440ux43eux437ux434ux456ux43bux435ux43dux43dux44f}

МЕТА ЛЕКЦІЇ: Познайомити читача з потужним методом розділення змінних і
показати, як цей метод можна використати для розв'язання добре відомої
проблеми дифузії. Оскільки цей метод створює труднощі для деяких учнів
(головним чином через громіздкі алгебранічні перетворення), ми
супроводжуємо презентацію інтуїтивними міркуваннями.

Основна ідея методу полягає в тому, щоб розкласти початкову умову на
найпростіші компоненти, знайти відгук системи на кожен найпростіший
компонент і підсумувати всі відповіді. Таким чином можна знайти
відповідь на довільну початкову умову.

Існуючий метод поділу змінних певною мірою зменшує можливість такої
інтерпретації методу, але все ж не виключає його.

Метод розділення змінних є одним із найвідоміших методів розв'язання
зсувних задач і використовується, якщо:

\begin{itemize}
\tightlist
\item
  \begin{enumerate}
  \def\labelenumi{\arabic{enumi}.}
  \tightlist
  \item
    Рівняння є лінійним і однорідним (не обов'язково з постійними
    коефіцієнтами).
  \end{enumerate}
\item
  \begin{enumerate}
  \def\labelenumi{\arabic{enumi}.}
  \setcounter{enumi}{1}
  \tightlist
  \item
    Граничні умови визначаються як
  \end{enumerate}
\end{itemize}

\begin{verbatim}
\alpha u_x(0, t) + \beta u(0, t) = 0,
\end{verbatim}

\begin{verbatim}
\gamma u_x(1, t) + \delta u(1, t) = 0, де \alpha , \beta , \gamma і \delta — константи (граничні умови, задані в такій формі, називаються лінійними однорідними граничними умовами).
\end{verbatim}

Метод був створений за часів Фур'є (зазвичай відомий як метод Фур'є) і,
здається, є найпопулярнішим (де це застосовно).

\pandocbounded{\includegraphics[keepaspectratio]{images/_page_34_Figure_8.jpeg}}

Рис. 5.1. Схематичне представлення задачі дифузії.

Замість того, щоб вивчати метод у загальному випадку, давайте спочатку
розглянемо конкретну задачу (загальний випадок обговоримо пізніше).
Розглянемо змішану задачу типу дифузії: щоб знайти розв'язок

(PPP)

\begin{verbatim}
u_t = \alpha^2 u_{xx}
, 0 < x < 1 , 0 < t < \infty ,
\end{verbatim}

що задовольняє граничні умови

(GI)

\begin{verbatim}
\begin{cases} u(0, t) = 0, & 0 < t < \infty, \\ u(1, t) = 0, & 0 < t < \infty, \end{cases}
\end{verbatim}

та початкова умова

(ГУ)

\begin{verbatim}
u(x, 0) = \varphi(x), 0 \le x \le 1.
\end{verbatim}

\begin{verbatim}
Перш ніж почати розділяти змінні, давайте дамо фізичну інтерпретацію нашої проблеми. Отже, існує стрижень скінченної довжини, кінці якого підтримуються при сталій температурі, рівній нулю (насправді кінці можна підтримувати при значно вищій температурі, значення якої вважається початковою точкою). Додаткові дані про завдання подаються як початкова умова. Наша мета — знайти розподіл температури u\left(x,t\right) в наступний час.
\end{verbatim}

\section{Загальні принципи методу розділення
змінних}\label{ux437ux430ux433ux430ux43bux44cux43dux456-ux43fux440ux438ux43dux446ux438ux43fux438-ux43cux435ux442ux43eux434ux443-ux440ux43eux437ux434ux456ux43bux435ux43dux43dux44f-ux437ux43cux456ux43dux43dux438ux445}

Для найпростішого диференціального рівняння в похідних похідних
розділення змінних є пошуком розв'язків у вигляді

\begin{verbatim}
u(x, t) = X(x)T(t),
\end{verbatim}

де \texttt{X(x)} --- це функція, що залежить лише від змінної x, а
\texttt{T(t)} --- функція, що залежить лише від t. Це розв'язок у
певному сенсі є найпростішим, оскільки температура \texttt{u(x,\ t)},
представлена у такій формі, зберігатиме «форму» профілю у різний час
часу t (див. рисунок 5.2).

\pandocbounded{\includegraphics[keepaspectratio]{images/_page_35_Figure_6.jpeg}}

РИСУНОК 5.2 Графік функції X(x)

\begin{verbatim}
T(t) у різних точках часу t.

~~
Загальна ідея полягає в тому, щоб знайти нескінченну кількість таких часткових диференціальних розв'язків (які задовольняють граничні умови). Ці найпростіші функції u_n\left(x,\ t\right) = X_n\left(x\right)T_n\left(t\right) (які називаються фундаментальними розв'язками) є, так би мовити, елементарними будівельними блоками, з яких будується розв'язок нашої проблеми. Розв'язок нашої задачі u\left(x,\ t\right) у вигляді такої лінійної комбінації фундаментальних розв'язків X_n\left(x\right)T_n\left(t\right) , що отримана сума
~~
\end{verbatim}

\texttt{\textbackslash{}sum\_\{n=1\}\^{}\{\textbackslash{}infty\}\ A\_n\ X\_n(x)\ T\_n(t)}

\begin{verbatim}
задовольняє початкові умови. Оскільки ця сума задовольняє рівняння та граничні умови, вона є розв'язком початкової задачі. Ми маємо робити всі ці розрахунки достатньо детально.

# Розділення змінних

КРОК 1. (Пошук елементарних розв'язків рівняння з частковими похідними.)

Ми хочемо знайти функцію u(x, t), яка є розв'язком

Наступного завдання:
\end{verbatim}

\texttt{\textbackslash{}begin\{aligned\}}
\texttt{(\textbackslash{}text\{УЧП\})\textbackslash{}quad\ \&\ u\_t\ =\ \textbackslash{}alpha\^{}2\ u\_\{xx\},\ \&\&\ 0\ \textless{}\ x\ \textless{}\ 1,\textbackslash{}\ 0\ \textless{}\ t\ \textless{}\ \textbackslash{}infty,\ \textbackslash{}\textbackslash{}}
\texttt{(\textbackslash{}text\{ГУ\})\textbackslash{}quad\ \&\ u(0,t)\ =\ 0,\textbackslash{}quad\ u(1,t)\ =\ 0,\ \&\&\ 0\ \textless{}\ t\ \textless{}\ \textbackslash{}infty,\ \textbackslash{}\textbackslash{}}
\texttt{(\textbackslash{}text\{НУ\})\textbackslash{}quad\ \&\ u(x,0)\ =\ \textbackslash{}varphi(x),\ \&\&\ 0\ \textless{}\ x\ \textbackslash{}leqslant\ 1.}
\texttt{\textbackslash{}end\{aligned\}}

\begin{verbatim}
Пошукаємо розв'язки, які представлені у вигляді u(x, t) = X(x) T(t). Для цього підставляємо вираз X(x) T(t) у рівняння. Внаслідок підстановки отримуємо
\end{verbatim}

X(x) T'(t) = q\^{}2 X'\,'(x) T(t).

\begin{verbatim}
~~
**Тепер** виконаємо операцію, властиву цьому методу: *поділимо* обидві сторони останнього рівняння на \alpha^2 X(x) T(t) , у результаті чого отримуємо
~~
\end{verbatim}

\texttt{\textbackslash{}frac\{T\textquotesingle{}(t)\}\{\textbackslash{}alpha\^{}2\ T(t)\}\ =\ \textbackslash{}frac\{X\textquotesingle{}\textquotesingle{}(x)\}\{X(x)\}.}

\begin{verbatim}
Цей вираз вважається таким, що розділив змінні, тобто ліва частина рівняння залежить лише від t, а права — лише від x. Оскільки x і t не є розділенням комунікації один від одного, кожна частина рівняння має бути сталою. Позначимо цю константу k, тоді
\end{verbatim}

\texttt{\textbackslash{}frac\{T\textquotesingle{}\}\{\textbackslash{}alpha\^{}2\ T\}\ =\ \textbackslash{}frac\{X\textquotesingle{}\textquotesingle{}\}\{X\}\ =\ k}

\begin{verbatim}
або
\end{verbatim}

\texttt{T\textquotesingle{}-k\textbackslash{}alpha\^{}2T=0,} X'\,'-kX=0.

\begin{verbatim}
Тепер можливо розв'язати кожне з цих звичайних диференціальних рівнянь. Добуток відповідних розв'язків задовольнятиме початкове рівняння з похідними похідними. (Зверніть увагу, що ми спростили оригінальне диференціальне рівняння другого порядку до двох звичайних диференціальних рівнянь.)

Тепер звернемо увагу на таку важливу обставину: константа відстань k має бути від'ємною ^1 ) (іншими словами, функції T(t) мають прагнути до нуля в точці t \to \infty ). З огляду на це, давайте розглянемо позначення k = -\lambda^2 , де \lambda не дорівнює нулю (у цьому випадку вираз -\lambda^2 завжди буде від'ємним). Враховуючи нову нотацію для константи відстань двох звичайних диференціальних рівнянь,

<sup>1)</sup> Інакше рівняння X'' - kX = 0 з граничними умовами X(0) = X(1) = 0 має лише гривну розв'язку X(x) = 0. Ед.

записуються у формі
\end{verbatim}

\texttt{T\textquotesingle{}\ +\ \textbackslash{}lambda\^{}2\ \textbackslash{}alpha\^{}2\ T\ =\ 0,}
\texttt{X\textquotesingle{}\textquotesingle{}\ +\ \textbackslash{}lambda\^{}2\ X\ =\ 0.}

\begin{verbatim}
Розв'яжемо ці рівняння. Це рівняння стандартного типу. Їхні загальні рішення фіксуються у формі
\end{verbatim}

\texttt{T(t)\ =\ Ae\^{}\{-\textbackslash{}gamma\ *\_\{O\}\ *\_\{I\}\}}

\begin{verbatim}
(A — довільна константа),
\end{verbatim}

\texttt{X(x)\ =\ A\ \textbackslash{}sin(\textbackslash{}lambda\ x)\ +\ B\ \textbackslash{}cos(\textbackslash{}lambda\ x)}

\begin{verbatim}
(A, B — довільні константи).

Відповідно, функції виду
\end{verbatim}

\texttt{u(x,\ t)\ =\ e\^{}\{-\textbackslash{}lambda\^{}2\ \textbackslash{}alpha\^{}2\ t\}\ {[}A\ \textbackslash{}sin(\textbackslash{}lambda\ x)\ +\ B\ \textbackslash{}cos(\textbackslash{}lambda\ x){]}}

\begin{verbatim}
~~
(де A, B і \lambda — довільні константи) задовольняють NPP u_t = \alpha^2 u_{xx} (пропонується перевірити це у задачі 1, розміщеній наприкінці цієї лекції). Отже, маємо нескінченну множину функцій, які задовольняють початкове диференціальне рівняння в похідних похідних.
~~

КРОК 2. (Пошук розв'язків, що задовольняють граничним умовам.)

Зараз ситуація така: ми маємо нескінченну кількість розв'язків початкового рівняння, але не всі вони задовольняють граничні або початкові умови. Наступний крок — вибрати таку підмножину розв'язків вигляду

(5.1)
\end{verbatim}

\texttt{e\^{}\{-\textbackslash{}lambda\^{}2\ \textbackslash{}alpha\^{}\{2\textbackslash{}ell\}\}\ {[}A\ \textbackslash{}sin(\textbackslash{}lambda\ x)\ +\ B\ \textbackslash{}cos(\textbackslash{}lambda\ x){]},}

\begin{verbatim}
які задовольняють граничні умови
\end{verbatim}

u(0, t) = 0,

\begin{verbatim}
u(1, t) = 0. Для цього підставляємо розв'язки (5.1) у ці граничні умови. У результаті ми отримуємо
\end{verbatim}

\texttt{u(0,\ t)\ =\ Be\^{}\{-\textbackslash{}lambda\^{}2\ \textbackslash{}alpha\^{}\{2t\}\}\ =\ 0\ \textbackslash{}Rightarrow\ B\ =\ 0,}
\texttt{u(1,\ t)\ =\ Ae\^{}\{-\textbackslash{}lambda\^{}2\ \textbackslash{}alpha\^{}\{2t\}\}\ \textbackslash{}sin\ \textbackslash{}lambda\ =\ 0\ \textbackslash{}Rightarrow\ \textbackslash{}sin\ \textbackslash{}lambda\ =\ 0.}

\begin{verbatim}
Друга гранична умова накладає обмеження на можливі значення константи \lambda розділення: вона має бути коренем рівняння \sin\lambda=0 . Іншими словами, щоб виконати умову u(1, t)=0, необхідно вимагати виконання
\end{verbatim}

\texttt{\textbackslash{}lambda\ =\ \textbackslash{}pm\ \textbackslash{}pi}

\begin{verbatim}
, \pm 2\pi , ... або
\end{verbatim}

\texttt{\textbackslash{}lambda\_n\ =\ \textbackslash{}pm\ n\textbackslash{}pi}

\begin{verbatim}
, n = 1, 2, \ldots Зверніть увагу, що друга гранична умова може бути виконана, якщо поставити A = 0, але в цьому випадку розв'язок (5.1) буде ідентичним нулю.

![](images/_page_38_Figure_3.jpeg)

Рис. 5.3. Фундаментальні рішення n_n(x, t) = A_n e^{-(\pi n \alpha)^{n/t}} \sin{(\pi n x)} .

Тепер, коли ми зробили другий крок, маємо нескінченний набір ознак

(5.2)
\end{verbatim}

\texttt{u\_n(x,\ t)\ =\ A\_n\ e\^{}\{-(n\textbackslash{}pi\textbackslash{}alpha)\^{}2\ t\}\ \textbackslash{}sin(n\textbackslash{}pi\ x),\ \textbackslash{}quad\ n\ =\ 1,\ 2\ \textbackslash{}dots}

\begin{verbatim}
кожна з яких задовольняє рівняння з частковими похідними та граничними умовами). Отримані функції — це будівельні блоки, з яких ми будемо створювати рішення нашої проблеми. Це рішення буде певною сумою

<sup>2)</sup> Зверніть увагу, що функції u_n і u_{-n} відрізняються лише знаком.

цих найпростіших функцій. Конкретний тип суми залежить від початкового стану. Фіг. Рисунок 5.3 показує графіки кількох таких фундаментальних функцій

КРОК 3. (Пошук розв'язку, який задовольняє рівняння, межу та початкові умови.)

Останній (і, ймовірно, найцікавіший з математичної точки зору) крок — знайти таку суму фундаментальних розв'язків
\end{verbatim}

\texttt{u(x,\ t)\ =\ \textbackslash{}sum\_\{n=1\}\^{}\{\textbackslash{}infty\}\ A\_n\ e\^{}\{-(n\textbackslash{}pi\textbackslash{}alpha)\^{}n\ t\}\ \textbackslash{}sin(n\textbackslash{}pi\ x),}

\begin{verbatim}
тобто при виборі таких коефіцієнтів A_n , що функція задовольнятиме початкову умову
\end{verbatim}

\texttt{u\textbackslash{}left(\ x,\textbackslash{}\ 0\textbackslash{}right)\ =\textbackslash{}varphi\textbackslash{}left(\ x\textbackslash{}right)\ .}

\begin{verbatim}
Підставляючи суму в початковому стані, отримуємо

(5.3)
\end{verbatim}

\texttt{q(x)\ =\ \textbackslash{}sum\_\{n=1\}\^{}\{\textbackslash{}infty\}\ A\_n\ \textbackslash{}sin(n\textbackslash{}pi\ x).}

\begin{verbatim}
Це рівняння підводить нас до цікавого питання: чи можна розкласти початкову температуру \phi(x) на ряд елементарних функцій вигляду
\end{verbatim}

\texttt{A\_1\ \textbackslash{}sin(\textbackslash{}pi\ x)\ +\ A\_2\ \textbackslash{}sin(2\textbackslash{}pi\ x)\ +\ A\_3\ \textbackslash{}sin(3\textbackslash{}pi\ x)\ +\ \textbackslash{}dots}

\begin{verbatim}
?

Позитивну відповідь на це питання дав французький математик Жозеф Фур'є. Виявилося, що для достатньо хороших функцій така декомпозиція можлива ^1 ). Тоді виникає нове питання: знайти коефіцієнти розкладу A_n ?

Насправді це дуже просто. Для цього використаємо властивість системи функцій
\end{verbatim}

\texttt{\textbackslash{}\{\textbackslash{}sin(n\textbackslash{}pi\ x);\ n=1,\ 2,\ \textbackslash{}ldots\textbackslash{}\},}

\begin{verbatim}
відомий як ортогональність.

У нашому випадку (див. завдання 2) ці функції задовольняють умови
\end{verbatim}

\texttt{\textbackslash{}int\_\{0\}\^{}\{1\}\ \textbackslash{}sin(m\textbackslash{}pi\ x)\ \textbackslash{}sin(n\textbackslash{}pi\ x)\ dx\ =\ \textbackslash{}begin\{cases\}\ 0,\ \&\ m\ \textbackslash{}neq\ n,\ \textbackslash{}\textbackslash{}\ 1/2,\ \&\ m\ =\ n.\ \textbackslash{}end\{cases\}}

\begin{verbatim}
Умови розкладу функцій у рядах Фур'є наведені в лекції 11. Для глибшого ознайомлення з цим питанням можна звернутися до численних освітніх літератур (див., наприклад: Г. П. Толстов, Серія Фур'є. — Москва: Наука, 1980). — Ед. Ед.

Ця властивість чітко показана на рис. 5.4, яка показує графіки деяких функцій із цієї системи.

Отже, ми хочемо знайти коефіцієнти у розкладі
\end{verbatim}

\texttt{\textbackslash{}varphi(x)\ =\ A\_1\ \textbackslash{}sin(\textbackslash{}pi\ x)\ +\ A\_2\ \textbackslash{}sin(2\textbackslash{}pi\ x)\ +\ A\_3\ \textbackslash{}sin(3\textbackslash{}pi\ x)\ +\ \textbackslash{}dots}

\begin{verbatim}
![](images/_page_40_Figure_5.jpeg)

РИСУНОК 5.4. Послідовність ортогональних функцій

Помножте обидві сторони цього співвідношення на \sin{(m\pi x)} (m — довільне ціле число) і інтегруйте від нуля до одиниці. У результаті ми отримуємо
\end{verbatim}

\texttt{\textbackslash{}int\_\{0\}\^{}\{1\}\ \textbackslash{}varphi(x)\ \textbackslash{}sin(m\textbackslash{}pi\ x)\ \textbackslash{},\ dx\ =\ A\_\{m\}\ \textbackslash{}int\_\{0\}\^{}\{1\}\ \textbackslash{}sin\^{}\{2\}(m\textbackslash{}pi\ x)\ \textbackslash{},\ dx\ =\ \textbackslash{}frac\{A\_\{m\}\}\{2\}}

\begin{verbatim}
(усі інші члени зрівнялися з нулем через ортогональність). Розв'язуючи рівняння відносно A_m , отримуємо
\end{verbatim}

\texttt{A\_m\ =\ 2\ \textbackslash{}int\_0\^{}1\ \textbackslash{}varphi(x)\ \textbackslash{}sin(m\textbackslash{}pi\ x)\ dx.}

\begin{verbatim}
Отже, отримано, що розв'язок записується у вигляді

(5.4)
\end{verbatim}

\texttt{u(x,\ t)\ =\ \textbackslash{}sum\_\{n=1\}\^{}\{\textbackslash{}infty\}\ A\_n\ e\^{}\{-(n\textbackslash{}pi\textbackslash{}alpha)\^{}2\ t\}\ \textbackslash{}sin(\textbackslash{}pi\ n\ x),}

\begin{verbatim}
де коефіцієнти A_n визначаються формулами
\end{verbatim}

\texttt{A\_n\ =\ 2\textbackslash{}int\_0\^{}1\ \textbackslash{}varphi(x)\ \textbackslash{}sin(\textbackslash{}pi\ nx)\ dx}

\begin{verbatim}
Ви можете переконатися, що отримане рішення задовольняє всі умови початкової задачі. На цьому завершується третій крок.

Багато студентів розчаровані, коли бачать, наскільки складним виглядає рішення, і майже не намагаються дивитися на нього знову (виглядає дуже погано). Однак це рішення не здаватиметься таким складним, якщо ви витратите час на його аналіз. Насправді, складніша форма розв'язку робить його більш інформативним. Нижче наведено кілька нотаток, які допоможуть вам інтерпретувати це рішення.

# ЗАУВАЖЕННЯ

Зверніть увагу, що єдина різниця між синусоїдальним розкладом функції φ(x) і розв'язком (5.4) — це наявність часового фактора e<sup>-(nπα)<sup>3</sup>1</sup> у кожному члені ряду. Отже, якщо початкова умова має просту декомпозицію, наприклад, вигляду
\end{verbatim}

\texttt{\textbackslash{}varphi(x)\ =\ \textbackslash{}sin(\textbackslash{}pi\ x)\ +\ \textbackslash{}frac\{1\}\{2\}\textbackslash{}sin(3\textbackslash{}pi\ x),}

\begin{verbatim}
тоді рішення можна негайно записати
\end{verbatim}

\texttt{u(x,\ t)\ =\ e\^{}\{-(\textbackslash{}pi\textbackslash{}alpha)\^{}2\ t\}\ \textbackslash{}sin(\textbackslash{}pi\ x)\ +\ \textbackslash{}frac\{1\}\{2\}\ e\^{}\{-(3\textbackslash{}pi\textbackslash{}alpha)\^{}2\ t\}\ \textbackslash{}sin(3\textbackslash{}pi\ x).}

\begin{verbatim}
У цьому випадку очевидно, що якщо ми розкладемо \phi(x) на ряд Фур'є вздовж синусів, то отримуємо
\end{verbatim}

A\_1 = 1,

\begin{verbatim}
A_2 = 0, A_3 = 1/2, A_4 = A_5 = \dots = 0. - 2. Розв'язок (5.4) можна інтерпретувати так: ми представляємо початкову температуру \varphi(x) як суму найпростіших функцій виду A_n \sin(\pi nx) ; кожна така функція породжує відповідь A_n e^{-(\pi n\alpha)^2 t} \sin(\pi nx) . Сумуючи всі такі відповіді, отримуємо розв'язок, який відповідає початковій умові u(x, 0) = \varphi(x) .
- 3. Кожен член у декомпозиції
\end{verbatim}

\texttt{u(x,\ t)\ =\ A\_1\ e\^{}\{-(\textbackslash{}pi\textbackslash{}alpha)\^{}2\ t\}\ \textbackslash{}sin(\textbackslash{}pi\ x)\ +\ A\_2\ e\^{}\{-(2\textbackslash{}pi\textbackslash{}alpha)\^{}2\ t\}\ \textbackslash{}sin(2\textbackslash{}pi\ x)\ +\ \textbackslash{}cdots}

\begin{verbatim}
є функцією від x і t. Зверніть увагу, що внесок членів з великими числами в t>0 дуже малий через множник e^{-(n\pi\alpha)^2 t} . Відповідно, після закінчення досить великого Р и с

![](images/_page_42_Figure_2.jpeg)

. 5.5. У задачах дифузії вищі порядки терміни швидко затухають; \phi(x) — початкова температура. Слід зазначити, що по-перше, у розчині зникають високі частоти — по-перше, зникли найбільші коливання міді.

Повний розв'язок приблизно збігається з першим членом, який є півхвилі синусоїди, що з часом згасає (рис. 1). 5.5.)

# ЗАВДАННЯ

1. Показати, що функції
\end{verbatim}

\texttt{u(x,\ t)\ =\ e\^{}\{-\textbackslash{}lambda\^{}2\ \textbackslash{}alpha\textquotesingle{}\ t\}\ \textbackslash{}left{[}\ A\ \textbackslash{}sin(\textbackslash{}lambda\ x)\ +\ B\ \textbackslash{}cos(\textbackslash{}lambda\ x)\ \textbackslash{}right{]}}

\begin{verbatim}
задовольняють рівняння u_l = \alpha^2 u_{xx} для довільних значень A, B і \lambda .

2. Покажи це
\end{verbatim}

\texttt{\textbackslash{}int\_0\^{}1\ \textbackslash{}sin(m\textbackslash{}pi\ x)\ \textbackslash{}sin(n\textbackslash{}pi\ x)\ dx\ =\ \textbackslash{}begin\{cases\}\ 0,\ \&\ m\ \textbackslash{}neq\ n,\ \textbackslash{}\textbackslash{}\ 1/2,\ \&\ m\ =\ n.\ \textbackslash{}end\{cases\}}

\begin{verbatim}
ПРИМІТКА. Використання ідентичності
\end{verbatim}

\texttt{\textbackslash{}sin(mx)\textbackslash{}sin(nx)\ =\ \textbackslash{}frac\{1\}\{2\}{[}\textbackslash{}cos(m-n)x\ -\ \textbackslash{}cos(m+n)x{]}.}

\begin{verbatim}
3. Знайти розклад у рядах Фур'є синусами функції \varphi(x) = 1 на сегменті [0, 1]. Побудуйте перші три терміни розкладу.

Використовуючи результати завдання 3, знайдіть розв'язок наступної змішаної задачі:

(PPP)
\end{verbatim}

u\_t = u\_\{xx\}

\begin{verbatim}
0 < x < 1 , 
(GI)
\end{verbatim}

\texttt{\textbackslash{}begin\{cases\}\ u(0,\ t)\ =\ 0,\ \textbackslash{}\textbackslash{}\ u(1,\ t)\ =\ 0,\ \textbackslash{}end\{cases\}}

\begin{verbatim}
0 < t < \infty , 
(НУ) u(x, 0) = 1, \quad 0 \le x \le 1 .

(Зверніть увагу, що ця проблема фізично не має значення, оскільки припускається, що температура на кінцях стрижня миттєво знижується з одного до нуля \phi(x) .

5. Знайти розв'язок проблеми 4, якщо початкова умова встановлена у вигляді
\end{verbatim}

\texttt{u(x,\ 0)\ =\ \textbackslash{}sin(2\textbackslash{}pi\ x)\ +\ \textbackslash{}frac\{1\}\{3\}\textbackslash{}sin(4\textbackslash{}pi\ x)\ +\ \textbackslash{}frac\{1\}\{5\}\textbackslash{}sin(6\textbackslash{}pi\ x).}

\begin{verbatim}
6. Розв'яжіть задачу 4, якщо початкова умова u\left( x,\ 0\right) =x-x^{2},\quad 0< x<1. Лекція 6

# ПЕРЕТВОРЕННЯ НЕОДНОРІДНИХ ГРАНИЧНИХ УМОВ У ОДНОРІДНІ ГРАНИЧНІ УМОВИ

МЕТА ЛЕКЦІЇ: Показати, як виникає змішана задача
\end{verbatim}

\texttt{\textbackslash{}begin\{array\}\{lll\}\ (\textbackslash{}text\{УЧП\})\ \&\ u\_t\ -\ a\^{}2\ u\_\{xx\}\ =\ f\textbackslash{}left(x,\ \textbackslash{}\ t\textbackslash{}right)\ \&\ \textbackslash{}frac\{\textbackslash{}beta\_t\}\{97\}\ \textbackslash{}\textbackslash{}\ (\textbackslash{}text\{ГУ\})\ \&\ \textbackslash{}begin\{cases\}\ \textbackslash{}alpha\_1\ u\_x\textbackslash{}left(0,\ \textbackslash{}\ t\textbackslash{}right)\ +\ \textbackslash{}beta\_1\ u\textbackslash{}left(0,\ \textbackslash{}\ t\textbackslash{}right)\ =\ g\_1\textbackslash{}left(t\textbackslash{}right),\ \textbackslash{}\textbackslash{}\ a\_2\ u\_x\textbackslash{}left(L,\ \textbackslash{}\ t\textbackslash{}right)\ +\ \textbackslash{}beta\_2\ u\textbackslash{}left(L,\ \textbackslash{}\ t\textbackslash{}right)\ =\ g\_2\textbackslash{}left(t\textbackslash{}right),\ \textbackslash{}end\{cases\}\ \textbackslash{}\textbackslash{}\ (\textbackslash{}text\{HУ\})\ \&\ u\textbackslash{}left(x,\ \textbackslash{}\ 0\textbackslash{}right)\ =\ \textbackslash{}varphi\textbackslash{}left(x\textbackslash{}right)\ \textbackslash{}end\{array\}}

\begin{verbatim}
можна звести до задачі вигляду
\end{verbatim}

\texttt{U\_\{t\}\ -\ \textbackslash{}alpha\^{}\{2\}U\_\{xx\}\ =\ F(x,\ t),}
\texttt{\textbackslash{}alpha\_\{1\}U\_\{x\}(0,\ t)\ +\ \textbackslash{}beta\_\{1\}U(0,\ t)\ =\ 0,}
\texttt{\textbackslash{}alpha\_\{2\}U\_\{x\}(L,\ t)\ +\ \textbackslash{}beta\_\{2\}U(L,\ t)\ =\ 0,}
\texttt{U(x,\ 0)\ =\ \textbackslash{}varphi(x)}

\begin{verbatim}
з нульовими граничними умовами. Цей новий виклик можна подолати

- 1. Шляхом ділення змінних (якщо нова задача виявиться однорідною, тобто F(x, t) = 0).

~~
- 2. Методом інтегральних перетворень і розкладу власними функціями, якщо F(x, t) \neq 0 .
~~

Метод розділення змінних, про який йшлося в попередній лекції, є досить потужним, а розв'язки, отримані за його допомогою, подаються у зручній формі. Однак читач повинен розуміти, що цей метод не застосовний до всіх задач. Щоб метод розділення змінних був застосовним, граничні умови мають бути лінійними, однорідними, тобто представленими як

(6.1)
\end{verbatim}

\texttt{\textbackslash{}begin\{aligned\}\ \textbackslash{}alpha\_1\ u\_x(0,\ t)\ +\ \textbackslash{}beta\_1\ u(0,\ t)\ \&=\ 0,\ \textbackslash{}\textbackslash{}\ \textbackslash{}alpha\_2\ u\_x(L,\ t)\ +\ \textbackslash{}beta\_2\ u(L,\ t)\ \&=\ 0.\ \textbackslash{}end\{aligned\}}

\begin{verbatim}
Мета цієї лекції — показати, як виникає задача з неоднорідними граничними умовами форми

(6.2)
\end{verbatim}

\texttt{\textbackslash{}begin\{aligned\}\ \&\ (\textbackslash{}mathsf\{У\}\textbackslash{}mathsf\{Ч\}\textbackslash{}Pi)\ \&\ u\_t\ =\ \textbackslash{}alpha\^{}2\ u\_\{xx\},\ \textbackslash{}\textbackslash{}\ \&\ (\textbackslash{}mathsf\{\textbackslash{}Gamma\}\textbackslash{}mathsf\{Y\})\ \&\ \textbackslash{}begin\{cases\}\ \textbackslash{}alpha\_1\ u\_x\ \textbackslash{},\ (0,\ \textbackslash{},\ t)\ +\ \textbackslash{}beta\_1\ u\ \textbackslash{},\ (0,\ \textbackslash{},\ t)\ =\ g\_1\ \textbackslash{},\ (t)\ \textbackslash{}\textbackslash{}\ \textbackslash{}alpha\_2\ u\_x\ \textbackslash{},\ (L,\ \textbackslash{},\ t)\ +\ \textbackslash{}beta\_2\ u\ \textbackslash{},\ (L,\ \textbackslash{},\ t)\ =\ g\_2\ \textbackslash{},\ (t)\ \textbackslash{}end\{cases\}\ \textbackslash{}text\{\ (неоднородные\ \}\ \textbackslash{}mathsf\{\textbackslash{}Gamma\}\textbackslash{}mathsf\{Y\}),\ \textbackslash{}\textbackslash{}\ \&\ (\textbackslash{}mathsf\{H\}\textbackslash{}mathsf\{Y\})\ \&\ u\ \textbackslash{},\ (x,\ \textbackslash{},\ 0)\ =\ \textbackslash{}varphi\ \textbackslash{},\ (x)\ \textbackslash{}end\{aligned\}}

\begin{verbatim}
можна розв'язати, зведивши її до задачі з однорідними граничними умовами.

Нове завдання можна розв'язати іншими методами

(наприклад, шляхом розкладу за власними функціями).

Почнемо з надзвичайно простого прикладу перетворення гетерогенних граничних умов на однорідні.

# Перетворення неоднорідних граничних умов у однорідні граничні умови

Розглянемо проблему поширення тепла в термоізольованому стрижні, кінці якого підтримуються при сталих температурах k_1 і k_2, тобто
\end{verbatim}

\texttt{\textbackslash{}begin\{aligned\}}
\texttt{(\textbackslash{}text\{УЧП\})\textbackslash{}quad\ \&\ u\_t\ =\ \textbackslash{}alpha\^{}2\ u\_\{xx\},\ \&\&\ 0\ \textless{}\ x\ \textless{}\ L,\textbackslash{}\ 0\ \textless{}\ t\ \textless{}\ \textbackslash{}infty,\ \textbackslash{}\textbackslash{}}
\texttt{(6.3)\textbackslash{}\ (\textbackslash{}text\{ГУ\})\textbackslash{}quad\ \&\ u(0,t)\ =\ k\_1,\textbackslash{}quad\ u(L,t)\ =\ k\_2,\ \&\&\ 0\ \textless{}\ t\ \textless{}\ \textbackslash{}infty,\ \textbackslash{}\textbackslash{}}
\texttt{(\textbackslash{}text\{НУ\})\textbackslash{}quad\ \&\ u(x,0)\ =\ \textbackslash{}varphi(x),\ \&\&\ 0\ \textbackslash{}leqslant\ x\ \textbackslash{}leqslant\ L.}
\texttt{\textbackslash{}end\{aligned\}}

\begin{verbatim}
Складність цієї задачі полягає в тому, що, оскільки граничні умови в ній є гетерогенними, ми не можемо розв'язати її шляхом розділення змінних. Однак очевидно, що в t \to \infty розв'язок нашої задачі схиляється до стаціонарного розв'язку, який лінійно (вздовж x) змінюється від температури k_1 до температури k_2 (рис. 6.1). Іншими словами, розумно припустити, що температуру в нашій задачі можна представити як суму
~~

# з двох термів
\end{verbatim}

\texttt{u\textbackslash{}left(x,\textbackslash{},t\textbackslash{}right)=}

\begin{verbatim}
~~
стаціонарна + тимчасова = \uparrow \uparrow Частина розв'язку, що залежить від HY і має тенденцію до нуля з помноженням на =\left[k_1+\frac{x}{L}\left(k_2-k_1\right)\right]+U\left(x,\,t\right). У цьому випадку наше завдання — знайти температуру переходу U\left( x,\,t\right) . Заміна
~~
\end{verbatim}

\texttt{u(x,\ t)\ =\ \textbackslash{}left{[}k\_1\ +\ \textbackslash{}frac\{x\}\{L\}(k\_2\ -\ k\_1)\textbackslash{}right{]}\ +\ U(x,\ t)}

\begin{verbatim}
~~
У початковій задачі (6.3) ми отримуємо нову задачу відносно невідомої функції U\left(x,\ t\right) . Ви можете вирішити проблему Р и с
~~

![](images/_page_45_Figure_7.jpeg)

~~
. 6.1. Розв'язання задачі (6.3) у різний час часу a- початкова температура u\left(x,\ 0\right)=\phi\left(x\right); b - стаціонарна температура u\left(x,\ \infty\right)=k_{1}+\frac{x}{L}\left(k_{2}-k_{1}\right); b - перехідна температура.
~~

~~
Відносно нова невідома функція U\left(x,\,t\right) і додає її до стаціонарного розв'язку, що дає бажану функцію u\left(x,\,t\right) . Виконавши ці прості перетворення з (6.3),
~~

Ми отримаємо

(6.4)
\end{verbatim}

\texttt{U\_\{t\}\ =\ \textbackslash{}alpha\^{}\{2\}U\_\{xx\}\ \textbackslash{}quad\ 0\ \textless{}\ x\ \textless{}\ L,}
\texttt{(\textbackslash{}Gamma\ Y)\ \textbackslash{}quad\ \textbackslash{}begin\{cases\}\ U(0,\ t)\ =\ 0,\ \textbackslash{}\textbackslash{}\ U(L,\ t)\ =\ 0,\ \textbackslash{}end\{cases\}\ \textbackslash{}quad\ 0\ \textless{}\ t\ \textless{}\ \textbackslash{}infty,}
\texttt{(HY)\ \textbackslash{}quad\ U(x,\ 0)\ =\ \textbackslash{}varphi(x)\ -\ \textbackslash{}left{[}k\_\{1\}\ +\ \textbackslash{}frac\{x\}\{L\}(k\_\{2\}\ -\ k\_\{1\})\textbackslash{}right{]}\ =\ \textbackslash{}overline\{\textbackslash{}varphi\}(x),}

\begin{verbatim}
~~
де \overline{\phi}(x) — нова, але відома початкова умова.
~~

Ця задача стосується не лише однорідного рівняння, а, на щастя, однорідних граничних умов, що дозволяє розв'язувати її методом розділення змінних. Насправді, читач, ймовірно, пам'ятає рішення:

(6.5)
\end{verbatim}

\texttt{U(x,\ t)\ =\ \textbackslash{}sum\_\{n=1\}\^{}\{\textbackslash{}infty\}\ a\_n\ e\^{}\{-(n\textbackslash{}pi\textbackslash{}alpha)\^{}4\ t\}\ \textbackslash{}sin(n\textbackslash{}pi\ x/L),}

\begin{verbatim}
рей
\end{verbatim}

\texttt{a\_n\ =\ \textbackslash{}frac\{2\}\{L\}\ \textbackslash{}int\_\{-\textbackslash{}infty\}\^{}\{L\}\ \textbackslash{}overline\{\textbackslash{}phi\}(\textbackslash{}xi)\ \textbackslash{}sin(\textbackslash{}pi\ n\textbackslash{}xi/L)\ d\textbackslash{}xi.}

\begin{verbatim}
Це все для сердечників із фіксованою температурою на кінцях. А як щодо більш реалістичних GU з правою стороною, залежною від часу? Основні ідеї тут ті ж, що й у попередній задачі, але вони набагато складніші.

# Перетворення граничних умов, залежних від часу, у нульові

Розглянемо типову задачу

(PPP)
\end{verbatim}

\texttt{u\_t\ =\ \textbackslash{}alpha\^{}2\ u\_\{xx\}}

\begin{verbatim}
~~
, 0 < x < L , 0 < t < \infty , 0 < t < \infty , 0 < t < \infty , 0 < t < \infty , 0 < t < \infty , 0 < t < \infty , 0 < t < \infty , 0 < t < \infty , 0 < t < \infty , 0 < t < \infty , 0 < t < \infty , 0 < t < \infty , 0 < t < \infty , 0 < t < \infty , 0 < t < \infty , 0 < t < \infty , 0 < t < \infty , 0 < t < \infty , 0 < t < \infty , 0 < t < \infty , 0 < t < \infty , 0 < t < \infty , 0 < t < \infty , 0 < t < \infty , 0 < t < \infty , 0 < t < \infty , , 0 < t < \infty , 0 < t < \infty , 0 < t < \infty , 0 < t < \infty , 0 < t < \infty , 0 < t < \infty , 0 < t < \infty , 0 < t < \infty , 0 < t < \infty , 0 < t < \infty , 0 < t < \infty , 0 < t < \infty , 0 < t < \infty , 0 < t < \infty , 0 < t < \infty , 0 < t < \infty , 0 < t < \infty , 0 < t < \infty , 0 < t < \infty , 0 < t < \infty , 0 < t < \infty , 0 < t < \infty , 0 < t < \infty , 0 < t < \infty , 0 < t < \infty , 0 < t < \infty , 0 < t < \infty , 0 < t < \infty , 0 < t < \infty , 0 < t < \infty , 0 < t < \infty , 0 < t < \infty , 0 < t < \infty , 0 < t < \infty , 0 < t < \infty , 0 < t < \infty , 0 < t < \infty , 0 < t < \infty , 0 < t < \infty , 0 < t < \infty , 0 < t < \infty , 0 < t < \infty , 0 < t < \infty , 0 < t < \infty , 0 < t < \infty , 0 < t < \infty , 0 < t < \infty , 0 < t < \infty , 0 < t < \infty , 0 < t < \infty , 0 < t < \infty , 0 < t < \infty , 0 < t < \infty , 0 < t < \infty , 0 < t < \infty , 0 < t < \infty , 0 < t < \infty , 0 < t < \infty , 0 < t < \infty , 0 < t < \infty , 0 < t < \infty , 0 < t < \infty , 0 < t < \infty , 0 < t < \infty , 0 < t < \infty , 0 < t < \infty , 0 < t < \infty , 0 < t < \infty , 0 < t < \infty , 0 < t < \infty , 0 < t < \infty , 0 < t < \infty , 0 < t < \infty , 0 < t < \infty , 0 < t < \infty , 0 < t < \infty , 0 < t < \infty , 0 < t < \infty , 0 < t < \infty , 0 < t < \infty , 0 < t < \infty , 0 < t < \infty , 0 < t < \infty , 0 < t < \infty , 0 < t < \infty , 0 < t < \infty , 0 < t < \infty , 0 < t < \infty , 0 < t < \infty , 0 < t < \infty , 0 < t < \infty , 0 < t < \infty , 0 < t < \infty , 0 < t < \infty , 0 < t < \infty , 0 < t < \infty , 0 < t < \infty , 0 < t < \infty , 0 < t < \infty , 0 < t < \infty , 0 < t < \infty , 0 < t < \infty , 0 < t < \infty , 0 < t < \infty , 0 < t < \infty , 0 < t < \infty , 0 < t < \infty , 0 < t < \infty , 0 < t < \infty , 0 < t < \infty , , 0 < t < \infty , 0 < t < \infty , 0 < t < \infty , 0 < t < \infty , 0 < t < \infty , 0 < t < \infty , 0 < t < \infty , 0 < t < \infty , 0 < t < \infty , 0 < t < \infty , 0 < t < \infty , 0 < t < \infty , 0 < t < \infty , 0 < t < \infty , 0 < t < \infty , 0 < t < \infty , 0 < t < \infty , 0 < t < \infty , 0 < Щоб перетворити ці граничні умови на нульові умови, ми (після деяких спроб і помилок) зупинилися на такій формі розв'язку:
~~

(6.7)
\end{verbatim}

u(x, t) = A(t){[}1-x/L{]} + B(t){[}x/L{]} + U(x, t),

\begin{verbatim}
~~
де функції A\left(t\right) і B\left(t\right) обираються так, що «квазі-стаціонарна» частина розв'язку (6.7)
~~

(6.8)
\end{verbatim}

S(x, t) = A(t){[}1 - x/L{]} + B(t){[}x/L{]}

\begin{verbatim}
~~
задовольняло граничні умови початкової задачі. У цьому випадку функція U\left(x,\ t\right) задовольнятиме однорідні граничні умови.
~~

Підставляючи функцію S(x, t) у граничні умови
\end{verbatim}

S(0, t) = g\_1(t),

\begin{verbatim}
S_x(L, t) + hS(L, t) = g_2(t) призводить до двох рівнянь, з яких можна визначити A(t) і B(t). У результаті ми отримуємо

(6.9)
\end{verbatim}

\texttt{A(t)\ =\ g\_1(t),\ B(t)\ =\ \textbackslash{}frac\{g\_1(t)\ +\ Lg\_2(t)\}\{1\ +\ Lh\}.}

\begin{verbatim}
Отже.
\end{verbatim}

\texttt{u(x,\ t)\ =\ g\_1(t)\ {[}1\ -\ x/L{]}\ +\ \textbackslash{}frac\{g\_1(t)\ +\ Lg\_2(t)\}\{1\ +\ Lh\}\ {[}x/L{]}\ +\ U(x,\ t).}

\begin{verbatim}
Якщо ми підставимо ці вирази для you(x, t) у початкову задачу (це просять читача), отримаємо нову задачу для невідомої функції U(x, t):

~~
\begin{array}{llllllllllllllllllllllllllllllllllll
~~

Тепер ми стикаємося з новою проблемою однорідних граничних умов (на жаль, рівняння стало гетерогенним). Ми не можемо розв'язати цю задачу методом розділення змінних, але якщо читач погоджується чекати, то за кілька лекцій він ознайомиться з розв'язком цієї задачі методом цілісних перетворень і розкладання за власними функціями.

# ЗАУВАЖЕННЯ .

1. Як зазначалося раніше, мета цієї лекції — навчити читача, як перетворювати задачі з неунірольними граничними умовами на задачі з нульовими граничними умовами. Після цього може статися, що нове диференціальне рівняння в похідних похідних стане однорідним. Якщо це станеться, ми можемо вважати себе щасливцями, оскільки в цьому випадку проблему можна розв'язати методом розділення змінних (див. перший приклад з цієї лекції). Якщо перетворене рівняння виявиться гетерогенним, тоді нову задачу потрібно розв'язати іншим методом.

2. Лінійні неоднорідні граничні умови найзагальної форми
\end{verbatim}

\texttt{\textbackslash{}alpha\_1\ u\_x(0,\ t)\ +\ \textbackslash{}beta\_1\ u(0,\ t)\ =\ g\_1(t),}

\begin{verbatim}
~~
\alpha_2 u_x(L, t) + \beta_2 u(L, t) = g_2(t) Ви також можете перетворити умови на нульові граничні умови, використовуючи метод, описаний у другому прикладі. Звісно, нове рівняння, ймовірно, буде гетерогенним.
~~

- 3. Деякі методи розв'язку не містять вимог щодо унікальності граничних умов і, отже, не потребують попередніх перетворень. Пізніше, коли ми ознайомимося з перетворенням Лапласа, побачимо, що нульові граничні умови є необов'язковими (у випадку нульових граничних умов розв'язок дещо простіший, ніж в інших випадках).

- 4. Для граничних умов типу
\end{verbatim}

u(0, t) = g\_1(t),

\begin{verbatim}
u(L, t) = g_2(t) Метод, описаний у другому прикладі, призводить до наступної форми розв'язання задач:
\end{verbatim}

\texttt{u(x,\ t)\ =\ \textbackslash{}left\textbackslash{}\{\ g\_1(t)\ +\ \textbackslash{}frac\{x\}\{L\}\ \textbackslash{}left{[}\ g\_2(t)\ -\ g\_1(t)\ \textbackslash{}right{]}\ \textbackslash{}right\textbackslash{}\}\ +\ U(x,\ t).}

\begin{verbatim}
# ЗАВДАННЯ

1. Знайти розв'язок наступної змішаної задачі:
\end{verbatim}

\texttt{\textbackslash{}begin\{array\}\{lll\}\ (\textbackslash{}text\{УЧП\})\ \&\ u\_t\ =\ \textbackslash{}alpha\^{}2\ u\_\{xx\}\ +\ x\ \textbackslash{}cos\ t,\ \&\ 0\ \textless{}\ x\ \textless{}\ 1,\ \textbackslash{}\textbackslash{}\ (\textbackslash{}Gamma\ \textbackslash{}text\{У\})\ \&\ \textbackslash{}begin\{cases\}\ u\ \textbackslash{}\ (0,\ \textbackslash{}\ t)\ =\ 1,\ \textbackslash{}\textbackslash{}\ u\_x\ \textbackslash{}\ (1,\ \textbackslash{}\ t)\ +\ hu\ \textbackslash{}\ (1,\ \textbackslash{}\ t)\ =\ 1,\ \textbackslash{}end\{cases\}\ \&\ 0\ \textless{}\ t\ \textless{}\ \textbackslash{}infty,\ \textbackslash{}\textbackslash{}\ (\textbackslash{}text\{HY\})\ \&\ u\ \textbackslash{}\ (x,\ \textbackslash{}\ 0)\ =\ \textbackslash{}sin\ (\textbackslash{}pi\ x),\ \&\ 0\ \textbackslash{}leqslant\ x\ \textbackslash{}leqslant\ 1,\ \textbackslash{}end\{array\}}

\begin{verbatim}
попередньо перетворюючи його на задачу з нульовими граничними умовами. Чи відповідає рішення, яке ви отримуєте, вашим задумам щодо цієї проблеми?

Завдання

(PPP)
\end{verbatim}

\texttt{u\_t\ =\ u\_\{xx\},\ \textbackslash{}quad\ 0\ \textless{}\ x\ \textless{}\ 1,}

\begin{verbatim}
(GI)
\end{verbatim}

\texttt{\textbackslash{}begin\{cases\}\ u(0,\ t)\ =\ 0,\ \&\ 0\ \textless{}\ t\ \textless{}\ \textbackslash{}infty,\ \textbackslash{}\textbackslash{}\ u(1,\ t)\ =\ 1,\ \&\ 0\ \textless{}\ t\ \textless{}\ \textbackslash{}infty,\ \textbackslash{}end\{cases\}}

\begin{verbatim}
~~
(НУ) u(x, 0) = x^2, \quad 0 \le x \le 1, Перетворіть у задачу з нульовими граничними умовами і розв'яжіть її. Як виглядатиме рішення у різних
~~

Миті часу? Чи відповідає це вашим інтуїтивним ідеям? Що таке стаціонарне рішення? Як виглядає перехідне рішення?

3. Завдання
\end{verbatim}

\texttt{\textbackslash{}begin\{aligned\}}
\texttt{(\textbackslash{}text\{УЧП\})\textbackslash{}quad\ \&\ u\_t\ =\ u\_\{xx\},\ \&\&\ 0\ \textless{}\ x\ \textless{}\ 1,\textbackslash{}\ 0\ \textless{}\ t\ \textless{}\ \textbackslash{}infty,\ \textbackslash{}\textbackslash{}}
\texttt{(\textbackslash{}text\{ГУ\})\textbackslash{}quad\ \&\ u\_x(0,t)\ =\ 0,\textbackslash{}quad\ u\_x(1,t)\ +\ h\ u(1,t)\ =\ 1,\ \&\&\ 0\ \textless{}\ t\ \textless{}\ \textbackslash{}infty,\ \textbackslash{}\textbackslash{}}
\texttt{(\textbackslash{}text\{НУ\})\textbackslash{}quad\ \&\ u(x,0)\ =\ \textbackslash{}sin(\textbackslash{}pi\ x),\ \&\&\ 0\ \textbackslash{}leqslant\ x\ \textbackslash{}leqslant\ 1.}
\texttt{\textbackslash{}end\{aligned\}}

\begin{verbatim}
Перетворимо на задачу з нульовими граничними умовами. Чи буде нове рівняння однорідним?

Лекція 7

# РОЗВ'ЯЗАННЯ СКЛАДНІШИХ ЗАДАЧ ШЛЯХОМ РОЗДІЛЕННЯ ЗМІННИХ

МЕТА ЛЕКЦІЇ: Показати, як більш складні проблеми теплопровідності можна розв'язати методом розділення змінних. У цій лекції ми головним чином розглядаємо такі задачі, розв'язання яких дозволить читачеві краще ознайомитися з методом ділення змінних. Сподіваємося, що читач зможе екстраполювати викладені тут ідеї для власних завдань.

У тій же лекції читач ознайомиться з задачею за її власними функціями, або, як її ще називають, проблемою Штурма-Ліувіля. Обговорюються деякі властивості проблеми Штурма-Ліувіля.

Як згадувалося вище, мета цієї лекції — розв'язати змішану задачу шляхом розділення змінних. Можливо, саме таким методом читач вирішить свою проблему. Водночас ми сподіваємося, що він зможе перенести метод, розглянутий у цій лекції, для власних завдань.

Почнемо з одномірної задачі теплопровідності, в якій одна з граничних умов містить похідну.

# Задача теплопровідності з похідною в граничній умові

Розглянемо пристрій, показаний на рис. 7.1. Припустимо, що верхній кінець стрижня підтримується при нульовій температурі, тобто u(0, t) = 0, а нижній кінець занурений у воду, температура якої також стала і дорівнює нулю (у цьому випадку — нуль Р и с

![](images/_page_50_Picture_4.jpeg)

. 7.1. Схематичне представлення зсувної задачі.

![](images/_page_50_Figure_6.jpeg)

Рис. 7.2. Природа кривих, які задовольняють
\end{verbatim}

\texttt{(\textbackslash{}Gamma\ Y)\ =\ \textbackslash{}begin\{cases\}\ u(0,\ t)\ =\ 0,\ \textbackslash{}\textbackslash{}\ u\_x(1,\ t)\ =\ -hu(1,\ t).\ \textbackslash{}end\{cases\}}

\begin{verbatim}
відповідає певній еталонній температурі). Теплообмін між стрижнем і водою відбувається згідно із законом Ньютона, який записується у x=1 як наступна гранична умова:
\end{verbatim}

u\_x(1, t) = -hu(1, t).

\begin{verbatim}
Тепер припустимо, що початкова температура стрижня дорівнює u(x, 0) = x, але одразу при t > 0 включаємо граничні умови. Щоб визначити температуру в стрижні, необхідно розв'язати проблему зміщення

~~
\begin{array}{llll} (\mathrm{УЧ}\Pi) & u_t = \alpha^2 u_{xx}, & 0 < x < 1, & 0 < t < \infty, \\ (7.1) & (\Gamma\mathrm{У}) & \begin{cases} u_t(0, t) = 0 \\ u_x(1, t) + hu_t(1, t) = 0 \end{cases} & \text{(однородные } \Gamma\mathrm{У}), \\ (\mathrm{H}\mathrm{У}) & u_t(x, 0) = x, & 0 \leq x \leq 1. \end{array}
~~

Щоб використати метод розділення змінних, розбиємо процес розв'язання на кілька етапів.

**КРОК** 1. (Перетворення диференціального рівняння в похідних у два звичайних диференціальних рівняння)

~~
Підставляючи u(x, t) = X(x)T(t) у рівняння (7.1), отримуємо XT' = \alpha^2 X''T Поділивши інші частини цього рівняння на \alpha^2 XT отримуємо
~~
\end{verbatim}

\texttt{\textbackslash{}frac\{T\textquotesingle{}\}\{\textbackslash{}alpha\^{}2\ T\}\ =\ \textbackslash{}frac\{X\textquotesingle{}\textquotesingle{}\}\{X\}\ .}

\begin{verbatim}
~~
Оскільки ліва частина останнього рівняння залежить лише від часу, а права — лише від x (і оскільки x і t незалежні), обидві сторони рівняння мають бути константами. Прирівнюючи обидві частини до однієї сталої \mu , отримуємо
~~

(7.2)
\end{verbatim}

\texttt{T\textquotesingle{}\ -\ \textbackslash{}mu\ \textbackslash{}alpha\^{}\{2\}\ T\ =\ 0,\ X\textquotesingle{}\textquotesingle{}\ -\ \textbackslash{}mu\ X\ =\ 0.}

\begin{verbatim}
Процес розщеплення змінних завершено.

КРОК 2. (Визначення константи відокремлення)

~~
По-перше, константа і не повинна бути додатною, інакше ^{1} ) функція T\left( t\right) буде зростати експоненційно, що призведе до нескінченного підвищення температури u=XT. Очевидно, що таке рішення не має фізичного значення.
~~

~~
По-друге, припустимо, що \mu = 0 . Тоді
~~
\end{verbatim}

X'\,' = 0

\begin{verbatim}
і, отже.
\end{verbatim}

\texttt{X(x)\ =\ A\ +\ \textbackslash{}hat\{B\}x.}

\begin{verbatim}
Однак, оскільки граничні умови вимагають, що
\end{verbatim}

u(0, t) = X(0) T(t) = 0,

\begin{verbatim}
u(1, t) = (hX(1) + X'(1)) T(t) = 0, Ми одразу дійшли висновку, що
\end{verbatim}

X(0) = 0
\texttt{\textbackslash{}Rightarrow\ A\ =\ 0\ X\textquotesingle{}(1)\ +\ hX(1)\ =\ 0\ \textbackslash{}Rightarrow\ B\ =\ 0\ .}

\begin{verbatim}
~~
Це означає, що u(x, t) \equiv 0 . Іншими словами, лише нульове розв'язання відповідає нульовій константі розділення. Ми відкинемо його, оскільки нас цікавлять ненульові розв'язки.
~~

~~
Якщо, нарешті, \mu < 0 , можна поставити \mu = -\lambda^2 і переписати рівняння (7.2) у вигляді
~~
\end{verbatim}

\texttt{T\textquotesingle{}\ +\ \textbackslash{}lambda\^{}2\ \textbackslash{}alpha\^{}2\ T\ =\ 0,}
\texttt{X\textquotesingle{}\ +\ \textbackslash{}lambda\^{}2\ X\ =\ 0.}

\begin{verbatim}
<sup>1)</sup> Розв'язок X(x) відповідної задачі з граничним значенням буде лише тривіальним, X(x) = 0. — Ед. Ед.

Загальні розв'язки цих рівнянь легко знайти:
\end{verbatim}

\texttt{T(t)\ =\ Ae\^{}\{-(\textbackslash{}lambda\ \textbackslash{}alpha)\^{}2\ t\},}
\texttt{X(x)\ =\ B\textbackslash{}sin(\textbackslash{}lambda\ x)\ +\ C\textbackslash{}cos(\textbackslash{}lambda\ x).}

\begin{verbatim}
Отже, будь-яка функція провідника

(7.3)
\end{verbatim}

\texttt{u(x,\ t)\ =\ e\^{}\{-(\textbackslash{}lambda\ \textbackslash{}alpha)\^{}2\ t\}\ {[}A\ \textbackslash{}sin(\textbackslash{}lambda\ x)\ +\ B\ \textbackslash{}cos(\textbackslash{}lambda\ x){]}}

\begin{verbatim}
~~
при будь-яких значеннях \lambda A і B задовольняє рівняння (читач може перевірити це самостійно простою заміною).
~~

Тепер нам потрібно з'ясувати, які з цих функцій задовольняють граничним умовам

(7.4)
\end{verbatim}

u(0, t) = 0, u\_x(1, t) + hu(1, t) = 0.

\begin{verbatim}
~~
Підстановка розв'язків виду (7.3) у граничні умови (7.4) дозволяє знайти відношення, яким мають задовольняти значення \lambda , \Lambda та B:
~~
\end{verbatim}

\texttt{Be\^{}\{-(\textbackslash{}lambda\textbackslash{}alpha)\^{}2\ t\}\ =\ 0\ \textbackslash{}Rightarrow\ B\ =\ 0,}

\begin{verbatim}
~~
A\lambda e^{-(\lambda\alpha)^2 t} \cos \lambda + hAe^{-(\lambda\alpha)^2 t} \sin \lambda = 0. Після простих алгебранхічних перетворень останнє рівняння зводиться до вигляду
~~
\end{verbatim}

\texttt{tg\ \textbackslash{}lambda\ =\ -\textbackslash{}lambda/\textbackslash{}hbar}

\begin{verbatim}
.

~~
Іншими словами, щоб визначити \lambda , потрібно знайти точки перетину графів функцій tg \lambda і — \lambda/h (рисунок 7.3).
~~

~~
Числові значення значень \lambda_1 , \lambda_2 , ... можна знайти за допомогою комп'ютера, якщо задано значення h Р и с
~~

![](images/_page_52_Figure_15.jpeg)

~~
. 7.3. Перетин графів функцій tg(\lambda) і -\lambda/h .
~~

викликати власні значення задачі граничного значення

(7.5)
\end{verbatim}

\texttt{X\textquotesingle{}\textquotesingle{}\ +\ \textbackslash{}lambda\^{}\{2\}X\ =\ 0,\ \textbackslash{}begin\{cases\}\ X(0)\ =\ 0\ \textbackslash{}\textbackslash{}\ X\textquotesingle{}(1)\ +\ hX(1)\ =\ 0.\ \textbackslash{}end\{cases\}}

\begin{verbatim}
~~
Іншими словами, з цими значеннями параметра \lambda існують нетривіальні розв'язки задачі (7.5). Власні значення задачі (7.5) визначаються як корені рівняння \log \lambda = -\lambda/h . Їхні числові значення знаходилися на комп'ютері для випадку h = 1. Таблиця 7.1 показує результати розрахунків перших кількох власних значень.
~~

~~
| n | \lambda_B | rt . | \lambda_n |
~~

|-------------|----------------------|--------|----------------|

| 1<br>2<br>3 | 2,02<br>4,91<br>7,98 | 4<br>5 | 11,08<br>14,20 |

~~
Розв'язки задачі (7.5), що відповідають власним значенням \lambda_n , називаються власними функціями X_n(x) . Для розглянутого завдання власні функції виглядають так:
~~
\end{verbatim}

\texttt{X\_n(x)\ =\ \textbackslash{}sin(\textbackslash{}lambda\_n\ x)}

\begin{verbatim}
і зображені на рис. 7.4.

![](images/_page_53_Figure_10.jpeg)

Рис. 7.4. Власні функції X_n(x) завдання (7.5) при h=1; Усі власні функції задовольняють граничні умови.

КРОК 3, (Пошук фундаментальних рішень)

Тепер у нас є нескінченний набір функцій (фундаментальних розв'язків) у нашому розпорядженні
\end{verbatim}

\texttt{u\_n(x,\ t)\ =\ X\_n(x)\ T\_n(t)\ =\ e\^{}\{-(\textbackslash{}lambda\_n\ x)\^{}2\ t\}\ \textbackslash{}sin(\textbackslash{}lambda\_n\ x),}

\begin{verbatim}
кожна з яких задовольняє рівняння та граничні умови. Залишився ще один крок — загалом
\end{verbatim}

\texttt{u(x,\ t)\ =\ \textbackslash{}sum\_\{n=1\}\^{}\{\textbackslash{}infty\}\ a\_n\ X\_n(x)\ T\_n(t)\ =}
\texttt{=\ \textbackslash{}sum\_\{n=1\}\^{}\{\textbackslash{}infty\}\ a\_n\ e\^{}\{-(\textbackslash{}lambda\_n\ c\ t\^{}2\ t)\}\ \textbackslash{}sin(\textbackslash{}lambda\_n\ x)}

\begin{verbatim}
Вибрати коефіцієнти a_n так, щоб функція u(x, t) задовольняла початкову умову u(x, 0) = x (у цьому випадку функція u(x, t) задовольнятиме рівняння та граничні умови, оскільки, по-перше, кожен член суми їх задовольняє, а по-друге, і рівняння, і граничні умови є лінійними та однорідними).

Отже, ймовірність a_n слід визначати з цієї умови

(7 6)
\end{verbatim}

\texttt{u(x,\ 0)\ =\ x\ =\ \textbackslash{}sum\_\{n=1\}\^{}\{\textbackslash{}infty\}\ a\_n\ \textbackslash{}sin(\textbackslash{}lambda\_n\ x).}

\begin{verbatim}
Ми маємо зробити останній крок.

КРОК 4. (Розклад початкової умови на ряд за власними функціями)

Щоб знайти коефіцієнти a_n у розкладі (7,6) на власні функції, помножте обидві частини цього співвідношення на \sin(\lambda_m x) і інтегруйте від 0 до 1. У результаті ми отримуємо
\end{verbatim}

\texttt{\textbackslash{}int\_\{0\}\^{}\{2\}\ \textbackslash{}xi\ \textbackslash{}sin(\textbackslash{}lambda\_\{m\}\ \textbackslash{}xi)\ d\textbackslash{}xi\ =\ \textbackslash{}sum\_\{n=1\}\^{}\{\textbackslash{}infty\}\ a\_\{n\}\ \textbackslash{}int\_\{0\}\^{}\{1\}\ \textbackslash{}sin(\textbackslash{}lambda\_\{n\}\ \textbackslash{}xi)\ \textbackslash{}sin(\textbackslash{}lambda\_\{m\}\ \textbackslash{}xi)\ d\textbackslash{}xi\ =}
\texttt{=\ a\_\{m\}\ \textbackslash{}int\_\{0\}\^{}\{1\}\ \textbackslash{}sin\^{}\{2\}(\textbackslash{}lambda\_\{m\}\ \textbackslash{}xi)\ d\textbackslash{}xi\ =}
\texttt{=\ a\_\{m\}\ \textbackslash{}left(\ \textbackslash{}frac\{\textbackslash{}lambda\_\{m\}\ -\ \textbackslash{}sin\ \textbackslash{}lambda\_\{m\}\ \textbackslash{}cos\ \textbackslash{}lambda\_\{m\}\}\{2\textbackslash{}lambda\_\{m\}\}\ \textbackslash{}right).}

\begin{verbatim}
Розв'язавши останнє відношення відносно a_m і замінивши m на n, отримуємо бажаний результат:

(7.7)
\end{verbatim}

\texttt{a\_n\ =\ \textbackslash{}frac\{2\textbackslash{}lambda\_n\}\{(\textbackslash{}lambda\_n\ -\ \textbackslash{}sin\ \textbackslash{}lambda\_n\ \textbackslash{}cos\ \textbackslash{}lambda\_n)\}\ \textbackslash{}int\_0\^{}\{\textbackslash{}infty\}\ \textbackslash{}xi\ \textbackslash{}sin\ (\textbackslash{}lambda\_n\ \textbackslash{}xi)\ d\textbackslash{}xi.}

\begin{verbatim}
Таблиця 7.2.

# Шанси a_n у формулі (7.8)

| n | an | n | a_n |
|-------------|-----------------------|---------------|----------------|
| 1<br>2<br>3 | 0,24<br>0,22-0,03 <br>| <b>4</b><br>5 | -0.11-0.09 <br>|

Отже, розв'язок задачі (7.1) записується у вигляді

(7.8)
\end{verbatim}

\texttt{u(x,\ t)\ =\ \textbackslash{}sum\_\{n=1\}\^{}\{\textbackslash{}infty\}\ a\_n\ e\^{}\{-i\textbackslash{}lambda\_n\ \textbackslash{}alpha\_n\^{}\{2\}\ t\}\ \textbackslash{}sin(\textbackslash{}lambda\_n\ x),}

\begin{verbatim}
де коефіцієнти a_n задаються формулами (7.7). Перші п'ять коефіцієнтів a_n нашої задачі були знайдені чисельно. Їхні значення наведені в Таблиці 1. 7.2.

![](images/_page_55_Figure_8.jpeg)

Рис. 7.5. Графічне представлення функції (7.8); межа умовно задовольняється у всі моменти часу.

Відповідно, перші три члени розв'язання змішаної задачі

(7.9)
\end{verbatim}

\texttt{\textbackslash{}begin\{aligned\}\ u\_t\ \&=\ u\_\{xx\},\ \&\ 0\ \textless{}\ x\ \textless{}\ 1,\ \&\ 0\ \textless{}\ t\ \textless{}\ \textbackslash{}infty,\ \textbackslash{}\textbackslash{}\ u\_t\ \&=\ 0,\ \&\ 0\ \textless{}\ t\ \textless{}\ \textbackslash{}infty,\ \textbackslash{}\textbackslash{}\ u\_x\ \&=\ 0,\ \&\ 0\ \textless{}\ t\ \textless{}\ \textbackslash{}infty,\ \textbackslash{}\textbackslash{}\ u\_x\ \&=\ 0,\ \&\ 0\ \textless{}\ t\ \textless{}\ \textbackslash{}infty,\ \textbackslash{}\textbackslash{}\ u\_x\ \&=\ 0,\ \&\ 0\ \textless{}\ t\ \textless{}\ \textbackslash{}infty,\ \textbackslash{}\textbackslash{}\ u\_x\ \&=\ 0,\ \&\ 0\ \textless{}\ t\ \textless{}\ \textbackslash{}infty,\ \textbackslash{}\textbackslash{}\ u\_x\ \&=\ 0,\ \&\ 0\ \textless{}\ t\ \textless{}\ \textbackslash{}infty,\ \textbackslash{}\textbackslash{}\ u\_x\ \&=\ 0,\ \&\ 0\ \textless{}\ t\ \textless{}\ \textbackslash{}infty,\ \textbackslash{}\textbackslash{}\ u\_x\ \&=\ 0,\ \&\ 0\ \textless{}\ t\ \textless{}\ \textbackslash{}infty,\ \textbackslash{}\textbackslash{}\ u\_x\ \&=\ 0,\ \&\ 0\ \textless{}\ t\ \textless{}\ \textbackslash{}infty,\ \textbackslash{}\textbackslash{}\ u\_x\ \&=\ 0,\ \&\ 0\ \textless{}\ t\ \textless{}\ \textbackslash{}infty,\ \textbackslash{}\textbackslash{}\ u\_x\ \&=\ 0,\ \&\ 0\ \textless{}\ t\ \textless{}\ \textbackslash{}infty,\ \textbackslash{}\textbackslash{}\ u\_x\ \&=\ 0,\ \&\ 0\ \textless{}\ t\ \textless{}\ \textbackslash{}infty,\ \textbackslash{}\textbackslash{}\ u\_x\ \&=\ 0,\ \&\ 0\ \textless{}\ t\ \textless{}\ \textbackslash{}infty,\ \textbackslash{}\textbackslash{}\ u\_x\ \&=\ 0,\ \&\ 0\ \textless{}\ t\ \textless{}\ \textbackslash{}infty,\ \textbackslash{}\textbackslash{}\ u\_x\ \&=\ 0,\ \&\ 0\ \textless{}\ t\ \textless{}\ \textbackslash{}infty,\ \textbackslash{}\textbackslash{}\ u\_x\ \&=\ 0,\ \&\ 0\ \textless{}\ t\ \textless{}\ \textbackslash{}infty,\ \textbackslash{}\textbackslash{}\ u\_x\ \&=\ 0,\ \&\ 0\ \textless{}\ t\ \textless{}\ \textbackslash{}infty,\ \textbackslash{}\textbackslash{}\ u\_x\ \&=\ 0,\ \&\ 0\ \textless{}\ t\ \textless{}\ \textbackslash{}infty,\ \textbackslash{}\textbackslash{}\ u\_x\ \&=\ 0,\ \&\ 0\ \textless{}\ t\ \textless{}\ \textbackslash{}infty,\ \textbackslash{}\textbackslash{}\ u\_x\ \&=\ 0,\ \&\ 0\ \textless{}\ t\ \textless{}\ \textbackslash{}infty,\ \textbackslash{}\textbackslash{}\ u\_x\ \&=\ 0,\ \&\ 0\ \textless{}\ t\ \textless{}\ \textbackslash{}infty,\ \textbackslash{}\textbackslash{}\ u\_x\ \&=\ 0,\ \&\ 0\ \textless{}\ t\ \textless{}\ \textbackslash{}infty,\ \textbackslash{}\textbackslash{}\ u\_x\ \&=\ 0,\ \&\ 0\ \textless{}\ t\ \textless{}\ \textbackslash{}infty,\ \textbackslash{}\textbackslash{}\ u\_x\ \&=\ 0,\ \&\ 0\ \textless{}\ t\ \textless{}\ \textbackslash{}infty,\ \textbackslash{}\textbackslash{}\ u\_x\ \&=\ 0,\ \&\ 0\ \textless{}\ t\ \textless{}\ \textbackslash{}infty,\ \textbackslash{}\textbackslash{}\ u\_x\ \&=\ 0,\ \&\ 0\ \textless{}\ t\ \textless{}\ \textbackslash{}infty,\ \textbackslash{}\textbackslash{}\ u\_x\ \&=\ 0,\ \&\ 0\ \textless{}\ t\ \textless{}\ \textbackslash{}infty,\ \textbackslash{}\textbackslash{}\ u\_x\ \&=\ 0,\ \&\ 0\ \textless{}\ t\ \textless{}\ \textbackslash{}infty,\ \textbackslash{}\textbackslash{}\ u\_x\ \&=\ 0,\ \&\ 0\ \textless{}\ t\ \textless{}\ \textbackslash{}infty,\ \textbackslash{}\textbackslash{}\ u\_x\ \&=\ 0,\ \&\ 0\ \textless{}\ t\ \textless{}\ \textbackslash{}infty,\ \textbackslash{}\textbackslash{}\ u\_x\ \&=\ 0,\ \&\ 0\ \textless{}\ t\ \textless{}\ \textbackslash{}infty,\ \textbackslash{}\textbackslash{}\ u\_x\ \&=\ 0,\ \&\ 0\ \textless{}\ t\ \textless{}\ \textbackslash{}infty,\ \textbackslash{}\textbackslash{}\ u\_x\ \&=\ 0,\ \&\ 0\ \textless{}\ t\ \textless{}\ \textbackslash{}infty,\ \textbackslash{}\textbackslash{}\ u\_x\ \&=\ 0,\ \&\ 0\ \textless{}\ t\ \textless{}\ \textbackslash{}infty,\ \textbackslash{}\textbackslash{}\ u\_x\ \&=\ 0,\ \&\ 0\ \textless{}\ t\ \textless{}\ \textbackslash{}infty,\ \textbackslash{}\textbackslash{}\ u\_x\ \&=\ 0,\ \&\ 0\ \textless{}\ t\ \textless{}\ \textbackslash{}infty,\ \textbackslash{}\textbackslash{}\ u\_x\ \&=\ 0,\ \&\ 0\ \textless{}\ t\ \textless{}\ \textbackslash{}infty,\ \textbackslash{}\textbackslash{}\ u\_x\ \&=\ 0,\ \&\ 0\ \textless{}\ t\ \textless{}\ \textbackslash{}infty,\ \textbackslash{}\textbackslash{}\ u\_x\ \&=\ 0,\ \&\ 0\ \textless{}\ t\ \textless{}\ \textbackslash{}infty,\ \textbackslash{}\textbackslash{}\ u\_x\ \&=\ 0,\ \&\ 0\ \textless{}\ t\ \textless{}\ \textbackslash{}infty,\ \textbackslash{}\textbackslash{}\ u\_x\ \&=\ 0,\ \&\ 0\ \textless{}\ t\ \textless{}\ \textbackslash{}infty,\ \textbackslash{}\textbackslash{}\ u\_x\ \&=\ 0,\ \&\ 0\ \textless{}\ t\ \textless{}\ \textbackslash{}infty,\ \textbackslash{}\textbackslash{}\ u\_x\ \&=\ 0,\ \&\ 0\ \textless{}\ t\ \textless{}\ \textbackslash{}infty,\ \textbackslash{}\textbackslash{}\ u\_x\ \&=\ 0,\ \&\ 0\ \textless{}\ t\ \textless{}\ \textbackslash{}infty,\ \textbackslash{}\textbackslash{}\ u\_x\ \&=\ 0,\ \&\ 0\ \textless{}\ t\ \textless{}\ \textbackslash{}infty,\ \textbackslash{}\textbackslash{}\ u\_x\ \&=\ 0,\ \&\ 0\ \textless{}\ t\ \textless{}\ \textbackslash{}infty,\ \textbackslash{}\textbackslash{}\ u\_x\ \&=\ 0,\ \&\ 0\ \textless{}\ t\ \textless{}\ \textbackslash{}infty,\ \textbackslash{}\textbackslash{}\ u\_x\ \&=\ 0,\ \&\ 0\ \textless{}\ t\ \textless{}\ \textbackslash{}infty,\ \textbackslash{}\textbackslash{}\ u\_x\ \&=\ 0,\ \&\ 0\ \textless{}\ t\ \textless{}\ \textbackslash{}infty,\ \textbackslash{}\textbackslash{}\ u\_x\ \&=\ 0,\ \&\ 0\ \textless{}\ t\ \textless{}\ \textbackslash{}infty,\ \textbackslash{}\textbackslash{}\ u\_x\ \&=\ 0,\ \&\ 0\ \textless{}\ t\ \textless{}\ \textbackslash{}infty,\ \textbackslash{}\textbackslash{}\ u\_x\ \&=\ 0,\ \&\ 0\ \textless{}\ t\ \textless{}\ \textbackslash{}infty,\ \textbackslash{}\textbackslash{}\ u\_x\ \&=\ 0,\ \&\ 0\ \textless{}\ t\ \textless{}\ \textbackslash{}infty,\ \textbackslash{}\textbackslash{}\ u\_x\ \&=\ 0,\ \&\ 0\ \textless{}\ t\ \textless{}\ \textbackslash{}infty,\ \textbackslash{}\textbackslash{}\ u\_x\ \&=\ 0,\ \&\ 0\ \textless{}\ t\ \textless{}\ \textbackslash{}infty,\ \textbackslash{}\textbackslash{}\ u\_x\ \&=\ 0,\ \&\ 0\ \textless{}\ t\ \textless{}\ \textbackslash{}infty,\ \textbackslash{}\textbackslash{}\ u\_x\ \&=\ 0,\ \&\ 0\ \textless{}\ t\ \textless{}\ \textbackslash{}infty,\ \textbackslash{}\textbackslash{}\ u\_x\ \&=\ 0,\ \&\ 0\ \textless{}\ t\ \textless{}\ \textbackslash{}infty,\ \textbackslash{}\textbackslash{}\ u\_x\ \&=\ 0,\ \&\ 0\ \textless{}\ t\ \textless{}\ \textbackslash{}infty,\ \textbackslash{}\textbackslash{}\ u\_x\ \&=\ 0,\ \&\ 0\ \textless{}\ t\ \textless{}\ \textbackslash{}infty,\ \textbackslash{}\textbackslash{}\ u\_x\ \&=\ 0,\ \&\ 0\ \textless{}\ t\ \textless{}\ \textbackslash{}infty,\ \textbackslash{}\textbackslash{}\ u\_x\ \&=\ 0,\ \&\ 0\ \textless{}\ t\ \textless{}\ \textbackslash{}infty,\ \textbackslash{}\textbackslash{}\ u\_x\ \&=\ 0,\ \&\ 0\ \textless{}\ t\ \textless{}\ \textbackslash{}infty,\ \textbackslash{}\textbackslash{}\ u\_x\ \&=\ 0,\ \&\ 0\ \textless{}\ t\ \textless{}\ \textbackslash{}infty,\ \textbackslash{}\textbackslash{}\ u\_x\ \&=\ 0,\ \&\ 0\ \textless{}\ t\ \textless{}\ \textbackslash{}infty,\ \textbackslash{}\textbackslash{}\ u\_x\ \&=\ 0,\ \&\ 0\ \textless{}\ t\ \textless{}\ \textbackslash{}infty,\ \textbackslash{}\textbackslash{}\ u\_x\ \&=\ 0,\ \&\ 0\ \textless{}\ t\ \textless{}\ \textbackslash{}infty,\ \textbackslash{}\textbackslash{}\ u\_x\ \&=\ 0,\ \&\ 0\ \textless{}\ t\ \textless{}\ \textbackslash{}infty,\ \textbackslash{}\textbackslash{}\ u\_x\ \&=\ 0,\ \&\ 0\ \textless{}\ t\ \textless{}\ \textbackslash{}infty,\ \textbackslash{}\textbackslash{}\ u\_x\ \&=\ 0,\ \&\ 0\ \textless{}\ t\ \textless{}\ \textbackslash{}infty,\ \textbackslash{}\textbackslash{}\ u\_x\ \&=\ 0,\ \&\ 0\ \textless{}\ t\ \textless{}\ \textbackslash{}infty,\ \textbackslash{}\textbackslash{}\ u\_x\ \&=\ 0,\ \&\ 0\ \textless{}\ t\ \textless{}\ \textbackslash{}infty,\ \textbackslash{}\textbackslash{}\ u\_x\ \&=\ 0,\ \&\ 0\ \textless{}\ t\ \textless{}\ \textbackslash{}infty,\ \textbackslash{}\textbackslash{}\ u\_x\ \&=\ 0,\ \&\ 0\ \textless{}\ t\ \textless{}\ \textbackslash{}infty,\ \textbackslash{}\textbackslash{}\ u\_x\ \&=\ 0,\ \&\ 0\ \textless{}\ t\ \textless{}\ \textbackslash{}infty,\ \textbackslash{}\textbackslash{}\ u\_x\ \&=\ 0,\ \&\ 0\ \textless{}\ t\ \textless{}\ \textbackslash{}infty,\ \textbackslash{}\textbackslash{}\ u\_x\ \&=\ 0,\ \&\ 0\ \textless{}\ t\ \textless{}\ \textbackslash{}infty,\ \textbackslash{}\textbackslash{}\ u\_x\ \&=\ 0,\ \&\ 0\ \textless{}\ t\ \textless{}\ \textbackslash{}infty,\ \textbackslash{}\textbackslash{}\ u\_x\ \&=\ 0,\ \&\ 0\ \textless{}\ t\ \textless{}\ \textbackslash{}infty,\ \textbackslash{}\textbackslash{}\ u\_x\ \&=\ 0,\ \&\ 0\ \textless{}\ t\ \textbackslash{}end\{aligned\}}

\begin{verbatim}
мають вигляд
\end{verbatim}

\texttt{u(x,\ t)\ =\ =\ 0.24e\^{}\{-4t\}\ \textbackslash{}sin(2.02x)\ +\ 0.22e\^{}\{-24t\}\ \textbackslash{}sin(4.9x)\ -\ 0.03e\^{}\{-48.3t\}\ \textbackslash{}sin(7.98x)\ +\ \textbackslash{}dots}

\begin{verbatim}
Фіг. На рисунку 7-5 показані графіки цього розв'язку для кількох точок часу. Рекомендуємо читачеві відповісти на такі питання: Чи відповідає тип рішення, яке ви отримуєте, з вашими інтуїтивними ідеями? Чи виконуються граничні умови?

# ЗАУВАЖЕННЯ

Задача власних значень (7.5) є особливим випадком більш загальної задачі, яку зазвичай називають проблемою Штурма-Ліувіля:

(7.10)
\end{verbatim}

\texttt{\textbackslash{}begin\{aligned\}\ \&\ (\textbackslash{}text\{ОДУ\})\ \textbackslash{}quad\ {[}p(x)y\textquotesingle{}{]}\textquotesingle{}\ -\ q(x)y\ +\ \textbackslash{}lambda\ r(x)y\ =\ 0,\ \textbackslash{}quad\ 0\ \textless{}\ x\ \textless{}\ 1,\ \textbackslash{}\textbackslash{}\ \&\ (\textbackslash{}alpha\_1\ y(0)\ +\ \textbackslash{}beta\_1\ y\textquotesingle{}(0)\ =\ 0,\ \textbackslash{}\textbackslash{}\ \&\ (\textbackslash{}alpha\_2\ y(1)\ +\ \textbackslash{}beta\_1\ y\textquotesingle{}(1)\ =\ 0.\ \textbackslash{}end\{aligned\}}

\begin{verbatim}
Якщо розв'язати задачу з однорідними граничними умовами методом розділення змінних, то функції X(x) завжди знаходяться як розв'язки деякої задачі Штурма–Ліувіля. Очевидно, що задача власних значень (7.5) є особливим випадком задачі (7.10).

Перелічимо основні властивості задачі (7.10), які були встановлені Штурмом і Ліувілем за певних природних обмежень.

функції p(x), q(x) та r(x).

1. Існує нескінченна послідовність власних значень, що задовольняють нерівності
\end{verbatim}

\texttt{\textbackslash{}lambda\_1\ \textless{}\ \textbackslash{}lambda\_2\ \textless{}\ \textbackslash{}lambda\_3\ \textless{}\ \textbackslash{}ldots\ \textless{}\ \textbackslash{}lambda\_n\ \textless{}\ \textbackslash{}ldots\ \textbackslash{}to\ \textbackslash{}infty}

\begin{verbatim}
.

2. Кожне власне значення \lambda_n відповідає одній (точній до константного множника) ненульовій власній функції y_n(x) .

3. Якщо y_n(x) і y_m(x) дві різні власні функції, що відповідають різним власним значенням \lambda_n і \lambda_m , то вони є ортогональними з вагою r(x) на сегменті [0, 1], тобто вони задовольняють це відношення
\end{verbatim}

\texttt{\textbackslash{}int\_\{0\}\^{}\{1\}\ r(x)\ y\_\{n\}(x)\ y\_\{m\}(x)\ dx\ =\ 0.}

\begin{verbatim}
Більше деталей про проблему Штурма-Ліувіля можна знайти в літературі, наведеній наприкінці книги.

<sup>1)</sup> Див., наприклад, Тіхоног А. Н., Самарський А. А. Рівняння математичної фізики. — Москва: Наука, 1972. — ред. Ед.

# ЗАВДАННЯ

1. Розв'язати наступну задачу для рівняння теплопровідності методом розкладу

(PPP)
\end{verbatim}

\texttt{u\_1\ =\ u\_\{xx\},\ \textbackslash{}quad\ 0\ \textless{}\ x\ \textless{}\ 1,\ \textbackslash{}quad\ 0\ \textless{}\ t\ \textless{}\ \textbackslash{}infty,}

\begin{verbatim}
(GI)
\end{verbatim}

\texttt{\textbackslash{}begin\{cases\}\ u(0,\ t)\ =\ 0,\ \textbackslash{}\textbackslash{}\ u\_x(1,\ t)\ =\ 0,\ \textbackslash{}\textbackslash{}\ u(x,\ 0)\ =\ x,\ \textbackslash{}quad\ 0\ \textbackslash{}le\ x\ \textbackslash{}le\ 1.\ \textbackslash{}end\{cases\}}

\begin{verbatim}
Чи відповідає рішення, яке ви знаходите, з інтонативними представленнями? Що таке стаціонарне рішення?

2. Знайдіть власні цінності та функції проблеми Ассолта-Ліувіля

(ODY)
\end{verbatim}

\texttt{X\textquotesingle{}\ +\ \textbackslash{}lambda\ X\ =\ 0,\ \textbackslash{}quad\ 0\ \textless{}\ x\ \textless{}\ 1,}

\begin{verbatim}
(FY)
\end{verbatim}

\texttt{\textbackslash{}begin\{cases\}\ X(0)\ =\ 0,\ \textbackslash{}\textbackslash{}\ X\textquotesingle{}(1)\ =\ 0.\ \textbackslash{}end\{cases\}}

\begin{verbatim}
Задайте специфічні вирази для функцій p(x), q(x) та r(x) із загальної задачі Штурма–Ліувіля для цього випадку.

3. Розв'язати задачу теплопровідності для планки з термоізольованими межами

(PPP)
\end{verbatim}

\texttt{u\_t\ =\ u\_\{xx\},\ \textbackslash{}quad\ 0\ \textless{}\ x\ \textless{}\ 1,\ \textbackslash{}quad\ 0\ \textless{}\ t\ \textless{}\ \textbackslash{}infty,}

\begin{verbatim}
(GI)
\end{verbatim}

\texttt{\textbackslash{}begin\{cases\}\ u\_x\ (0,\ t)\ =\ 0,\ \textbackslash{}\textbackslash{}\ u\_x\ (1,\ t)\ =\ 0,\ \textbackslash{}end\{cases\}\ \textbackslash{}quad\ 0\ \textless{}\ t\ \textless{}\ \textbackslash{}infty,}

\begin{verbatim}
(ГУ) u(x, 0) = x, \quad 0 \le x \le 1. Чи відповідає рішення, яке ви отримуєте, з вашою інтерпретацією цієї проблеми? Що таке стаціонарне рішення? Яке його фізичне значення?

4. Знайдіть власні цінності та функції завдань

(ОДА)
\end{verbatim}

\texttt{X\textquotesingle{}\textquotesingle{}\ +\ \textbackslash{}lambda\ X\ =\ 0,\ \textbackslash{}quad\ 0\ \textless{}\ x\ \textless{}\ 1,}

\begin{verbatim}
(GI)
\end{verbatim}

\texttt{\textbackslash{}begin\{cases\}\ X\textquotesingle{}(0)\ =\ 0,\ \textbackslash{}\textbackslash{}\ X\textquotesingle{}(1)\ =\ 0.\ \textbackslash{}end\{cases\}}

\begin{verbatim}
## ПЕРЕТВОРЕННЯ СКЛАДНИХ РІВНЯНЬ У ПРОСТУ ФОРМУ

МЕТА ЛЕКЦІЇ Показати, як PPP для функції u(x, t) можна перетворити на простіший PPP для нової функції w(x, t). Вибір правильної трансформації зазвичай базується на ненавмисному намірі. У нашій лекції
\end{verbatim}

\texttt{u\_t\ =\ \textbackslash{}alpha\^{}2\ u\_\{xx\}\ -\ \textbackslash{}beta\ u,}
\texttt{u\_t\ =\ \textbackslash{}alpha\^{}2\ u\_\{xx\}\ -\ v\ u\_x,\ \textbackslash{}quad\ (v\ =\ \textbackslash{}text\{const\})}

\begin{verbatim}
переходимо до найпростішого рівняння теплопровідності
\end{verbatim}

\texttt{w\_1\ =\ \textbackslash{}alpha\^{}2\ w\_\{xx\}}

\begin{verbatim}
внаслідок відповідних перетворень
\end{verbatim}

\texttt{u(x,\ t)\ =\ e\^{}\{-\textbackslash{}beta\ t\}\ w(x,\ t)}
\texttt{u(x,\ t)\ =\ e\^{}\{\textbackslash{}epsilon\ t\ (1\ -\ \textbackslash{}delta\ t\ -\ 2)\ 2\textbackslash{}alpha\^{}2\}\ w(x,\ t).}

\begin{verbatim}
Розв'язавши просте рівняння теплопровідності за функцією \omega(x,t) , ми можемо знайти розв'язки початкових рівнянь, перетворивши функцію \omega(x,t) відповідно до цих формул (звісно, ці перетворення слід застосовувати і до початкових і граничних умов).

Після прочитання двох останніх лекцій читач може створити враження, що метод розділення змінних застосовується *лише* до рівняння
\end{verbatim}

\texttt{u\_1\ =\ \textbackslash{}alpha\^{}2\ u\_\{xx\}}

\begin{verbatim}
Насправді рівняння теплопотоку є найпростішим, але не єдиним параболічним рівнянням, яке можна розв'язати методом розділення змінних. Як уже згадувалося, якщо рівняння є лінійним і однорідним, його можна розв'язати, поділивши змінні. Наприклад, двовимірна задача теплопровідності в колі описується рівнянням
\end{verbatim}

\texttt{u\_t\ =\ \textbackslash{}alpha\^{}2\ \textbackslash{}left{[}\ u\_\{rr\}\ +\ \textbackslash{}frac\{1\}\{r\}\ u\_r\ +\ \textbackslash{}frac\{1\}\{r\^{}2\}\ u\_\{\textbackslash{}theta\textbackslash{}theta\}\ \textbackslash{}right{]},}

\begin{verbatim}
І хоча це рівняння з *змінними коефіцієнтами*, його можна звести до трьох звичайних диференціальних рівнянь.

У цій лекції ми покажемо читачеві, що іноді, замість того, щоб безпосередньо розв'язувати диференціальне рівняння в похідних, краще перетворити його на простіше рівняння, яке можна розв'язати методом розділення змінних або іншим методом. Наступний приклад ілюструє цей підхід.

# Перетворення проблеми з перенесенням бічної поверхні у проблему з термічно ізольованою бічною поверхнею

Розглянемо наступне завдання:

(8.1)
\end{verbatim}

\texttt{u\_\{t\}\ =\ \textbackslash{}alpha\^{}\{2\}u\_\{xx\}\ -\ \textbackslash{}beta\ u,\ \textbackslash{}quad\ 0\ \textless{}\ x\ \textless{}\ 1,\ \textbackslash{}quad\ 0\ \textless{}\ t\ \textless{}\ \textbackslash{}infty,}
\texttt{\textbackslash{}begin\{cases\}\ u(0,\ t)\ =\ 0,\ \textbackslash{}\textbackslash{}\ u(1,\ t)\ =\ 0,\ \textbackslash{}\textbackslash{}\ u(x,\ 0)\ =\ \textbackslash{}varphi(x),\ \textbackslash{}quad\ 0\ \textbackslash{}le\ x\ \textbackslash{}le\ 1,\ \textbackslash{}end\{cases\}}

\begin{verbatim}
(8.1)
\end{verbatim}

\texttt{(FY)\ \textbackslash{}quad\ \textbackslash{}begin\{cases\}\ u(0,\ t)\ =\ 0,\ \textbackslash{}\textbackslash{}\ u(x,\ 0)\ =\ \textbackslash{}varphi(x),\ \textbackslash{}quad\ 0\ \textbackslash{}le\ x\ \textbackslash{}le\ 1,\ \textbackslash{}end\{cases\}}

\begin{verbatim}
~~
де термін -\beta u описує потік тепла через бічну поверхню стрижня (рис. 8.1).
~~

![](images/_page_59_Figure_7.jpeg)

~~
Рис. 81. Задача теплопровідності, описана рівнянням u_t = \alpha^t u_{\lambda\lambda} - \beta u : \alpha — це втрата гепле через поверхню Боха; \delta – дифузія вздовж стрижня.
~~

~~
Наша мета — ввести нову температурну w\left(x,\ t\right) замість u\left(x,\ t\right) , щоб рівняння для w\left(x,\ t\right) стало простішим за початкове рівняння
~~
\end{verbatim}

\texttt{u\_t\ =\ \textbackslash{}alpha\^{}2\ u\_\{xx\}\ -\ \textbackslash{}beta\ u}

\begin{verbatim}
.

Це досить загальний метод розв'язання диференціальних рівнянь у частинних похідних, але вибір відповідного перетворення базується на інтуїтивних уявленнях про поведінку розв'язків початкового рівняння. Наприклад, у задачі (8.1) температура u(x, t) у кожній точці x_0 змінюється внаслідок двох факторів:

~~
- 1) дифузія тепла вздовж стрижня (через компонентний \alpha^2 u_{xx} ),
~~

- 2) теплопередача через бічну поверхню стрижня (через член RI).

~~
Найважливішим елементом нашого міркування є наступне: якщо не було дифузії вздовж стрижня ( \alpha = 0 ),
~~

тоді температура в кожній точці стрижня експоненційно знижується до нуля згідно із законом
\end{verbatim}

\texttt{u(x\_0\ t)\ =\ u(x\_0,\ 0)\ e\^{}\{-\textbackslash{}beta\ t\}}

\begin{verbatim}
.

Виходячи з цього спостереження, запитаємо, чи можна розв'язок проблеми (8.1) представити як утворення двох факторів
\end{verbatim}

\texttt{(8.2)\ u(x,\ t)\ =\ e\^{}\{-\textbackslash{}beta\ t\}w(x,\ t)}

\begin{verbatim}
або:

~~
\begin{pmatrix}
~~

Температура елемента з \\

Неізольовані боком

Поверхня

де w(x, t) описує розподіл температури, що виникає лише через дифузію. Давайте подивимось, що станеться, якщо підставити цей вираз для температури у задачу (8.1). Після виконання відповідних розрахунків (читач може зробити їх сам), ми отримаємо

~~
Нова задача майже точно така ж, як і оригінальна, але більше не має додатку 6u. Тому замість розв'язання задачі (8.1) ми розв'язуємо перетворену задачу (8.3), а потім множимо отриманий розв'язок w(x, t) на e^{-\beta t} . У нашому випадку розв'язок задачі (8.3) вже було отримано раніше методом розділення змінних, тому ми можемо одразу записати
~~

(8.4) \end{pmatrix}
\end{verbatim}

\texttt{w(x,\ t)\ =\ \textbackslash{}sum\_\{n=1\}\^{}\{\textbackslash{}infty\}\ a\_n\ e\^{}\{-(n\textbackslash{}pi\textbackslash{}alpha)\^{}\{\textbackslash{}frac\{n\}\{2\}\}t\}\ \textbackslash{}sin(n\textbackslash{}pi\ x),}
\texttt{a\_n\ =\ 2\ \textbackslash{}int\_0\^{}\{\textbackslash{}infty\}\ \textbackslash{}phi(\textbackslash{}xi)\ \textbackslash{}sin(n\textbackslash{}pi\ \textbackslash{}xi)\ d\textbackslash{}xi,}

\begin{verbatim}
і, отже, розв'язок початкової задачі (8.1) має вигляд
\end{verbatim}

\texttt{u(x,\ t)\ =\ e\^{}\{-\textbackslash{}beta\ t\}w(x,\ t).}

\begin{verbatim}
Нижче, у розділі «Нотатки», наведено ще один приклад проблеми, розв'язаної цим методом.

# ПРИМІТКИ

1. Вирішити проблему
\end{verbatim}

\texttt{\textbackslash{}begin\{aligned\}}
\texttt{u\_t\ \&=\ u\_\{xx\}\ -\ u,\ \&\&\ 0\ \textless{}\ x\ \textless{}\ 1,\textbackslash{}\ 0\ \textless{}\ t\ \textless{}\ \textbackslash{}infty,\ \textbackslash{}\textbackslash{}}
\texttt{u(0,t)\ \&=\ 0,\textbackslash{}quad\ u(1,t)\ =\ 0,\ \&\&\ 0\ \textless{}\ t\ \textless{}\ \textbackslash{}infty,\ \textbackslash{}\textbackslash{}}
\texttt{u(x,0)\ \&=\ \textbackslash{}sin(\textbackslash{}pi\ x)\ +\ 0.5\textbackslash{}sin(3\textbackslash{}pi\ x),\ \&\&\ 0\ \textbackslash{}leqslant\ x\ \textbackslash{}leqslant\ 1.}
\texttt{\textbackslash{}end\{aligned\}}

\begin{verbatim}
За методом, описаним вище, ми

- a) тимчасово відкинути конвективний член u у початковому рівнянні.
- b) записати розв'язок змішаної задачі без конвективного члена
\end{verbatim}

\texttt{u(x,\ t)\ =\ e\^{}\{-\textbackslash{}pi\^{}2\ t\}\ \textbackslash{}sin(\textbackslash{}pi\ x)\ +\ 0.5e\^{}\{-(3\textbackslash{}pi)\^{}2\ t\}\ \textbackslash{}sin(3\textbackslash{}pi\ x),}

\begin{verbatim}
c) помножимо це розв'язання на конвективний фактор e^{-\beta t}=e^{-t} отримаємо остаточне розв'язання початкової задачі
\end{verbatim}

\texttt{u(x,\ t)\ =\ e\^{}\{-t\}\ \textbackslash{}left{[}\ e\^{}\{-\textbackslash{}pi\^{}2\ t\}\ \textbackslash{}sin(\textbackslash{}pi\ x)\ +\ 0.5e\^{}\{-(3\textbackslash{}pi)\^{}2\ t\}\ \textbackslash{}sin(3\textbackslash{}pi\ x)\ \textbackslash{}right{]}.}

\begin{verbatim}
2. Рівняння конвективної дифузії
\end{verbatim}

\texttt{u\_t\ =\ \textbackslash{}alpha\^{}2\ u\_\{xx\}\ -\ v\ u\_x}

\begin{verbatim}
(коефіцієнт і є константою) також можна перетворити у вигляд
\end{verbatim}

\texttt{w\_t\ =\ \textbackslash{}alpha\^{}2\ w\_\{xx\}.}

\begin{verbatim}
Відповідне перетворення має вигляд
\end{verbatim}

\texttt{u(x,\ t)\ =\ e\^{}\{v(x-vt/2)/2\textbackslash{}alpha\^{}2\}w(x,\ t).}

\begin{verbatim}
Це перетворення підкреслює експоненціальний фактор, який повністю зумовлений рухом середовища. Зверніть увагу, що цей множник описує хвилю, що поширюється зліва направо зі швидкістю \upsilon/2 . Читач матиме можливість скористатися цією трансформацією.

# ЗАВДАННЯ

1. Розв'язати проблему конвективної дифузії

~~
\begin{array}{llll} (\mathrm{Y}\mathrm{H}\Pi) & u_t = u_{xx} - u_x, & 0 < x < 1, & 0 < t < \infty, \\ (\mathrm{\Gamma}\mathrm{Y}) & \left\{ \begin{array}{ll} u \ (0, \ t) = 0, & 0 < t < \infty, \\ u \ (1, \ t) = 0, & 0 < t < \infty, \\ u \ (x, \ 0) = e^{x/2}, & 0 \leqslant x \leqslant 1, \end{array} \end{array}
~~

~~
\Rightarrow. \\ \end{array}
~~

метод перетворення у простіше рівняння. Яке виглядає рішення? Можна дати таку інтерпретацію цієї задачі: u(x, t) — це концентрація певної речовини у рухомому середовищі (середовище рухається зліва направо зі швидкістю v = 1), початковий розподіл копконцентрації задається функцією e^{x/2} , концентрація підтримується на межах на нулі (можна припустити, що там є фільтри). Чи відповідає ваше рішення цій інтерпретації?

2. Решите залачу

~~
\begin{array}{llll} \text{(YUII)} & u_t = u_{xx} - u + x, & 0 < x < 1, & 0 < t < \infty, \\ \text{(FY)} & \left\{ \begin{array}{ll} u \ (0, \ t) = 0, \\ u \ (1, \ t) = 1, \\ u \ (x, \ 0) = 0, & 0 \le x \le 1, \end{array}
~~

~~
\Rightarrow. \\ \end{array}
~~

- а) преобразовав исоднородные граничные условия к однородным;
- б) преобразовав исходное уравнение к повому уравнению, не содержащему члена -u;
- в) решив получившуюся в результате эт іх преобразований задачу.
- 3. Решите задачу
\end{verbatim}

\texttt{\textbackslash{}begin\{array\}\{lll\}\ (\textbackslash{}text\{Y\ 4II\})\ \&\ u\_t\ =\ u\_\{\textbackslash{}text\{vx\}\}\ -\ u,\ \&\ 0\ \textless{}\ x\ \textless{}\ 1,\ \&\ 0\ \textless{}\ t\ \textless{}\ \textbackslash{}infty,\ \textbackslash{}\textbackslash{}\ (\textbackslash{}text\{\textbackslash{}Gamma\ Y\})\ \&\ \textbackslash{}begin\{cases\}\ u(0,\textbackslash{},t)\ =\ 0,\ \textbackslash{}\textbackslash{}\ u(1,\textbackslash{},t)\ =\ 0,\ \textbackslash{}\textbackslash{}\ u(x,\textbackslash{},0)\ =\ \textbackslash{}sin\{(\textbackslash{}pi\ x)\},\ \&\ 0\ \textbackslash{}leqslant\ x\ \textbackslash{}leqslant\ 1,\ \textbackslash{}end\{cases\}\ \textbackslash{}end\{array\}}

\begin{verbatim}
применяя метод поділу змінних непосредственно к исходному уравнению и не проводя никаких предварительных преобразований. Совпадает ли полученное вами розв'язок с решением, полученным ранее після преобразования
\end{verbatim}

u(x, t) = e\^{}\{-t\}w(x, t)

\begin{verbatim}
?

# РОЗВ'ЯЗОК НЕОДНОРОДНЫХ УЧП МЕТОДОМ РАЗЛОЖЕНИЯ ПО СОБСТВЕННЫМ ФУНКЦИЯМ.

ЦЕЛЬ ЛЕКЦИИ: Научиться решать смешанные задачи вида

~~
\begin{array}{llll} \text{(YUII)} & u_t = \alpha^2 u_{xx} + f(x,\ t), & 0 < x < 1, & 0 <t < \infty, \\ \text{(FY)} & \left\{ \begin{array}{ll} \alpha_1 u_x \left( 0,\ t \right) + \beta_1 u \left( 0,\ t \right) = 0, \\ \alpha_2 u_x \left( 1,\ t \right) + \beta_2 u \left( 1,\ t \right) = 0, \\ u \left( x,\ 0 \right) = q \left( x \right), & 0 \leqslant x \leqslant 1, \end{array} \end{array}
~~

~~
\right. \\ \end{array}
~~

Розв'язок гетерогенного UBP можна знайти у вигляді серії
\end{verbatim}

\texttt{u\textbackslash{}left(\ x,\textbackslash{}\ t\textbackslash{}right)\ =\textbackslash{}sum\_\{n=1\}\^{}\{\textbackslash{}infty\}T\_\{n\}\textbackslash{}left(\ t\textbackslash{}right)\ X\_\{n\}\textbackslash{}left(\ x\textbackslash{}right)\ ,}

\begin{verbatim}
~~
де X_n\left(x\right) — власні функції відповідної однорідної задачі
~~

~~
\begin{array}{llll} (\text{YH}\Pi) & u_t = \alpha^2 u_{x_1}, & 0 < x < 1, & 0 < t < \infty, \\ (\text{\GammaV}) & \left\{ \begin{array}{ll} \alpha_1 u_x (0, \ t) + \beta_1 u \ (0, \ t) = 0, \\ \alpha_2 u_x \ (1, \ t) + \beta_2 u \ (1, \ t) = 0, \\ u \ (x, \ 0) = \phi \ (x), & 0 \leqslant x \leqslant 1, \end{array}
~~

~~
\right. \\ \end{array}
~~

А функції T_n(t) визначаються як розв'язки деяких звичайних диференціальних рівнянь.

На лекції 6 ми ознайомилися з тим, як перетворювати неоднорідні граничні умови на однорідні. На жаль, рівняння залишається гетерогенним, і ми стоїмо перед завданням
\end{verbatim}

\texttt{(9.1)\ \textbackslash{}quad\ u\_\{t\}\ =\ \textbackslash{}alpha\^{}\{2\}\ u\_\{xx\}\ +\ f(x,\ t),\ \textbackslash{}quad\ 0\ \textless{}\ x\ \textless{}\ 1,\ \textbackslash{}quad\ 0\ \textless{}\ t\ \textless{}\ \textbackslash{}infty,}
\texttt{(9.1)\ \textbackslash{}quad\ (\textbackslash{}Gamma\ Y)\ \textbackslash{}quad}
\texttt{\textbackslash{}begin\{cases\}\ \textbackslash{}alpha\_\{1\}\ u\_\{x\}(0,\ t)\ +\ \textbackslash{}beta\_\{1\}\ u(0,\ t)\ =\ 0,\ \textbackslash{}\textbackslash{}\ \textbackslash{}alpha\_\{2\}\ u\_\{x\}(1,\ t)\ +\ \textbackslash{}beta\_\{2\}\ u(1,\ t)\ =\ 0,\ \textbackslash{}\textbackslash{}\ u(x,\ 0)\ =\ \textbackslash{}varphi(x),\ \textbackslash{}quad\ 0\ \textbackslash{}le\ x\ \textbackslash{}le\ 1.\ \textbackslash{}end\{cases\}}

\begin{verbatim}
Мета нашої лекції — опанувати метод розв'язання таких задач. Цей метод аналогічний методу варіації довільних констант.

теорія звичайних диференціальних рівнянь називається методом розкладу за власними функціями.

Ідея методу надзвичайно проста. Відомо, що розв'язок задачі (9.1) при f(x, t) = 0 (так звана відповідна однорідна задача) записується як ряд
\end{verbatim}

\texttt{u\textbackslash{}left(x,\textbackslash{}\ t\textbackslash{}right)=\textbackslash{}sum\_\{n=1\}\^{}\{\textbackslash{}infty\}a\_\{n\}e\^{}\{-\textbackslash{}left(\textbackslash{}lambda\_\{n\}\textbackslash{}alpha\textbackslash{}right)\^{}\{2\}\}/X\_\{n\}\textbackslash{}left(x\textbackslash{}right),}

\begin{verbatim}
де X_n(x) — власні функції задачі Штурма-Ліувіля,

(9.2)
\end{verbatim}

\texttt{X\textquotesingle{}\textquotesingle{}\ +\ \textbackslash{}lambda\^{}2\ X\ =\ 0,\ \textbackslash{}\textbackslash{}\ \textbackslash{}alpha\_1\ X\textquotesingle{}(0)\ +\ \textbackslash{}beta\_1\ X(0)\ =\ 0,\ \textbackslash{}\textbackslash{}\ \textbackslash{}alpha\_2\ X\textquotesingle{}(1)\ +\ \textbackslash{}beta\_2\ X(1)\ =\ 0.}

\begin{verbatim}
Запитаймо себе, чи можна також шукати розв'язок гетерогенної проблеми (9.1) у вигляді серії власних функцій, але більш загальної форми
\end{verbatim}

\texttt{u(x,\ t)\ =\ \textbackslash{}sum\_\{n=1\}\^{}\{\textbackslash{}infty\}\ T\_n(t)\ X\_n(x).}

\begin{verbatim}
Існують певні фізичні підстави для такого припущення. Насправді терміни f(x,t) описують джерело тепла, розташоване всередині стрижня. Це означає, що залежність рішення від часу не буде виражена показниками
\end{verbatim}

\texttt{e\^{}\{-(\textbackslash{}lambda\_n\textbackslash{}alpha)\^{}2\ t\}}

\begin{verbatim}
,

Це стосувалося однорідного завдання, але з деякими іншими функціями.

Щоб показати цей метод у дії і не ускладнювати подання деталями, ми застосуємо його до розв'язання однієї дуже простої задачі.

# Построзв'язання методом демаркації за власними функціями,

Розглянемо гетерогенну проблему

(9.3) (PY)
\end{verbatim}

\texttt{\textbackslash{}begin\{aligned\}\ u\_t\ \&=\ \textbackslash{}alpha\^{}2\ u\_\{xx\}\ +\ f(x,\ t),\ \&\ 0\ \textless{}\ x\ \textless{}\ 1,\ \&\ 0\ \textless{}\ t\ \textless{}\ \textbackslash{}infty,\ \textbackslash{}\textbackslash{}\ u(0,\ t)\ \&=\ 0,\ \&\ 0\ \textless{}\ t\ \textless{}\ \textbackslash{}infty,\ \textbackslash{}\textbackslash{}\ u(1,\ t)\ \&=\ 0,\ \&\ 0\ \textless{}\ t\ \textless{}\ \textbackslash{}infty,\ \textbackslash{}\textbackslash{}\ u(x,\ 0)\ \&=\ \textbackslash{}varphi(x),\ \&\ 0\ \textbackslash{}leqslant\ x\ \textbackslash{}leqslant\ 1.\ \textbackslash{}end\{aligned\}}

\begin{verbatim}
Процес розв'язання задачі подається у вигляді послідовності наступних кроків.

КРОК 1. Основна ідея методу полягає в розкладанні густини джерела f(x, t) на ряд власними функціями
\end{verbatim}

\texttt{f(x,\ t)\ =\ f\_1(t)\ X\_1(x)\ +\ f\_2(t)\ X\_2(x)\ +\ \textbackslash{}dots\ +\ f\_n(t)\ X\_n(x)\ +\ \textbackslash{}dots}

\begin{verbatim}
а також визначити реакції системи u_n\left(x,\,t\right) на вплив кожного компонента f_n\left(t\right)X_n\left(x\right) . Якщо знайти всі відповіді, рішення проблеми виглядатиме так:
\end{verbatim}

\texttt{u(x,\ t)\ =\ \textbackslash{}sum\_\{n=1\}\^{}\{\textbackslash{}infty\}\ u\_n(x,\ t).}

\begin{verbatim}
Головна складність цього методу — розклад густини джерела на компоненти f_n(t)\,X_n(x) . Виявляється, що X_n(x) факторами в цій задачі є власними векторами системи Штурма–Ліувіля, які виникають при розв'язанні відповідної однорідної задачі методом розділення змінних, а саме
\end{verbatim}

\texttt{\textbackslash{}begin\{array\}\{c\}\ u\_t\ =\ \textbackslash{}alpha\^{}2\ u\_\{xx\}\ \textbackslash{}quad\ \textbackslash{}text\{(отметим,\ что\ \}\ f(x,\ \textbackslash{}\ t)\ =\ 0),\ \textbackslash{}\textbackslash{}\ u(0,\ \textbackslash{}\ t)\ =\ 0,\ \textbackslash{}\textbackslash{}\ u(1,\ \textbackslash{}\ t)\ =\ 0,\ \textbackslash{}\textbackslash{}\ u(x,\ \textbackslash{}\ 0)\ =\ \textbackslash{}varphi(x).\ \textbackslash{}end\{array\}}

\begin{verbatim}
Завдання Штурма-Ліувіля в цьому випадку виглядає так:
\end{verbatim}

\texttt{X\textquotesingle{}\textquotesingle{}\ +\ \textbackslash{}lambda\^{}2\ X\ =\ 0,}

\begin{verbatim}
X(0) = 0, X(1) = 0, а його розв'язки є функціями
\end{verbatim}

\texttt{X\_n(x)\ =\ \textbackslash{}sin(n\textbackslash{}pi\ x),\ \textbackslash{}quad\ n\ =\ 1,\ 2,\ 3\ \textbackslash{}dots}

\begin{verbatim}
Отже, розклад на густину джерела представлений як ряд
\end{verbatim}

\texttt{(9.5)\ \textbackslash{}\ f(x,t)\ =\ f\_1(t)\ \textbackslash{}sin(\textbackslash{}pi\ x)\ +\ f\_2(t)\ \textbackslash{}sin(2\textbackslash{}pi\ x)\ +\ \textbackslash{}dots\ +\ f\_n(t)\ \textbackslash{}sin(n\textbackslash{}pi\ x)\ +\ \textbackslash{}dots}

\begin{verbatim}
\mathcal{U} нарешті, щоб знайти функції f_n(t) , помножте обидві сторони останнього відношення на \sin(m\pi x) і інтегруйте від 0 до 1 за координатою x.

У результаті ми отримуємо
\end{verbatim}

\texttt{\textbackslash{}int\_\{0\}\^{}\{1\}\ f(x,\ t)\ \textbackslash{}sin(m\textbackslash{}pi\ x)\ dx\ =\ \textbackslash{}sum\_\{n=1\}\^{}\{\textbackslash{}infty\}\ f\_\{n\}(t)\ \textbackslash{}int\_\{0\}\^{}\{1\}\ \textbackslash{}sin(m\textbackslash{}pi\ x)\ \textbackslash{}sin(n\textbackslash{}pi\ x)\ dx\ =}
\texttt{=\ \textbackslash{}frac\{1\}\{2\}\ f\_\{m\}(t)}

\begin{verbatim}
або (після заміни t на p)

(9.6)
\end{verbatim}

\texttt{f\_\{n\}(t)\ =\ 2\ \textbackslash{}int\_\{0\}\^{}\{1\}\ f(x,\ t)\ \textbackslash{}sin(n\textbackslash{}pi\ x)\ dx.}

\begin{verbatim}
Останнє співвідношення встановлює співвідношення між коефіцієнтами f_n(t) і часткою джерела тепла f(x, t).

КРОК 2. (Пошук u_n(x, t) = X_n(x) T_n(t) реакції на вхідний вплив f_n(t) X_n(x) .)

Замінимо густину джерела на його розклад у рядок
\end{verbatim}

\texttt{f(x,\ t)\ =\ \textbackslash{}sum\_\{n=1\}\^{}\{\textbackslash{}infty\}\ f\_n(t)\ \textbackslash{}sin(n\textbackslash{}pi\ x)}

\begin{verbatim}
~~
і намагатися знайти індивідуальні відповіді, тобто функції, T_{\mathbf{n}}(t) у
~~
\end{verbatim}

\texttt{u(x,\ t)\ =\ \textbackslash{}sum\_\{n=1\}\^{}\{\textbackslash{}infty\}\ T\_n(t)\ \textbackslash{}sin(n\textbackslash{}pi\ x).}

\begin{verbatim}
Оскільки розв'язок задачі має вигляд
\end{verbatim}

\texttt{u(x,\ t)\ =\ \textbackslash{}sum\_\{n=1\}\^{}\{\textbackslash{}infty\}\ T\_n(t)\ \textbackslash{}sin(n\textbackslash{}pi\ x),}

\begin{verbatim}
потім підставимо цей вираз у початкову задачу
\end{verbatim}

\texttt{u\_\{t\}\ =\ \textbackslash{}alpha\^{}\{2\}\ u\_\{xx\}\ +\ \textbackslash{}sum\_\{n=1\}\^{}\{\textbackslash{}infty\}\ f\_\{n\}(t)\ \textbackslash{}sin(n\textbackslash{}pi\ x),}
u(0, t) = 0, u(1, t) = 0,
\texttt{u(x,\ 0)\ =\ \textbackslash{}varphi(x),}

\begin{verbatim}
У результаті ми отримаємо
\end{verbatim}

\texttt{\textbackslash{}sum\_\{n=1\}\^{}\{\textbackslash{}infty\}\ T\textquotesingle{}\_n(t)\ \textbackslash{}sin(n\textbackslash{}pi\ x)\ =\ -\textbackslash{}alpha\^{}2\ \textbackslash{}sum\_\{n=1\}\^{}\{\textbackslash{}infty\}\ (n\textbackslash{}pi)\^{}2\ T\_n(t)\ \textbackslash{}sin(n\textbackslash{}pi\ x)\ +\ +\ \textbackslash{}sum\_\{n=1\}\^{}\{\textbackslash{}infty\}\ f\_n(t)\ \textbackslash{}sin(\textbackslash{}pi\ n\textbackslash{}lambda),}

\begin{verbatim}
(9.7)
\end{verbatim}

\texttt{\textbackslash{}sum\_\{n=1\}\^{}\{\textbackslash{}infty\}\ T\_n(t)\ \textbackslash{}sin\ 0\ =\ 0}

\begin{verbatim}
(задовольняється ідентично), \sum_{n=1}^{\infty} T_n(t) \sin (n\pi) = 0 (задовольняється ідентично), \sum_{n=1}^{\infty} T_n(0) \sin (n\pi x) = \phi(x) .

Переписування рівняння та початкової умови

(PPP)
\end{verbatim}

\texttt{\textbackslash{}sum\_\{n=1\}\^{}\{\textbackslash{}infty\}\ {[}T\textquotesingle{}\_n\ +\ (n\textbackslash{}pi\textbackslash{}alpha)\^{}2\ T\_n\ -\ f\_n(t){]}\ \textbackslash{}sin(n\textbackslash{}pi\ x)\ =\ 0,}

\begin{verbatim}
(НУ)
\end{verbatim}

\texttt{\textbackslash{}sum\_\{n=1\}\^{}\{\textbackslash{}infty\}\ T\_n(0)\ \textbackslash{}sin(n\textbackslash{}pi\ x)\ =\ \textbackslash{}varphi(x),}

\begin{verbatim}
можна побачити, що функції T_n(\ell) є розв'язками задачі Коші

(9.8)
\end{verbatim}

\texttt{T\_\{n\}\ +\ (n\textbackslash{}pi\textbackslash{}alpha)\^{}\{2\}\ T\_\{n\}\ =\ f\_\{n\}(t),}
\texttt{T\_\{n\}(0)\ =\ 2\ \textbackslash{}int\_\{0\}\^{}\{1\}\ \textbackslash{}varphi(\textbackslash{}xi)\ \textbackslash{}sin(n\textbackslash{}pi\textbackslash{}xi)\ d\textbackslash{}xi\ =\ a\_\{n\}.}

\begin{verbatim}
Усі ці TAC легко розв'язати, а їхні рішення записуються у вигляді

(99)
\end{verbatim}

\texttt{T\_n(t)\ =\ a\_n\ e\^{}\{-(n\textbackslash{}pi\textbackslash{}alpha)\^{}2\ t\}\ +\ \textbackslash{}int\_0\^{}t\ e\^{}\{-(n\textbackslash{}pi\textbackslash{}alpha)\^{}2\ (t-\textbackslash{}tau)\}\ f\_n(\textbackslash{}tau)\ d\textbackslash{}tau,}

\begin{verbatim}
З того місця, де легко отримати повне розв'язання проблеми (9.3):

(9.10)
\end{verbatim}

\texttt{u(x,\ t)\ =\ \textbackslash{}sum\_\{n=1\}\^{}\{\textbackslash{}infty\}\ T\_n(t)\ \textbackslash{}sin(n\textbackslash{}pi\ x)\ =}
\texttt{=\ \textbackslash{}sum\_\{n=1\}\^{}\{\textbackslash{}infty\}\ \textbackslash{}left{[}\ a\_n\ e\^{}\{-(n\textbackslash{}pi\textbackslash{}alpha)\^{}n\ t\}\ \textbackslash{}sin(n\textbackslash{}pi\ x)\ \textbackslash{}right{]}\ +\ \textbackslash{}sum\_\{n=1\}\^{}\{\textbackslash{}infty\}\ \textbackslash{}left{[}\ \textbackslash{}sin(n\textbackslash{}pi\ x)\ \textbackslash{}int\_0\^{}t\ e\^{}\{-(n\textbackslash{}pi\textbackslash{}alpha)\^{}n\ (t-\textbackslash{}tau)\}\ f\_n(\textbackslash{}tau)\ d\textbackslash{}tau\ \textbackslash{}right{]}}

\begin{verbatim}
Перехідна частина (обумовлена стаціонарною частиною (умовлена початковою умовою) правою частиною f(x, t)).

Розв'язок (9.10) показує, що температурна реакція складається з двох частин: перша частина зумовлена початковими умовами, а друга — густиною джерела f(x,t). Зверніть увагу, що термін «стаціонарна частина» не є найкращим описом другої частини температурної реакції, оскільки він не означає, що ця частина не залежить від часу (стаціонарний стан може бути періодичним, якщо f(x, t) є періодичною функцією часу t).

Проблема повністю вирішена. Перш ніж перейти до наступної теми, давайте покажемо, як використовувати описаний вище метод для розв'язання

Конкретне завдання.

70

# Розв'язання конкретної задачі методом декомпозиції за власними функціями

Розглянемо просту задачу

(9.11)
\end{verbatim}

\texttt{\textbackslash{}begin\{aligned\}\ u\_t\ \&=\ \textbackslash{}alpha\^{}2\ u\_\{xx\}\ +\ \textbackslash{}sin\{(3\textbackslash{}pi\ x)\},\ \&\ 0\ \textless{}\ x\ \textless{}\ 1,\ \textbackslash{}\textbackslash{}\ (\textbackslash{}mathbf\{FY\})\ \&\ \textbackslash{}begin\{cases\}\ u(0,\ t)\ \&=\ 0,\ \textbackslash{}\textbackslash{}\ u(1,\ t)\ \&=\ 0,\ \textbackslash{}end\{cases\}\ \&\ 0\ \textless{}\ t\ \textless{}\ \textbackslash{}infty,\ \textbackslash{}\textbackslash{}\ (\textbackslash{}mathbf\{HY\})\ \&\ u(x,\ 0)\ \&=\ \textbackslash{}sin\{(\textbackslash{}pi\ x)\},\ \&\ 0\ \textbackslash{}leq\ x\ \textbackslash{}leq\ 1.\ \textbackslash{}end\{aligned\}}

\begin{verbatim}
Наша мета — обчислити коефіцієнти T_n(t) при розкладі
\end{verbatim}

\texttt{u\textbackslash{}left(x,\textbackslash{},t\textbackslash{}right)=\textbackslash{}sum\_\{n=1\}\^{}\{\textbackslash{}infty\}\textbackslash{},T\_\{n\}\textbackslash{}left(t\textbackslash{}right)\textbackslash{}sin\textbackslash{}left(n\textbackslash{}pi\ x\textbackslash{}right)}

\begin{verbatim}
(Власні функції цієї задачі X_n(x) збігатися з власними функціями попереднього). Якщо підставити цю декомпозицію на початкову задачу, отримаємо наступну задачу Коші для визначення функцій T_n(t) :
\end{verbatim}

\texttt{T\_n\ +\ (n\textbackslash{}pi\textbackslash{}alpha)\^{}2\ T\_n\ =\ \textbackslash{}int\_n\ (t)\ =\ 2\ \textbackslash{}int\_0\^{}1\ \textbackslash{}sin(3\textbackslash{}pi\ x)\ \textbackslash{}sin(n\textbackslash{}pi\ x)\ dx\ =}
\texttt{\textbackslash{}begin\{cases\}\ 1\ \&n\ =\ 3,\ \textbackslash{}\textbackslash{}\ 0\ \&\ n\ \textbackslash{}neq\ 3,\ \textbackslash{}end\{cases\}}
\texttt{T\_n(0)\ =\ 2\ \textbackslash{}int\_0\^{}1\ \textbackslash{}sin(n\textbackslash{}xi)\ \textbackslash{}sin(n\textbackslash{}pi\textbackslash{}xi)\ d\textbackslash{}xi\ =}
\texttt{\textbackslash{}begin\{cases\}\ 1\ \&n\ =\ 1,\ \textbackslash{}\textbackslash{}\ 0\ \&\ n\ \textbackslash{}neq\ 1.\ \textbackslash{}end\{cases\}}

\begin{verbatim}
Перепишем эти уравнения последовательно для n=1, 2...

~~
\begin{aligned}
(n=1)\quad & T_1' + (\pi\alpha)^2 T_1 = 0, && T_1(0) = 1, && T_1(t) = e^{-(\pi\alpha)^2 t}, \\
(n=2)\quad & T_2' + (2\pi\alpha)^2 T_2 = 0, && T_2(0) = 0, && T_2(t) = 0, \\
(n=3)\quad & T_3' + (3\pi\alpha)^2 T_3 = 4, && T_3(0) = 0, && T_3(t) = \frac{1}{(3\pi\alpha)^2}\left[1 - e^{-(3\pi\alpha)^2 t}\right], \\
(n\geqslant 4)\quad & T_n' + (n\pi\alpha)^2 T_n = 0, && T_n(0) = 0, && T_n(t) = 0.
\end{aligned}
~~

Отже, розв'язок задачі можна записати у вигляді
\end{verbatim}

\texttt{u(x,\ t)\ =\ e\^{}\{-(\textbackslash{}pi\textbackslash{}alpha)\^{}2\ t\}\ \textbackslash{}sin(\textbackslash{}pi\ x)\ +\ \textbackslash{}frac\{1\}\{(3\textbackslash{}pi\textbackslash{}alpha)\^{}2\}\ {[}1\ -\ e\^{}\{-(3\textbackslash{}pi\textbackslash{}alpha)\^{}2\ t\}{]}\ \textbackslash{}sin(3\textbackslash{}pi\ x).}

\begin{verbatim}
(9.12)

# ЗАУВАЖЕННЯ

1. Метод розкладу власних функцій є одним із найпотужніших методів розв'язання неоднорідних диференціальних рівнянь у похідних похідних. Пізніше, коли ми ознайомимося з інтегральними перетвореннями, побачимо, що існують й інші методи розв'язання задач такого типу.

2. Власні функції X_n(x) , за допомогою якої виконується розклад, залежать від задачі, яку потрібно розв'язати, точніше від рівнянь і граничних умов. Читача просять розв'язати задачу 4 наприкінці лекції, щоб перевірити, чи дійсно він навчився знаходити власні функції X_n(x) .

3. Якщо читач пам'ятає основні факти з теорії звичайних диференціальних рівнянь, то він знає, що розв'язки, що відповідають неоднорідним членам вигляду
\end{verbatim}

\texttt{P\_n(x)\ e\^{}\{\textbackslash{}alpha\ x\}}
\texttt{\textbackslash{}begin\{cases\}\ \textbackslash{}sin\{(\textbackslash{}beta\ x)\},\ \textbackslash{}\textbackslash{}\ \textbackslash{}cos\{(\textbackslash{}beta\ x)\},\ \textbackslash{}end\{cases\}}

\begin{verbatim}
можно найти методом неопределенных коэффициентов. То же самое и здесь. Задача (9.11) может быть решена таким же способом.

# ЗАДАЧИ

1. Подставьте разложение (46.8) в задачу (46.7) и получите последовательность задач P_0,\ P_1,\ P_2,\ \dots.

2. Покажите, что нелинейную задачу
\end{verbatim}

\texttt{\textbackslash{}begin\{aligned\}}
\texttt{\textbackslash{}Delta\ u\ +\ u\^{}2\ \&=\ 0,\ \&\&\ 0\ \textbackslash{}leqslant\ r\ \textless{}\ 1,\textbackslash{}\ 0\ \textbackslash{}leqslant\ \textbackslash{}theta\ \textless{}\ 2\textbackslash{}pi,\ \textbackslash{}\textbackslash{}}
\texttt{u(1,\textbackslash{}theta)\ \&=\ \textbackslash{}cos\textbackslash{}theta,\ \&\&\ 0\ \textbackslash{}leqslant\ \textbackslash{}theta\ \textless{}\ 2\textbackslash{}pi,}
\texttt{\textbackslash{}end\{aligned\}}

\begin{verbatim}
можно свести к последовательности линейных задач P_0,\ P_1,\ P_2.

3. Подставьте приближенное решение (46.5) в нелинейную задачу 2 и оцените его точность.

4. Решите задачу P_1 из примера с возмущенной границей и проверьте, насколько хорошо приближенное решение
\end{verbatim}

\texttt{u(r,\textbackslash{}theta)\ =\ u\_0(r,\textbackslash{}theta)\ +\ \textbackslash{}frac\{1\}\{4\}u\_1(r,\textbackslash{}theta)}

\begin{verbatim}
удовлетворяет соотношениям
\end{verbatim}

\texttt{\textbackslash{}Delta\ u\ =\ 0,}
\texttt{u\textbackslash{}left(1\ +\ \textbackslash{}frac\{1\}\{4\}\textbackslash{}sin\textbackslash{}theta,\textbackslash{}\ \textbackslash{}theta\textbackslash{}right)\ =\ \textbackslash{}cos\textbackslash{}theta.}

\begin{verbatim}
# Лекція 47

# РОЗВ'ЯЗАННЯ ДИФЕРЕНЦІАЛЬНИХ РІВНЯНЬ У ПОХІДНИХ ПОХІДНИХ ЗА ДОПОМОГОЮ МЕТОДУ КОНФОРМНИХ ВІДОБРАЖЕНЬ

МЕТА ЛЕКЦІЇ: Показати, як деякі двовимірні граничні задачі можна перетворити на інші, простіші задачі за допомогою конформного відображення. Припустимо, наприклад, що ми хочемо розв'язати рівняння Лапласа
\end{verbatim}

\texttt{\textbackslash{}begin\{cases\}\ u(0,\ t)\ =\ 0,\ \textbackslash{}\textbackslash{}\ u\_x(1,\ t)\ =\ -u(1,\ t),\ \textbackslash{}\textbackslash{}\ u(x,\ 0)\ =\ x,\ \textbackslash{}\textbackslash{}\ u\_t(x,\ 0)\ =\ 0.\ \textbackslash{}end\{cases\}}

\begin{verbatim}
з певною граничною умовою в області комплексної форми (площина змінних x і y). Цю задачу можна перетворити на нову, в якій потрібно знайти розв'язок рівняння Лапласа

 у простішій області (площина змінних і і і).

![](images/_page_353_Figure_9.jpeg)

У конформному відображенні рівняння Лапласа \varphi_{xx}+\varphi_{yy}=0 у координатній площині (x,y) повертається у рівняння Лапласа \varphi_{nn}\dashv -\varphi_{vo}=0 у координатній площині (u,v). Інші рівняння, такі як \varphi_{xx}+2\varphi_{xy}+\varphi_{yy}=0 , не мають такої властивості. Після розв'язання рівняння Лапласа \varphi(u,v) у простій області (

наприклад, у колі, напівплощині, квадраті) знайдено, достатньо підставити вирази U=U(X, Y), V=V(X, Y) у цей розв'язок, і ми отримаємо розв'язок нашої задачі \phi[u(x, y), v(x, y)] , вираженого через початкові змінні.

У цій лекції ми розглянемо:

1) про властивості конформних відображень,

2) як їх будувати (для деяких районів),

3) як розв'язувати задачі методом конформних відображень.

Однією з головних труднощів, що виникають при розв'язанні задач меж, є складна форма кордонів. Навіть межі відносно простої форми часто роблять задачу дуже складною для розв'язання. Звісно, можна спробувати розв'язати задачу за допомогою методів теорії збурень, але теорія збурень працює добре лише у випадках, коли межа близька до простої.

Однак існує спосіб розв'язати двовимірне рівняння Лапласа в областях з комплексною межею. Цей метод базується на конформних відображеннях.

Перш ніж почати розв'язувати задачі цим методом, давайте трохи ознайомимося з поняттям конформного відображення та комплексних змінних загалом.

# Комплексні функції та конформні відображення

У цій лекції ми ознайомимося лише з тими поняттями з теорії функцій комплексної змінної, які нам знадобляться пізніше. Представимо комплексне число z = x + iy як точку на площині комплексного xy (рис. 47.1).

Щоб ввести поняття конформного відображення, спершу ознайомимося з поняттям функції комплексної змінної

.

Тут z — це комплексна змінна (з деякої області площини z), а w — нова комплексна змінна, яка отримає значення відповідно до формули w = f(z) у новій комплексній площині w.

Наприклад, комплексна функція, визначена в *першому квадранті* комплексної z-площини, w=z^2 відображає цю область у всю півплощину v>0 у комплексній площині w.

Таблиця 47 2

~~
| Площина d | z \longrightarrow z^2 | Площина S |
~~

|--------------------------|-------------------------|------------------------------------------------------------|

| 0 | | 0 |

| <i>t</i><br>1 → <i>t</i> | | -1<br>(2) |

| <br>Вісь x Перший квадрант | | Позитивна дійсна піввісь верхньої півплощини |

Таблиця 47.2 показує відповідність, яка встановлюється шляхом відображення w=z^2 між деякими точками на площинах z і w.

![](images/_page_355_Figure_5.jpeg)

~~
Рис. 47.1. Площина комплексної змінної \boldsymbol{z} та некольорових формул бігунів Р И С У Н
~~

![](images/_page_355_Figure_7.jpeg)

~~
ОК 47.2. Відображення першого квадранта площини {\bf z} у верхню півплощину {\bf w} .
~~

Інакше кажучи, якщо ми говоримо про комплексісні функції змінної w = f(z), то можна розглянути питання, як криві площини z перетворюються на криві в площині

Це питання не ставиться в теорії функцій дійсної змінної, з якою читач уже добре знайомий.

Перш ніж рухатися далі, давайте розглянемо, як w=z^2 складне відображення типів у еквівалентній дійсній формі. Дійсна форма показує, як початкові координати (x,y) у площині z перетворюються на нові координати (u,v) у площині w w=z^2 .
\end{verbatim}

\texttt{\textbackslash{}begin\{cases\}\ u(x,\ 0)\ =\ x,\ \textbackslash{}\textbackslash{}\ u\_t(x,\ 0)\ =\ 0.\ \textbackslash{}end\{cases\}}

\begin{verbatim}
І давайте окремо прирівняємо реальні та уявні частини в цьому співвідношенні. У результаті ми отримуємо
\end{verbatim}

\texttt{\textbackslash{}begin\{array\}\{ll\}\ u=x\^{}2-y\^{}2,\textbackslash{}\textbackslash{}\ v=2xy,\ \textbackslash{}end\{array\}}

\begin{verbatim}
(действительная форма отображения  w=z^2 ).

Эта форма понадобится нам позже. В этом примере она интересна тем, что показывает, как гиперболы в плоскости z переходят в координатные линии u = const и v = const в плоскости w.

Теперь мы перейдем к изучению одного частного вида функций комплексного переменного, которые осуществляют конформные отображения.

# Определение конформного отображения

~~
Отображение w=f(z) комплексиой плоскости z на комплексную плоскость w называется конформным в точке  z_{\bullet}  плоскости z, если производная  f'(z_{\bullet}) \neq 0 . Отображение f(z) называется конформным в области D, если  f'(z) \neq 0  в каждой точке области D.
~~

~~
Например, отображение  f(z) = z^2  конформно всюду, за исключением точки z = 0, поскольку  z'(z) = 2z \neq 0  для всех  z \neq 0 . C другой стороны, отображение  e^z  конформно во всей плоскости z, так как всюду  f'(z) = e^z \neq 0 . Какая же польза от конформных отображений? Ответ состоит в том, что, решая уравнение Лапласа  \phi_{xx} + \phi_{yy} = 0  внутри какой-то области плоскости переменных x и y, мы можем рассматривать эту плоскость как плоскость комплексного переменного z. Рассмотрим комплексное отображение w = f(z) плоскости z на плоскость w, в которой введены координаты (u, v). При отображении w = f(z) уравнение Лапласа  \varphi_{xx} + \varphi_{yy} = 0 преобразуется в некоторое новое уравнение с частными производными, зависящее от новых координат и и v. Замечательным является следующее: если отображение w = f(z) является конформным в области, где задано уравнение  \phi_{xy} + \phi_{uy} = 0 , то новое уравнение также будет уравнением Лапласа в координатах и и v. To есть уравнение  \phi_{xx} + \phi_{uu} = 0  перейдет снова в уравнение  \varphi_{nn} + \varphi_{nn} = 0 . Так мы приходим к идее найти такое конформное отображение, которое переводит область со сложной границей в область с простой границей (напомним, что отображение  w=z^*  переводит границу первого квадранта плоскости z в действительную ось плоскости w).
~~

Приведем несколько примеров, в которых будет показано, как можно решать уравнения с частными производными с помощью конформных отображений.

Уравнение Лапласа в верхней полуплоскости.

Предположим, что нам необходимо решить следующую задачу Дирихле в верхней полувлоскости (рис. 47.3):

(47.1)

можно найти стандартным методом разделения переменных. При этом оказывается, что решение  u\left(x,\,t\right) , представленное в виде ряда

Один із способів розв'язання цієї проблеми — застосувати перетворення Фур'є до змінної x, але набагато краще конформно відобразити верхню половину площини у площину (u, v). Читачеві не буде складно побачити, що конформне відображення

является суперпозицией простейших колебаний. Пространственная часть  X_n(x)  каждого элементарного колебания является решением (собственной функцией) некоторой краевой задачи Штурма—Лиувилля.

Итак, мы научились решать волиовое уравнение  u_{II}=c^2u_{xx}  на всей прямой  -\infty < x < \infty . Формула Даламбера представляет решение в виде суммы двух бегущих волн, распространяющихся в противоположные стороны. Если же рассматривать это же уравнение на ограниченном промежутке 0 < x < L, то бегущих волн уже не станет, так как они будут взаимодействовать с границами. Вместо них возникают другие волны, которые называют стоячими волнами. Посмотрим, например, что случится, если гитарную струну, закрепленную в точках x=0, x=L, привести в движение. Для ответа на этот вопрос нам необходимо решить следующую смешанную задачу гиперболического типа:

Вона має такі властивості (див. RNS. 47.4).

~~
1. Вона відображає верхню півплощину площини z у смуги -\infty < u < \infty , 0 < v < \pi , у площині w.
~~

![](images/_page_357_Figure_11.jpeg)

РИС 47.3 Задача граничного значення у верхній половині площини.

~~
- 2. Ребро прямої y=0, -1 < x < 1, у площині z (де потенціал дорівнює 1), переходить у пряму v=\pi, - \infty < u < \infty , у площині w. - 3. Два промені x>1, y=0, і x<-1, y=0, у площині z проходять відповідно у позитивний і негативний
~~

Читач повинен бути впевнений, що може підтвердити ці твердження.

![](images/_page_358_Figure_2.jpeg)

Рис. 47.4. Перетворення складної задачі на легку за допомогою конформного відображення.

Важливість цих властивостей полягає в тому, що вони перетворюють початкову вадаху (47.1) на дуже просту задачу

(47.2)

 Розв'язок очевидний: \varphi(u, v) = \frac{1}{\pi}v .

Тепер, щоб знайти розв'язок початкової задачі, потрібно пояснити v через x і y і підставити цей вираз на q(u,v)=v'\pi . Після цього ми отримуємо

Оказывается, в этом случае бегущие волны отражаются от границ таким образом, что результирующие колебания становятся не бегущими, а сохраняющими форму в одном положении, т. е.

![](images/_page_147_Picture_5.jpeg)

РИС. 20.1. Три типичных стоячих волны X(x) T(t).

превращаются в стоячие волны Несколько примеров стоячих волн показано на рис. 20.1. Если известны профили этих стоячих волн   X_n(x)   и характер колебаний каждой волны   T_n(t)  , то решение задачи о колебаниях гитарной струны представляет собой суперпозицию простейших колебаний   X_n(x)\,T_n(t)   и запи-

сывается в виде
\end{verbatim}

\texttt{=\ \textbackslash{}ln\textbackslash{}left{[}\textbackslash{}frac\{z-1\}\{z+1\}\textbackslash{}right{]}\ +\ i\ \textbackslash{}arg\textbackslash{}left{[}\textbackslash{}frac\{z-1\}\{z+1\}\textbackslash{}right{]}}

\begin{verbatim}
~~
Единственное, что остается сделать нам,— это подобрать коэффициенты   c_n   так, чтобы при t=0 удовлетворялись начальные условия   u\left(x,\ 0\right)=f\left(x\right)   и   u_t\left(x,\ 0\right)=g\left(x\right)  .
~~

Задачу о колебаниях гитарной струны мы решим методом разделения переменных.

# Решение задачи о колебаниях ограниченной струны методом разделения переменных

Чтобы решить смешанную задачу

(УЧП)
\end{verbatim}

\texttt{v\ =\ \textbackslash{}arg\textbackslash{}left{[}\textbackslash{}frac\{z-1\}\{z+1\}\textbackslash{}right{]}\ =\ \textbackslash{}arg\textbackslash{}left{[}\textbackslash{}frac\{x+iy-1\}\{x+1+iy\}\textbackslash{}right{]}\ =}
\texttt{=\ \textbackslash{}arg\textbackslash{}left{[}\textbackslash{}frac\{2y\}\{x\^{}2+y\^{}2-1\}\textbackslash{}right{]}}
\texttt{\textbackslash{}varphi(x,\ y)\ =\ \textbackslash{}frac\{1\}\{\textbackslash{}pi\}\ \textbackslash{}arctan\textbackslash{}left{[}\textbackslash{}frac\{2y\}\{x\^{}2\ +\ y\^{}2\ -\ 1\}\textbackslash{}right{]}.}
\texttt{\textbackslash{}begin\{cases\}\ u(0,\ t)\ =\ 0,\ \textbackslash{}\textbackslash{}\ u(L,\ t)\ =\ 0,\ \textbackslash{}end\{cases\}}
x\^{}\{2\} + y\^{}\{2\} = 1,

\begin{verbatim}
якщо потенціал на внутрішньому колі дорівнює одинці, а на зовнішньому колі дорівнює двом (рис. 47.5). Іншими словами, нам потрібно розв'язати проблему

(47.3)

 Проблема полягає в тому, щоб знайти конформне відображення, яке перетворює цю область на щось простіше, де нашу проблему легко розв'язати. У цьому випадку очевидно, що ця проста ділянка має бути кільцем. Відповідна експозиція Р И С У Н

![](images/_page_359_Figure_10.jpeg)

ОК 47.5. Конформне відображення на кільце.

Це вже не так очевидно, але існують довідники, де можна знайти дисплей із вказаними властивостями. Одним із найзагальних є «Довідник конформних відображень»

G. Kober 1), яка містить сотні конформних відображень між різними регіонами. Якщо читач витратить трохи часу і краще ознайомиться з конформними таблицями відображення, він знайде відповідне конформне відображення для задачі (47.3). У нашому конкретному випадку це буде образ виду

(47.4)

де s=-0,146, а t=-6,85. У своїй еквівалентній дійсній формі це буде

(47.5)
\end{verbatim}

\texttt{\textbackslash{}begin\{cases\}\ u(x,\ 0)\ =\ f(x),\ \textbackslash{}\textbackslash{}\ u\_t(x,\ 0)\ =\ g(x),\ \textbackslash{}end\{cases\}}
\texttt{v\ =\ \textbackslash{}gamma\ {[}y(x-t){]},}
\texttt{\textbackslash{}gamma\ =\ 2t/{[}(x-t)\^{}2\ +\ y\^{}2{]}.}

\begin{verbatim}
найдем сначала стоячне волны, т. е. решения вида

~~
\begin{aligned} \varphi_{uu} + \varphi_{vv} &= 0, & \text{в кольче,} \\ (\Gamma Y) & & & & & & & \\ (\Gamma Y) & & & & & & \\ (\Gamma Y) & & & & & \\ & & & & & \\ (\Gamma Y) & & & & & \\ & & & & \\ & & & & \\ & & & & \\ & & & & \\ & & & \\ & & & \\ & & & \\ & & & \\ & & & \\ & & & \\ & & \\ & & \\ & & \\ & & \\ & & \\ & & \\ & & \\ & & \\ & & \\ & & \\ & & \\ & & \\ & & \\ & & \\ & & \\ & & \\ & & \\ & & \\ & & \\ & & \\ & & \\ & & \\ & & \\ & & \\ & & \\ & & \\ & & \\ & & \\ & & \\ & & \\ & & \\ & & \\ & & \\ & & \\ & & \\ & & \\ & & \\ & & \\ & & \\ & & \\ & & \\ & & \\ & & \\ & & \\ & & \\ & & \\ & & \\ & & \\ & & \\ & & \\ & & \\ & & \\ & & \\ & & \\ & & \\ & & \\ & & \\ & & \\ & & \\ & & \\ & & \\ & & \\ & & \\ & & \\ & & \\ & & \\ & & \\ & & \\ & & \\ & & \\ & & \\ & & \\ & & \\ & & \\ & & \\ & & \\ & & \\ & & \\ & & \\ & & \\ & & \\ & & \\ & & \\ & & \\ & & \\ & & \\ & & \\ & & \\ & & \\ & & \\ & & \\ & & \\ & & \\ & & \\ & & \\ & & \\ & & \\ & & \\ & & \\ & & \\ & & \\ & & \\ & & \\ & & \\ & & \\ & & \\ & & \\ & & \\ & & \\ & & \\ & & \\ & & \\ & & \\ & & \\ & & \\ & & \\ & & \\ & & \\ & & \\ & & \\ & & \\ & & \\ & & \\ & & \\ & & \\ & & \\ & & \\ & & \\ & & \\ & & \\ & & \\ & & \\ & & \\ & & \\ & & \\ & & \\ & & \\ & & \\ & & \\ & & \\ & & \\ & & \\ & & \\ & & \\ & & \\ & & \\ & & \\ & & \\ & & \\ & & \\ & & \\ & & \\ & & \\ & & \\ & & \\ & & \\ & & \\ & & \\ & & \\ & & \\ & & \\ & & \\ & & \\ & & \\ & & \\ & & \\ & & \\ & & \\ & & \\ & & \\ & & \\ & & \\ & & \\ & & \\ & & \\ & & \\ & & \\ & & \\ & & \\ & & \\ & & \\ & & \\ & & \\ & & \\ & & \\ & & \\ & & \\ & & \\ & & \\ & & \\ & & \\ & & \\ & & \\ & & \\ & & \\ & & \\ & & \\ & & \\ & & \\ & & \\ & & \\ & & \\ & & \\ & & \\ & & \\ & & \\ & & \\ & & \\ & & \\ & & \\ & & \\ & & \\ & & \\ & & \\ & & \\ & & \\ & & \\ & & \\ & & \\ & & \\ & & \\ & & \\ & & \\ & & \\ & & \\ & & \\ & & \\ & & \\ & & \\ & & \\ & & \\ & & \\ & & \\ & & \\ & & \\ & & \\ & & \\ & & \\ & & \\ & & \\ & & \\ & & \\ & & \\ & & \\ & & \\ & & \\ & & \\ & & \\ & & \\ & & \\ & & \\ & & \\ & & \\ & & \\ & & \\ & & \\ & & \\ & & \\ & & \\ & & \\ & & \\ & & \\ & & \\ & & \\ & & \\ & & \\ & & \\ & & \\ & & \\ & & \\ & & \\ & & \\ & & \\ & & \\ & & \\ & & \\ & & \\ & & \\ & & \\ & & \\ & & \\ & & \\ & & \\ & & \\ & & \\ & & \\ & & \\ & & \\ & & \\ & & \\ & & \\ & & \\ & & \\ & & \\ & & \\ & & \\ & & \\ & & \\ & & \\ & & \\ & & \\ & & \\ & &
~~

Подставив это выражение в уравнение (20.1) и разделив переменные, получим два обыкновенных дифференциальных уравнения \end{aligned}
\end{verbatim}

\texttt{\textbackslash{}varphi(r)\ =\ a\ \textbackslash{}ln\ r\ +\ b,\ \textbackslash{}quad\ r\ =\ u\^{}2\ +\ v\^{}2.}
\texttt{\textbackslash{}varphi(u,\ v)\ =\ 0.57\ \textbackslash{}ln(u\^{}2\ +\ v\^{}2)\ +\ 1}

\begin{verbatim}
где константа  \lambda  — пока что любое действительное число —  \infty < \lambda < \infty  .

Результаты исследования решений этих двух уравнений при различных значениях  \lambda  приведены на рис 20.2.

Теперь мы должны *опбросить* те решения, которые либо неограничены при  t \to \infty , либо становятся тождественно равными нулю, если потребовать выполнения граничных условий u(0, t) = u(L, t) = 0.

В качестве упражнения читателю предлагается проверить, что нетривиальные и ограниченные решения можно получить только при отрицательных значениях параметра  \lambda . Таким образом, наша задача— найти такие значения констант A, B, C, D и отрица-

тельной константы  \lambda = -\beta^2 , чтобы выражение

(20.2)
\end{verbatim}

\texttt{\textbackslash{}varphi(x,\ y)\ =\ 0.57\ \textbackslash{}ln(u\^{}2\ +\ v\^{}2)\ +\ 1,}

\begin{verbatim}
удовлетворяло граничным условиям. Если это удастся, то мы будем располагать множеством простейших колебаний струны, каждое из которых удовлетворяет уравнению с частными произ- Р И С

![](images/_page_149_Figure_5.jpeg)

20.2. Построение стоячих воли при раздичных значениях параметра λ.

водными и граничным условиям. Затем останется только найти такую линейную комбинацию этих колебаний, чтобы сумма удовлетворяла начальным условиям при t=0.

Подстановка (20.2) в граничные условия u(0, t) = u(L, t) = 0 дает
\end{verbatim}

\texttt{u(0,\ t)\ =\ X(0)T(t)\ =\ D\textbackslash{}bigl{[}A\textbackslash{}sin(\textbackslash{}alpha\textbackslash{}beta\ t)\ +\ B\textbackslash{}cos(\textbackslash{}alpha\textbackslash{}beta\ t)\textbackslash{}bigr{]}\ =\ 0\ \textbackslash{}Rightarrow\ D\ =\ 0,}

\begin{verbatim}
и
\end{verbatim}

\texttt{u(L,\ t)\ =\ X(L)T(t)\ =\ C\textbackslash{}sin(\textbackslash{}beta\ L)\textbackslash{}bigl{[}A\textbackslash{}sin(\textbackslash{}alpha\textbackslash{}beta\ t)\ +\ B\textbackslash{}cos(\textbackslash{}alpha\textbackslash{}beta\ t)\textbackslash{}bigr{]}\ =\ 0\ \textbackslash{}Rightarrow\ \textbackslash{}sin(\textbackslash{}beta\ L)\ =\ 0.}

\begin{verbatim}
Другими словами, константа разделения \beta должна удовлетворять уравнению \sin(\beta L) = 0, откуда
\end{verbatim}

\texttt{\textbackslash{}beta\_n\ =\ \textbackslash{}frac\{n\textbackslash{}pi\}\{L\},\ \textbackslash{}qquad\ n\ =\ 0,\ 1,\ 2,\ \textbackslash{}ldots}

\begin{verbatim}
Заметим, что если во втором уравнении (20.3) положить C=0, то получится тривиальное решение X(x)T(t) \equiv 0. Следовательно, мы нашли последовательность элементарных колебаний струны:
\end{verbatim}

\texttt{u\_n(x,\ t)\ =\ X\_n(x)T\_n(t)\ =\ \textbackslash{}sin\textbackslash{}!\textbackslash{}left(\textbackslash{}frac\{n\textbackslash{}pi\ x\}\{L\}\textbackslash{}right)\textbackslash{}left{[}a\_n\ \textbackslash{}sin\textbackslash{}!\textbackslash{}left(\textbackslash{}frac\{n\textbackslash{}pi\ \textbackslash{}alpha\ t\}\{L\}\textbackslash{}right)\ +\ b\_n\ \textbackslash{}cos\textbackslash{}!\textbackslash{}left(\textbackslash{}frac\{n\textbackslash{}pi\ \textbackslash{}alpha\ t\}\{L\}\textbackslash{}right)\textbackslash{}right{]}.}

\begin{verbatim}
(20.4) або
\end{verbatim}

\texttt{u\_n(x,\ t)\ =\ R\_n\ \textbackslash{}sin\textbackslash{}!\textbackslash{}left(\textbackslash{}frac\{n\textbackslash{}pi\ x\}\{L\}\textbackslash{}right)\textbackslash{}cos\textbackslash{}!\textbackslash{}left(\textbackslash{}frac\{n\textbackslash{}pi\ \textbackslash{}alpha\ (t\ -\ \textbackslash{}delta\_n)\}\{L\}\textbackslash{}right).}

\begin{verbatim}
Тут a_n, b_n, R_n і \delta_n — довільні сталі. Кожне таке елементарне коливання є стоячою хвилею.

![](images/_page_150_Picture_5.jpeg)

Рис. 20.3. Стоячі хвилі u_n(x, t) = X_n(x)T_n(t).

В нашей задаче уравнение и граничные условия линейны и однородны, поэтому произвольная сумма элементарных стоячих волн также является решением волнового уравнения и удовлетворяет граничным условиям. Мы должны теперь найти такую суперпозицию стоячих волн, чтобы она удовлетворяла также и начальным условиям, тогда получим решение нашей задачи. Подстановка суммы
\end{verbatim}

Зверніть увагу на зображення.

\pandocbounded{\includegraphics[keepaspectratio]{images/_page_362_Figure_1.jpeg}}

\pandocbounded{\includegraphics[keepaspectratio]{images/_page_362_Figure_2.jpeg}}

\begin{enumerate}
\def\labelenumi{\arabic{enumi}.}
\setcounter{enumi}{3}
\tightlist
\item
  Розв'язати змішану задачу Діріхле-Неймана з кутовим порошком за
  допомогою розчину під кутом 45°
\end{enumerate}

\pandocbounded{\includegraphics[keepaspectratio]{images/_page_362_Figure_5.jpeg}}

ПРИМІТКА. Комплексна функція
\texttt{w\ =\ \textbackslash{}ln\ z\ =\ \textbackslash{}ln\ \textbar{}z\textbar{}\ +\ i\ \textbackslash{}arg(z)}
відображає промінь \texttt{\textbackslash{}theta\ =\ c\_1} у площині z у
пряму \texttt{v\ =\ c\_1} , а коло \texttt{r\ =\ c\_2} у пряму
\texttt{u\ =\ \textbackslash{}ln\ c\_2} у площині w (див. рисунки).

\pandocbounded{\includegraphics[keepaspectratio]{images/_page_362_Figure_7.jpeg}}

\section{Джерела ()}\label{ux434ux436ux435ux440ux435ux43bux430}

\section{Головне}\label{ux433ux43eux43bux43eux432ux43dux435}

\begin{itemize}
\tightlist
\item
  Бідадзе А. В. Рівняння математичної фізики.--- Москва: Наука, 1976.
\item
  Володимиров В. С. Рівняння математичної фізики.--- Москва: Наука,
  1971. Кошляков Н. С., Глінер Е. Б., Смірнов М. М. Рівняння в зокрема
  математичній фізиці.--- Москва: Вища школа, 1970.
\item
  \begin{enumerate}
  \def\labelenumi{\arabic{enumi}.}
  \setcounter{enumi}{3}
  \tightlist
  \item
    Самарський, А. А. Вступ до чисельних методів. Москва, Наукове
    видавництво, 1982.
  \end{enumerate}
\item
  \begin{enumerate}
  \def\labelenumi{\arabic{enumi}.}
  \setcounter{enumi}{4}
  \tightlist
  \item
    Смирнов В. І. Курс, виший математіки, т. 2.--- М.: Наука, 1967.
  \end{enumerate}
\item
  \begin{enumerate}
  \def\labelenumi{\arabic{enumi}.}
  \setcounter{enumi}{5}
  \tightlist
  \item
    Тихонов А. Н., Самарський А. А. Рівняння математичної фізики.
    Москва, Наукове видання, 1972.
  \end{enumerate}
\end{itemize}

\section{Додатково}\label{ux434ux43eux434ux430ux442ux43aux43eux432ux43e}

\begin{itemize}
\tightlist
\item
  \begin{enumerate}
  \def\labelenumi{\arabic{enumi}.}
  \setcounter{enumi}{6}
  \tightlist
  \item
    Задача Адамара Дж. Коші для лінійних диференціальних рівнянь у
    похідних похідних гіперболічного типу. Москва, Наукове видання,
    1978.
  \end{enumerate}
\item
  Арсенін В. Я. Методи математичної фізики та спеціальні функції.
  Москва; Наука, 1974.
\item
  \begin{enumerate}
  \def\labelenumi{\arabic{enumi}.}
  \setcounter{enumi}{8}
  \tightlist
  \item
    Арфкен Г. Математичні методи у фізиці.--- Москва: Атомиздат, 1970.
  \end{enumerate}
\item
  \begin{enumerate}
  \def\labelenumi{\arabic{enumi}.}
  \setcounter{enumi}{9}
  \tightlist
  \item
    Берс Л., Джон Ф., Шехтер Б. Диференціальні рівняння з частинними
    похідними. Москва, видавництво Мир, 1966.
  \end{enumerate}
\item
  \begin{enumerate}
  \def\labelenumi{\arabic{enumi}.}
  \setcounter{enumi}{10}
  \tightlist
  \item
    Бінадзе, А. В. Крайові задачі для еліптичних рівнянь другого
    порядку.--- Москва: Наука, 1966.
  \end{enumerate}
\item
  \begin{enumerate}
  \def\labelenumi{\arabic{enumi}.}
  \setcounter{enumi}{11}
  \tightlist
  \item
    Бідадзе А. В. Деякі класи диференціальних рівнянь у похідних
    похідних. Москва, Наукове видання, 1981.
  \end{enumerate}
\item
  \begin{enumerate}
  \def\labelenumi{\arabic{enumi}.}
  \setcounter{enumi}{12}
  \tightlist
  \item
    Біцадзе А. В., Калініченко Д. К. Збірка задач з рівнянь математичної
    фізики. Москва, Наука, 1977.
  \end{enumerate}
\item
  \begin{enumerate}
  \def\labelenumi{\arabic{enumi}.}
  \setcounter{enumi}{13}
  \tightlist
  \item
    Будак Б. М., Самарський А. А., Тихонов А. Н. Сборник задач по
    математичній фізиці {[}Збірка задач математичної фізики{]}. Москва,
    Наукове видання, 1972.
  \end{enumerate}
\item
  Вазов В., Форсайт Дж. Методи різниці для розв'язання диференціальних
  рівнянь у похідних похідних.--- Москва: Іл, 1963
\item
  Нові методи розв'язання еліптичних рівнянь {[}Нові методи розв'язання
  еліптичних рівнянь{]}. Москва, Гостехіддат, 1948. Векуа, І. Н.
  Узагальнені аналітичні функції. Москва, Фізматгиз, 1959.
\item
  Володимиров В. С. Узагальнені функції в математичній фізиці, Москва,
  Наука, 1976.
\item
  \begin{enumerate}
  \def\labelenumi{\arabic{enumi}.}
  \setcounter{enumi}{18}
  \tightlist
  \item
    Володимиров В. С. та ін. Збірка задач з рівнянь математичної фізики.
    --- Москва: Наука, 1974.
  \end{enumerate}
\item
  Годунов С. К. Рівняння математичної фізики.--- Москва: Наука, 1971.
  Годунов С. К., Рябенський В. С. Схеми різниці.--- Москва: Наука, 1977.
\item
  Гордін: задача Л. Коші для гіперболічних рівнянь.--- Москва: IL, 1961.
  Джеффріс Г., Спірлс Б. Методи математичної фізики. т. 1-3, --- Москва:
  Мир,
\item
  \begin{enumerate}
  \def\labelenumi{\arabic{enumi}.}
  \setcounter{enumi}{1969}
  \tightlist
  \item
  \end{enumerate}
\item
  \begin{enumerate}
  \def\labelenumi{\arabic{enumi}.}
  \setcounter{enumi}{23}
  \tightlist
  \item
    Зоммерфельд А. Рівняння в часткових похідних математичної фізики.---
    Москва: Іл, 1950. 25. Кализкін, Н. Н. Численільна методика.--- М.:
    Наука, 1978.
  \end{enumerate}
\item
  25 Курант Р., Гілберт Д. Методи математичної фізики. т. 1-2, - Москва:
  Гостехиздат, 1951.
\item
  \begin{enumerate}
  \def\labelenumi{\arabic{enumi}.}
  \setcounter{enumi}{26}
  \tightlist
  \item
    Курант Р. Рівняння з частковими виробництвами. Москва, Видавництво
    Мир, 1964.
  \end{enumerate}
\item
  \begin{enumerate}
  \def\labelenumi{\arabic{enumi}.}
  \setcounter{enumi}{27}
  \tightlist
  \item
    Лаврентьєв М. А., Шабат Б. В. Методи теорії функцій комплексної
    змінної. - Москва: Наука, 1965.
  \end{enumerate}
\item
  \begin{enumerate}
  \def\labelenumi{\arabic{enumi}.}
  \setcounter{enumi}{28}
  \tightlist
  \item
    Ладиженська О. А. Прикордонні задачі математичної фізики. Москва:
    Наука,
  \end{enumerate}
\item
  Ладиженська О. А. Змішана задача для гіперболічних рівнянь.--- Москва:
  Фізматгиз, 1953.
\end{itemize}

Список був складений редактором перекладу.

\begin{itemize}
\tightlist
\item
  Ладиженська О. А., Солоняков В. А., Уралцева Н. Н. Лінійні та
  кварілінні рівняння параболічного типу.--- Москва: Наука, 1967.
\item
  \begin{enumerate}
  \def\labelenumi{\arabic{enumi}.}
  \setcounter{enumi}{31}
  \tightlist
  \item
    Ладиженська О. А., Уралцева Н. Н. Лінійні та квазілінійні рівняння
    еліптичного типу. - Москва: Наука, 1964, 1973.
  \end{enumerate}
\item
  Лебедєв Н. Н., Скальська М. П., Уфляйнд Ю. С. Сборник задах по
  математичній фізиці.--- Москва: Гостехиддат, 1965.
\item
  \begin{enumerate}
  \def\labelenumi{\arabic{enumi}.}
  \setcounter{enumi}{33}
  \tightlist
  \item
    Лайонс, Дж.--- Л. Деякі методи розв'язання задач з нелінійними
    краями.--- М.З. Мір. 1972.
  \end{enumerate}
\item
  \begin{enumerate}
  \def\labelenumi{\arabic{enumi}.}
  \setcounter{enumi}{34}
  \tightlist
  \item
    Лайонс Дж.---Л., Maggenes E. Neobreed Grancinum Problems та їх
    застосування. Москва, Видавництво Мир, 1971. 36. Марчук Г. І. Методи
    обчислювальної математики. Москва, Наукове видання, 1980.
  \end{enumerate}
\item
  Мізохата С. Теорія рівнянь з приватними виробниками. --- Москва: Мир,
  1977.
\item
  \begin{enumerate}
  \def\labelenumi{\arabic{enumi}.}
  \setcounter{enumi}{37}
  \tightlist
  \item
    Міранда К. Рівняння з частковими похідними еліптичного типу.---
    Москва: Іл, 1957.
  \end{enumerate}
\item
  Михайлов, В. П. Диференціальні рівняння в похідних з похідними. ---
  Москва: Наука, 1976.
\item
  Міклін С. Г. Курс математичної фізики.--- Москва: Наука, 1968. Міхлін,
  С. Г. Лінійні диференціальні рівняння. --- Москва; Вишня школа, 1977.
\item
  Морс П. М., Фецибах X, Методи теоретичної фізики. т. 1-2.--- Москва:
  IL, 1958.
\item
  \begin{enumerate}
  \def\labelenumi{\arabic{enumi}.}
  \setcounter{enumi}{42}
  \tightlist
  \item
    Меттьюз Дж., Вокер Р. Математичні методи фізики.--- Москва:
    Атомиздат, 1972.
  \end{enumerate}
\item
  Лекції Нагумо М. з сучасної теорії рівнянь у приватних виробництвах
  Москва, видавництво Мвр, 1967. 45. Обен, Дж.-П. Приблизне розв'язання
  еліптичних граничних задач.--- Москва:
\item
  Мір, 1977. 46. Петровський І. Г. Лекції з рівнянь із частковими
  похідними. --- Москва: Фізматгиз, 1961.
\item
  \begin{enumerate}
  \def\labelenumi{\arabic{enumi}.}
  \setcounter{enumi}{46}
  \tightlist
  \item
    Положий Г. Н. Еквалізація математичної фізики.--- Москва: Вища
    школа.
  \end{enumerate}
\item
  \begin{enumerate}
  \def\labelenumi{\arabic{enumi}.}
  \setcounter{enumi}{47}
  \tightlist
  \item
    Рід М., Саймон Б. Методи сучасної математичної фізики.--- Москва:
    Мир,
  \end{enumerate}
\item
  т. 1-1977, т. 2-1978, т. 3, 4-1982. Ріхамаєр Р. Принципи сучасної
  математичної фізики.--- Москва: Мир,
\item
  \begin{enumerate}
  \def\labelenumi{\arabic{enumi}.}
  \setcounter{enumi}{49}
  \tightlist
  \item
    Ріхтмаєр Р., Мортон К. Методи різниці для розв'язання задач з
    граничними значеннями.--- Москва:
  \end{enumerate}
\item
  Мір, 1972,
\item
  Рождественський Б. Л., Япенко, Н. Н. Системи квазіполіцейських рівнянь
  та їх застосування в газовій динаміці.--- Москва: Наука, 1978.
\item
  52, Самарський А. А. Теорія різницьких схем. Москва, Наука, 1977.
\item
  \begin{enumerate}
  \def\labelenumi{\arabic{enumi}.}
  \setcounter{enumi}{52}
  \tightlist
  \item
    Самарський А. А., Андрєєв В. Б. Різниці для еліптичних рівнянь---
    Москва, Наука, 1976.
  \end{enumerate}
\item
  \begin{enumerate}
  \def\labelenumi{\arabic{enumi}.}
  \setcounter{enumi}{53}
  \tightlist
  \item
    Смірнов В. І. Курс вищої математики.--- Москва: Наука, 1974.
  \end{enumerate}
\item
  \begin{enumerate}
  \def\labelenumi{\arabic{enumi}.}
  \setcounter{enumi}{54}
  \tightlist
  \item
    Смірнов М. М. Дегенеративні еліптичні та гіперболічні рівняння.--
    Москва: Наука, 1966.
  \end{enumerate}
\item
  56, Смірнов М. М. Диференціальні рівняння у похідних другого порядку,
  -- Мінськ: Видавництво БДУ, 1974.
\item
  \begin{enumerate}
  \def\labelenumi{\arabic{enumi}.}
  \setcounter{enumi}{56}
  \tightlist
  \item
    Смірнов М. М. Проблеми рівнянь математичної фізики. Москва, Наукове
    видавництво, 1965.
  \end{enumerate}
\item
  \begin{enumerate}
  \def\labelenumi{\arabic{enumi}.}
  \setcounter{enumi}{57}
  \tightlist
  \item
    Соболєв С. Л. Пекаки primeneniya funktsional'nogo analiza v
    matematicheskoi fizike {[}Застосування функціонального аналізу в
    математичній фізиці{]}. Ленінград: Видавництво ЛГУ, 1950,
    Новосибірськ: СБ АН РСРСР, 1962.
  \end{enumerate}
\item
  Соболєв, С. Л. Рівняння математичної фізики.--- М · Наука, 1966.
\item
  Траптер К. Дж. Інтегральні перетворення у математичній фізиці.---
  Москва: Гостехиддат, 1957.
\item
  \begin{enumerate}
  \def\labelenumi{\arabic{enumi}.}
  \setcounter{enumi}{60}
  \tightlist
  \item
    Трев Ф. Лекції з теорії диференціальних рівнянь у похідних
    похідних.--- Москва: Мир, 1965.
  \end{enumerate}
\item
  Трікомі Ф. Лекції з диференціальних рівнянь у похідних похідних.
  Москва, IL Publ., 1957. Фрідман А. Рівняння з частковими похідними
  параболічного типу. --- Москва: Мир, 1968.
\item
  Германдер Л. Лінійні диференціальні оператори з частковими
  похідними.--- Москва: Мир. 1965.
\item
  \begin{enumerate}
  \def\labelenumi{\arabic{enumi}.}
  \setcounter{enumi}{64}
  \tightlist
  \item
    Шварц Л. Математичні методи фізичних наук.--- Москва: Мир, 1965.
  \end{enumerate}
\item
  \begin{enumerate}
  \def\labelenumi{\arabic{enumi}.}
  \setcounter{enumi}{65}
  \tightlist
  \item
    Шилов Г. Е. Математичний аналіз. Спеціальний курс. Москва,
    Фізмат-гіз видавництво, 1960.
  \end{enumerate}
\item
  \begin{enumerate}
  \def\labelenumi{\arabic{enumi}.}
  \setcounter{enumi}{66}
  \tightlist
  \item
    Щилов, Г. Е. Математичний аналіз. Другий спеціальний курс.---
    Москва: Наука, 1965.
  \end{enumerate}
\item
  \begin{enumerate}
  \def\labelenumi{\arabic{enumi}.}
  \setcounter{enumi}{67}
  \tightlist
  \item
    Яненко Н. Н. Метод дробових щагів для розв'язання багатовимірних
    задач математичної фізики.--- Новосибірськ: Наука, 1967.
  \end{enumerate}
\end{itemize}

\section{ІНДЕКС
ІМЕНІ}\label{ux456ux43dux434ux435ux43aux441-ux456ux43cux435ux43dux456}

Бернуллі 329, Бессель 106, Буффон 322, Ніколсон 296, Ньютон 296

Гамільтон 335, 336 Гаусс 286 Гельмгольц 223 Грін 10 Гук 147 Гюйгенс 177,
181 Пуассон 180, 235

Д'Аламбер 126--129, Діріхле 85, 240, Дюамель 107, 108, 110--112, Рітц
338, 340, Робін 240, Родюг 268

** Seidel 286

Сімсон 320 Стокс 8, 180

Кіз 305, 306 Крхгоф 124

Кірхгоф 124 Кранк 296 Крістоффель 363 Ферма 335 Фур'є 24, 25, 36, 41,
74, 84, 91

Гакслі 304, Ханкел 105, 106, Хевісайд 82, 115, 116, Ходжкін 304-306

Лагранж 329 Ламм 307 Лаплас 8, 50, 98-107 Лежандр 235 Лельбніц 180
Лібман 285 Ліувіль 51, 58, 66, 148

Шварц 363 Шредінгер 8 Штурм 51, 58, 66, 148

Ейлер 84, 215, 329

\textbf{М}аксвелл 8 Меллін 105, 106

10sr (Молодий) 124

Navier 8 Neumann 240, 286

Jacobi 286

\section{ІНДЕКС
ПРЕДМЕТІВ}\label{ux456ux43dux434ux435ux43aux441-ux43fux440ux435ux434ux43cux435ux442ux456ux432}

Матрична алгебра 215, 216 Аналітичне розв'язання 301 Наближення сплайну
10

Позачасовий час 165

Калькулюс варіацій 329 Методи варіацій 10 Зовнішня похідна 79 Хвилі, що
рухаються, 149

-- квартира 177 -- стоячи 222

сферичний 178

--- ударні 210 циліндричні 177, 182

Гармонік 84, 154 Сферичний 265 градієнів 22 Граничні умови 17, 22, 23,
26

--- лінійний 26

--- однорідні 36, 46 --- гетерогенні 26, 49, 50

Дельта-функція 276 Діагоналізація матриці 216 Задача дифузії 113, 163
Дифузійний потік 114 --- член 113 Дифузія 113

Задача ** 147 -- Діріхле 242, 243, 246

\begin{itemize}
\item
  Діріхле - Нойман 364 - Кірхгоф 124
\item
  Коуу 69, 95, 116, 126, 207
\end{itemize}

--- мінімізація 10 --- Нойман 243---245

Ньютон 52, 121

--- на брахістохроні 329 Штурму --- Ліувіль 51, 58, 59, 66,

Закон збереження енергії 31, 114, 205

--- Теплопередача Ньютона 8 246

--- Фур'є 24, 25 Звукові хвилі 124 Момент згинання 157 Інтеграл Фур'є 87
Інтегральні перетворення 9, 73.76. Рівняння 10

Класифікація UCP 308 Крутильні осциляції 123 Поздовжні 123 Конвективні
хвилі 116 Конвективний транспорт 113 --- Член 113 Конвекція 113
Скінченні різниці 282 Канонічні координати 127, 132 --- Полярні 233,
235-238 Сферичні 233, 235, 236, 239, 266 Циліндричні 233, 235, 236, 239
--- еліптичний 239 Неперервний набір 96 Косинусне перетворення 73 ---
скінченне 185, 186 --- Обернене 75, 185, 186 Коефіцієнт дифузії 119 ---
теплопередача 23 Задача граничного значення 240, 241

Лапласіан 233 Лінійність 11

--- --- Еліптичний тренажер 241

Метод дедукції 319 -- Лібман 285, 286 -- Монте-Карло 316 -- Рітц 338,
340 -- Сімпсон 320 спуск 182 -- Фур'є 36 Режим коливання 154 Модуль Юнга
124, 145

Початковий зсув 231 Початкові умови 18

Оператор Лапласа 233 Ортогональність 41, 84, 90

\begin{itemize}
\tightlist
\item
  Ходжкіна - Гакслі 304
\end{itemize}

Власні вектори 67, 215, 217

Випадкові номери 319

Змішане завдання 19

--- enachenia 54, 215, 216 Відгук 96 функція 10, 65 Конформне
відображення 253, 356 Крістоффель-Шварц 363 --- частота 162 Спектральна
функція 76, 77 --- неперервна 77 --- частота 85 Stock 274 Параметрична
ідентифікація 304 Змінна залежна 8 Стоячі хвилі 148, 149 Суперпозиція
97, 193, 223 --- незалежна 8 Схема підрахунку тікання 289 Щільність
матеріалу 30 --- явна 289 Многочлен Лежандра 265 Кранк--Ніколсон 296,
299, 328 Порядок рівняння 10 Потенціал Pewtop 266 з круговою симетрією
240. сферично-симетрична 240 циліндрично симетрична 270 Температурна
характеристика 69 Теорема Діріхле 85 Правило Лейбніца 180--Стокс 180
Перетворення залежних змінних Теорія збурень 10 Тепловий імпульс 96
інтеграл 9, 73, 76 Теплоємність 30 --- координати 9 --- Лаплас 50, 75,
98---107, 115, 116, теплопровідність 25, 30, 33 126 Точкове джерело 274
--- обернене 75 --- Меллін 105, 106 --- Ханкель 76, 105, 106 ---
обернене 76. Питома теплоємність 34 --- Фур'є 74, 75, 83, 85, 89, 91
Вирівнювання Бесселя 106, 224, 239 --- хвиля 8, 142, 169, 234 --- ---
скінченне 185 --- нелінійне 304 --- пряме 88, 116, 126 --- обернене 88,
74, 75 гармонічні коливання 223 --- Гельмгольц 223, 224, 235 Принцип
Гамільтона 335, 336 --- Гюйгенс 177, 181 --- Дюамель 107, 108, 110, 112
Рух Ньютона 8 дифузія 18, 197 - Суперлосоцитності 191, 196, 250
конвектива 19, 63, 113, 119, 120 ** - Ферма 335** - Ламм 307 - Лаплас 8,
170, 235, 258, 285, 363 - Лежандр 239, 268 Змінне розділення 9, 191
лінійне 11 - гіперболічне 11, 12, 121, 169 Розклад власних функцій 10,
66 - параболічне 11, 12, 14, 169 - еліптичне 11, 12, 169 - Максвелл 8 -
Нав'є - Стокс 8 Згортка 90, 93 неліцейних {[}{]} гетерогенних 11, 19,
65. --- скінченна 102 Гравітаційна сила 122 гомогенна P Пуассон 235. --
--- ранкового перетворення 122 синусоїдального перетворення 75 постпост
285 скінченне 185, 186, 193 зворотне 75, 185, 186 реакції з нелінійною
дифузією Система телеграфних рівнянь 124 --- з прямокутними
коефіцієнтами 60.

170

з частковими похідними 7

Телеграф 8, 123, 125, 215

Зберігання тепла 16

{\def\LTcaptype{none} % do not increment counter
\begin{longtable}[]{@{}llllllll@{}}
\toprule\noalign{}
\_ & tep & Lonre & Доступність води & 8, & 16, & 30, & 169. \\
\midrule\noalign{}
\endhead
\bottomrule\noalign{}
\endlastfoot
& & 234. & & • & • & • & , \\
\end{longtable}
}

\begin{itemize}
\tightlist
\item
  характеристики 202, 208
\end{itemize}

--- Шредінгер, 8 років

--- Ейлер 215, 256, 268

--- Ейлер-Лагранж 329, 330, \textbf{33}2

Стан Діріхле 240

\begin{verbatim}
- Neumann 240 - Robin 240
\end{verbatim}

Сингулярний диференціал 268

Гріна 10, 96, 196, 273, 276

д'Аламбер 126-132

\begin{itemize}
\tightlist
\item
  додатково до ймовірнісного інтеграла 80
\end{itemize}

Неточника 96, 276

-- неперіодичні 77 -- періодичні 77 -- Хевісайд 82, 115, 116, 365.
Фур'є-образ 88, 89, 91, 93

Формула д'Аламбера 180, 183

--- Лейбніц 184

Джерела 99, 100

\begin{itemize}
\tightlist
\item
  Пуассон 180
\end{itemize}

--- --- інтеграл 248, 249, 251

--- Родріга 268

--- Ейлер 84, 85 Фундаментальне рішення 37 Функціональність 329

Функція Бесселя 224, 225

Чисельні методи 9 Кількість персемних 11

Рівнопотенціальні лінії 277

Ядро перетворення 73 -- перетворення Пуассона 252

\section{ЗМІСТ}\label{ux437ux43cux456ux441ux442}

{\def\LTcaptype{none} % do not increment counter
\begin{longtable}[]{@{}
  >{\raggedright\arraybackslash}p{(\linewidth - 4\tabcolsep) * \real{0.1856}}
  >{\raggedright\arraybackslash}p{(\linewidth - 4\tabcolsep) * \real{0.1237}}
  >{\raggedright\arraybackslash}p{(\linewidth - 4\tabcolsep) * \real{0.6907}}@{}}
\toprule\noalign{}
\begin{minipage}[b]{\linewidth}\raggedright
З редагування
\end{minipage} & \begin{minipage}[b]{\linewidth}\raggedright
Гора
\end{minipage} & \begin{minipage}[b]{\linewidth}\raggedright
Переклад
\end{minipage} \\
\midrule\noalign{}
\endhead
\bottomrule\noalign{}
\endlastfoot
Передмова & viye & ****** \\
& & \\
Частина 1. & Введено & \\
Лекція & 1. & Вступ до теорії диференціальних рівнянь з похідними
похідними \\
Частина 2. & Діф & Завдання з синтезу \\
Лекція & 2. & Задачі типу дифузії (параболічні рівняння) \\
Лекція & 3. & Граничні умови в задачах дифузії \\
Лекція & 4. & Виведення рівняння теплопровідності , \\
Лекція & 5. & Розщеплення змінних \\
Лекція & 6. & Перетворення неоднорідних граничних умов у одно- \\
& & Родичі \\
Лекція & 7. & Розв'язання складніших задач методом розділення змін- \\
Тторене \textasciitilde{} & o & Перетворити комплексні рівняння у просту
форму \\
Лекція-лекція & o. & Перетворення комплексних рівнянь у простий
вигляд \\
Мексика & σ, & Власні функції , \\
Pection & 10 & Інтегральні перетворення (перетворення синусів і
козіпус- \\
Боулз & 10. & ) \\
Лекція & 11. & Перетворення рядів і четвірок \\
Лекіня & 12 & Перетворення Фур'є та його застосування до розв'язку \\
* 10HM 11 /1 & , & Рівняння з похідними похідними \\
Лекція & 13. & Перетворення Лапласа \\
Лекція & 14. & Принцип Дюамеля \\
Лекчич & 15. & Принцип Дюамеля \\
& & \\
Частина 3. & Guipe & Проблеми \\
Лекпея & 16. & Одномірне хвильове рівняння (гіперболічні еквалізації) \\
& & \\
Лекція & 17. & Формула Д'Аламбера \\
Лекція & 18. & Формула Д'Аламбера (продовження) \\
Лекція & 19. & Рівняння води та граничні умови \\
Лекція & 20. & Вібрації обмеженої струни (стоячі хвилі) 14 \\
Лекція & 21. & Осциляції пучка (рівняння часткового виробника \\
& & Четвертий порядок) \\
& & Перехід на безрозмірні змінні \\
Лекція & 23. & Класифікація диференціальних рівнянь у похідних \\
& & (Канонічна форма гіперболічного рівняння) 10 \\
Лекція & 24. & Хвильове рівняння у вільному просторі (двовимірне \\
\_ & \textasciitilde{} - & та тривимірні задачі) \\
Лекція & 25. & Скінченні перетворення Фур'є (синус і косинус до - \\
& 00 & Освіта) \\
& & Принцип суперпозиції --- основи теорії динамічних систем 19 \\
& & Рівняння першого порядку (характеристичний метод) 19 \\
Лекція & 28. & Нединамічні рівняння першого порядку (закони
збереження \\
t. & 00 & Ньяя) , \\
Лекція & 29. & Системи диференціальних рівнянь у похідних 2 \\
Розділ & & Мембранні осциляції (покращення хвиль у Polar 2000
Godinates) \\
& & БУУУУММаалака аа аа аа аа аа \\
\end{longtable}
}

{\def\LTcaptype{none} % do not increment counter
\begin{longtable}[]{@{}
  >{\raggedright\arraybackslash}p{(\linewidth - 2\tabcolsep) * \real{0.7658}}
  >{\raggedright\arraybackslash}p{(\linewidth - 2\tabcolsep) * \real{0.2342}}@{}}
\toprule\noalign{}
\begin{minipage}[b]{\linewidth}\raggedright
Частина 4. Еліптичні задачі
\end{minipage} & \begin{minipage}[b]{\linewidth}\raggedright
233
\end{minipage} \\
\midrule\noalign{}
\endhead
\bottomrule\noalign{}
\endlastfoot
Лекція 31. Лапласнан (інтутивна олісанація) & 233240248255 \\
гармоніки) & 264273 \\
Частина 5. Чисельні та наближені методи & 282 \\
Лекція 37. Чисельні розв'язки (еліптичні задачі) & 282289 \\
Lsson) & 296301 \\
Рівняння) & 308316 \\
Монте-Карло & 322329 \\
Лекція 45. Варіаційні методи розв'язання диференціальних рівнянь у
похідних похідних & 338 \\
Лекція 46. Розв'язок рівнянь часткових похідних теоретиків збурень &
346 \\
Лекція 47. Розв'язання диференціальних рівнянь у похідних похідних за
методом конформних відображень & 355 \\
Додаток 1 & 365 \\
Додаток 2 & 371 \\
Spinning 3 & 373 \\
Джерела & 375 \\
Індекс назв & 378 \\
Індекс предметів & 379 \\
\end{longtable}
}

\section{Дорогий
читачу!}\label{ux434ux43eux440ux43eux433ux438ux439-ux447ux438ux442ux430ux447ux443}

Будь ласка, надсилайте свої коментарі щодо змісту приладів, його
конструкції, якості перекладу та іншого за адресою: 129820, Москва,
I-110, GSP. 1-й Рижський провулок, 2, видавництво «Мир».

\section{Фарлоу
Стейпле}\label{ux444ux430ux440ux43bux43eux443-ux441ux442ux435ux439ux43fux43bux435}

\section{ДИФЕРЕНЦІАЛЬНІ РІВНЯННЯ В ПОХІДНИХ ПОХІДНИХ ДЛЯ НАУКОВЦІВ І
ІНЖЕНЕРІВ}\label{ux434ux438ux444ux435ux440ux435ux43dux446ux456ux430ux43bux44cux43dux456-ux440ux456ux432ux43dux44fux43dux43dux44f-ux432-ux43fux43eux445ux456ux434ux43dux438ux445-ux43fux43eux445ux456ux434ux43dux438ux445-ux434ux43bux44f-ux43dux430ux443ux43aux43eux432ux446ux456ux432-ux456-ux456ux43dux436ux435ux43dux435ux440ux456ux432}

Науковий редактор М. В. Хатунцева-молодший науковий редактор Е. І.
Качулина Художник В. Е. Карпов Художній редактор В. І. Шаповалов
Технічний редактор Т. А. Максимова Коректор Т. І. Стіфсева

\section{UB No 5289}\label{ub-no-5289}

\texttt{Набор\ 27.12.84.\ Подано\ до\ публікації\ 01.07.85.\ Формат\ \textquotesingle{}60\textbackslash{}times90\^{}\{1\}\textquotesingle{}\ є.\ Папір\ книги.\ журналу.\ Печатка\ висока.\ Літературний\ Гарінтура.\ Том\ L2\ —\ бум.\ n,\ Usl.\ Печ.\ с.\ 24.\ Стан.\ Cr.-Ott.\ 24.\ Уч.-изд.\ л.\ 19.69,\ Ізд.\ No\ 1/3807.\ Тираж\ 5\ 000\ примірників.\ Замовлення\ L6\ 241.\ Ціна\ 1\ ранд.\ 70\ тис.}

MIR 129820, GSP, Москва, I-110, 1-й Рижський провулок, 2

Надруковано в Ленінградській друкарні No 6 Ордена Червоного Прапора
Праці Ленінградської асоціації «Технічний король» імені Євга. Соколова
Союзполіграф при Державному комітеті СРСР з видавничої справи,
друкарства та книжкової торгівлі, 193144, Ленінград, вулиця Моісеєнко,
10, з матрицями Ордена Жовтневого Регулюта та Ордена Червоного Прапора
Праці МПО «Перша зразкова друкарня» імені А. А. Жданога
Союзполіграфірому при Державному комітеті СССГ з видавничої, друкарської
та книжкової торгівлі, 113054, Москам, Валовія, 28

\end{document}
